% concatenated from main.tex
\documentclass[11pt, a4paper]{article}
\usepackage{amssymb,amsmath,amsthm}
\usepackage{fancyhdr}
\usepackage[british]{babel}
\usepackage{color}
\usepackage[title]{appendix}
\usepackage{amsmath,amsthm}


\usepackage{enumerate,enumitem}
\usepackage{enumitem,verbatim}

\usepackage{ulem}
\usepackage{hyperref}
\usepackage[margin=2.5cm]{geometry}
\usepackage{lineno}
\usepackage[ruled]{algorithm2e}
\usepackage{todonotes}

\newtheorem{theorem}{Theorem}[section]
\newtheorem{prop}[theorem]{Proposition}
\newtheorem{lemma}[theorem]{Lemma}
\newtheorem{cor}[theorem]{Corollary}
\newtheorem{assump}[theorem]{Setup}
\newtheorem{conjecture}[theorem]{Conjecture}
\newtheorem{definition}[theorem]{Definition}

\theoremstyle{definition}
\newtheorem{problem}{Problem}
\newtheorem{question}{Question}
\renewcommand{\thequestion}{\Alph{question}}
\newtheorem{expl}{Example}
\renewcommand{\theexpl}{\Alph{expl}}
\newtheorem{defn}[theorem]{Definition}
\newtheorem{rmk}{Remark}
\renewcommand{\thermk}{\Alph{rmk}}
\newtheorem{claim}[theorem]{Claim}
\newtheorem{fact}[theorem]{Fact}
\newtheorem{step}{Step}

\newenvironment{poc}{\begin{proof}[Proof of claim]\renewcommand*{\qedsymbol}{$\blacksquare$}}{\end{proof}}
\newcommand{\btheorem}{\begin{theorem}}
\newcommand{\etheorem}{\end{theorem}}
\newcommand{\bconjecture}{\begin{conjecture}}
\newcommand{\econjecture}{\end{conjecture}}
\newcommand{\bproposition}{\begin{proposition}}
\newcommand{\eproposition}{\end{proposition}}
\newcommand{\bdefinition}{\begin{definition}}
\newcommand{\edefinition}{\end{definition}}
\newcommand{\bcorollary}{\begin{corollary}}
\newcommand{\ecorollary}{\end{corollary}}
\newcommand{\bproof}{\begin{proof}}
\newcommand{\eproof}{\end{proof}}
\newcommand{\bclaim}{\begin{claim}}
\newcommand{\eclaim}{\end{claim}}
\newcommand{\bquestion}{\begin{question}}
\newcommand{\equestion}{\end{question}}
\newcommand{\bfact}{\begin{fact}}
\newcommand{\efact}{\end{fact}}
\newcommand{\bremark}{\begin{remark}}
\newcommand{\eremark}{\end{remark}}
\newcommand{\eexample}{\end{example}}
\newcommand{\bexample}{\begin{example}}
\newcommand{\la}{\langle}
\newcommand{\elemma}{\end{lemma}}
\newcommand{\blemma}{\begin{lemma}}
\newcommand*{\abs}[1]{\lvert #1\rvert}
\newcommand*{\floor}[1]{\lfloor #1\rfloor}
\newcommand{\xc}{\overline{X}}
\newcommand{\be}{\beta}

\newcommand{\cH}{\mathcal{H}}
\newcommand{\F}{\mathcal{F}}
\newcommand{\C}{\mathcal{C}}
\newcommand{\N}{\mathbb{N}}
\newcommand{\sfH}{\mathsf{H}}
\newcommand{\sfG}{\mathsf{G}}
\newcommand{\sfT}{\mathsf{T}}
\newcommand{\ex}{\mathrm{ex}}
\newcommand{\eg}{e.g.\ }
\newcommand{\ie}{i.e.\ }
\def \hcf{\mathrm{hcf}}
\newcommand{\eps}{\varepsilon}
%\newcommand{\exp}{\text{exp}}
\newcommand{\<}{\subseteq}
\newcommand{\blue}[1]{\textcolor[rgb]{0.00,0.00,1.00}{#1}}
\newcommand{\red}[1]{\textcolor[rgb]{1.00,0.00,0.00}{#1}}
\newcommand{\De}{\Delta}
\newcommand{\de}{\delta}

\title{Clique factors in random samplings of regular graphs}

\author{
Wanting Sun\thanks{Data Science Institute, Shandong University, Jinan, China. Email: {\tt wtsun@sdu.edu.cn}.},\ \ 
Shunan Wei\thanks{School of Mathematics, Shandong University, Jinan, China. Email: {\tt snwei@mail.sdu.edu.cn}.},\ \ 
Donglei Yang\thanks{School of Mathematics, Shandong University, Jinan, China. Email: {\tt dlyang@sdu.edu.cn}.}}

\date{\today}

\begin{document}
\maketitle
\linespread{1.2}
%\linenumbers

\begin{abstract}
We show that for any integer $r\ge 2$, there exists a constant $c>0$ such that for every sufficiently large integer $n$, every $((r-1)n+1)$-regular graph $G$ on $rn$ vertices has at least $c2^{rn}$ subsets $S\subseteq V(G)$ such that $G[S]$ contains a $K_r$-factor. This confirms a conjecture of Dragani\'c, Keevash and M\"uyesser for large $n$ [\textit{Cyclic subsets in regular Dirac graphs. Int. Math. Res. Not., 2025(14): 1-16, 2025}].
\end{abstract}

%%%%%%%%%%%%%%%%%%%%%%%%%%%%%%%%%%%%%%%%%%%%%%%%%%%%%%%%%%%

\section{Introduction}
A classical line of research in extremal graph theory concerns about the  minimum degree condition that guarantees the existence of certain spanning subgraphs.   Once such a result is established,  a natural further question is to measure the \textit{robustness} of the graph with respect to the property of containing the given spanning structure. This direction of research aims to strengthen classical results in extremal and probabilistic combinatorics.


%{Leave it for Wei to rephrase}~ 
Hamiltonicity is one of the
most central problems in graph theory. A classical result proved by Dirac in 1952~\cite{Dirac} states that for $n\geq 3$, every $n$-vertex graph $G$ with $\delta(G) \ge \frac{n}{2}$ (known as \textit{Dirac graph}) contains a Hamilton cycle.  
 %Investigating when the classical theorems hold in a \textit{robust} or \textit{resilient} way has emerged as a highly active research topic. 
% with Sudakov and Vu~\cite{Sudakov2019} pioneering its systematic study by first articulating its significance.
%Determining minimum degree conditions that guarantee the existence of a spanning structure is a central problem in extremal graph theory.  
%A cornerstone result in this direction is Dirac's theorem~\cite{Dirac}, which states that for $n\geq 3$, every $n$-vertex graph $G$ with $\delta(G) \ge \frac{n}{2}$ (known as \textit{Dirac graph}) contains a Hamilton cycle.  
%Angluin and Valiant~\cite{Angluin1979} provided  a robust version of Dirac's theorem, that is, the random graph $G(n,p)$ contains a Hamilton cycle with $p\ge \Omega(\frac{\log n}{n})$. 
A natural way to strengthen Dirac’s theorem is to ask how many Hamilton cycles a Dirac graph must contain. S\'{a}rk\"{o}zy, Selkow and Szemer\'{e}di~\cite{Sarkozy2003} proved that every Dirac graph contains at least $c^nn!$ Hamilton cycles for some small positive constant $c$, and conjectured that $c$ can be improved to $\frac{1}{2}-o(1)$. This conjecture was resolved by Cuckler and Kahn~\cite{Cuckler2009}. 
{We refer the reader to a survey of Sudakov~\cite{Sudakov2017} where various measures of robustness and relevant results are collected.}
%Cuckler and Kahn \cite{Cuckler2009} proved this  by showing a stronger result that every Dirac graph with minimum degree $\delta$ contains at least $(\frac{\delta}{e+o(1)})^n$ Hamilton cycles. 
%Determining sufficient minimum degree conditions for the existence of an $H$-factor is the fundamental research in extremal graph theory.  In this direction, the cornerstone theorem of Dirac~\cite{Dirac} states that every $n$-vertex graph $G$ with minimum degree $\delta(G)\ge \frac{n}{2}$ contains a Hamilton cycle. In particular, if $n$ is even then $G$ has a perfect matching (i.e. $K_2$-factor). A graph on $n$ vertices is called a \textit{Dirac graph} if its minimum degree is at least $n/2$.
%There are many further results to illustrate the Hamiltonicity problem, such as counting the number of Hamiltonian cycles in Dirac graphs~\cite{Sarkozy2003,Cuckler2009}, the robustness or resilience of Hamiltonicity and so on. 
% (such as the random graph $G(n,p)$, which is Hamiltonian with $p\ge \Omega(\frac{\log n}{n})$; see~\cite{Angluin1979, Shamir1983}), and so on. 

% where various measures of robustness and relevant results are collected, some further results on resilience can be found in~\cite{Krivelevich2011,Montgomery2019, Sudakov2019}.

We introduce the robustness problem from the following two perspectives. 
% A widely studied robustness problem concerns the Hamiltonicity of Dirac graphs or random graphs in which the edges are sampled at random and substantial progress has been made (see~\cite{Frieze2008,Lee2012}).
Let $G$ be a graph with a property $\mathcal{P}$.
On the one hand, we sample each edge uniformly at random in $G$ and study the threshold for the property $\mathcal{P}$ of the resulting random subgraph. 
% Recently, there has been increasing interest in the study of \textit{robustness} and \textit{resilience} of graph properties, aiming to strengthen classical results in extremal and probabilistic combinatorics.
% For a comprehensive overview about robustness, we refer the reader to the survey~\cite{Sudakov2017}.
% For example, given an $n$-vertex graph with minimum degree at least $\frac{n}{2}$, Dirac theorem guarantees that there are at least one Hamilton cycle, but can one determine the number of such cycles?
% Cuckler and Kahn~\cite{Cuckler2009} answer this question and they prove that Dirac graph with $n$ vertices contains at least $(1/2-o(1))^nn!$ Hamilton cycles.
{More precisely, let $G(p)$ be a random subgraph of $G$ such that every edge in $G$ is kept independently and randomly with probability $p$.} 
Krivelevich, Lee and Sudakov~\cite{Krivelevich2014} studied the threshold for hamiltonicity and showed that there exists a constant $C$ such that for any $n$-vertex Dirac graph $G$ and $p\ge \frac{C\log n}{n}$,  %if one takes edges of $G$ independently at random with probability $p$, 
\textit{with high probability}\footnote{We say that an event $E$ holds \textit{with high probability,} (w.h.p. for short) if $\mathbb{P}[E]\rightarrow 1$ as $n$ tends to infinity.} $G(p)$ is Hamiltonian.
On the other hand, one may sample vertices rather than edges at random and investigate the probability that the induced subgraph on such a random vertex set still possesses property \(\mathcal{P}\). {In this direction, Erd\H{o}s and Faudree~\cite{Erdos1999} asked whether every regular Dirac graph contains a constant proportion of Hamiltonian vertex subsets.} 
%\sout{This was motivated by a conjecture of Erd\H{o}s and Faudree~\cite{Erdos1999} as follows.}

\begin{conjecture}[\cite{Erdos1999}]\label{conj:Hamilton cycle}
Any $(n+1)$-regular graph $G$ on $2n$ vertices has at least $c2^{2n}$ subsets of $V(G)$ inducing a Hamiltonian subgraph for some absolute $c>0$.
\end{conjecture}

{This conjecture is tight in several ways: the \((n+1)\)-regularity cannot be relaxed to \(n\)-regularity (witnessed by \(K_{n,n}\)) nor weakened to a minimum degree condition (illustrated by adding a spanning star to each class of \(K_{n,n}\)). %On the one hand, the \((n+1)\)-regularity condition cannot be replaced by \(n\)-regularity, as shown by the complete bipartite graph $K_{n,n}$. On the other hand, the regularity assumption cannot be weakened to minimum degree, which can be seen from the graph $G$ obtained by adding a spanning star in each part of $K_{n,n}$.} %We remark that the $(n+1)$-regularity assumption cannot be weakened to the minimum degree $n+1$ or $n$-regularity.  In particular, the construction of complete bipartite graph $K_{n-1,n+1}$ by adding a $2$-factor in the part of size $n+1$ uniquely minimizes the number of Hamiltonian subsets. 
Conjecture~\ref{conj:Hamilton cycle} was recently solved for large $n$ by Dragani\'c, Keevash and M\"uyesser~\cite{Keevash2025}, who showed that a uniformly random vertex subset of an $(n+1)$-regular $2n$-vertex graph induces a Hamiltonian graph with probability at least $\frac{1}{2}$. The bound  $\frac{1}{2}$ is tight, as shown by the example of the  complete bipartite graph $K_{n-1,n+1}$ with a $2$-factor added to the larger side.    %Conjecture~\ref{conj:Hamilton cycle} for large $n$ by showing that and they determined the optimal value $c$ to be precisely $\frac{1}{2}$. We can phrase Conjecture~\ref{conj:Hamilton cycle} as taking a random vertex-subset by including each vertex independently with probability $\frac{1}{2}$ and to show that the subgraph induced on this set is Hamiltonian with constant probability.}
%Wei: Yes! This is their main result.
Very recently, Hunter, Liu, Milojevi\'c and Sudakov~\cite{hunter2025}  investigated a tournament analogue of Conjecture~\ref{conj:Hamilton cycle} under a minimum semi-degree condition.} % degree and give an optimal number of subsets of $V(T)$ inducing a Hamiltonian subgraph.}
%In the setting of digraph, Ghouila and Houri~\cite{} show an analogue of Dirac theorem, that is, every $n$-vertex digraph with minimum semi-degree at least $\frac{n}{2}$ contains a Hamilton cycle.



\subsection{The Hajnal-Szemer\'edi theorem}
%\sout{Given a graph $H$ that has a connected component of size at least $3$, Hell and Kirkpatrick~\cite{Hell1983} prove that the decision problem for $H$-factors is NP-complete.}
Given graphs $G$ and $H$, an \textit{$H$-tiling} in $G$ is a collection of vertex-disjoint copies of $H$. 
A \textit{perfect $H$-tiling} (or \textit{$H$-factor}) is one that covers every vertex of $G$.  A celebrated theorem of Hajnal and Szemer\'edi~\cite{HajnalSz} provides the best possible minimum degree condition that guarantees the existence of a $K_r$-factor (the case $r=3$ was previously obtained by Corr\'adi and Hajnal~\cite{Corradi}), {which answers a conjecture of Erd\H{o}s \cite{E1964}. }%{conjecture of Erdos}
\begin{theorem}[\cite{Corradi, HajnalSz}]
\label{thm:HSz}
Let $G$ be an $n$-vertex graph with $n\in r\mathbb{N}$. If $\delta(G)\ge \left(1-\frac{1}{r}\right)n$, then $G$ contains a $K_r$-factor.
\end{theorem}


The minimum degree condition is best possible by considering slightly unbalanced complete $r$-partite graphs. Later, Kierstead and Kostochka \cite{Kierstead2008} gave a short proof of the Hajnal-Szemer\'edi theorem phrased in terms of equitable colorings. {For more recent results on equitable colorings, see \cite{cheng2025,kierstead-survey}.} %\sout{and they also prove the analogue of Hajnal-Szemer\'edi theorem under an Ore-type degree condition~\cite{Kostochka2008}}.
%\sout{Here and through this paper, we always assume the necessary condition $n\in r\mathbb{N}$, otherwise we will draw special attention to it.}
% For a general graph $H$, Alon and Yuster~\cite{} gives an asymptotic result, that is, if $\delta(G)\ge (1-\frac{1}{\chi(H)})n+o(n)$, then $G$ contains an $H$-factor, where $\chi(H)$ is the chromatic number of $H$.
% After some excellent work (see e.g.~\cite{}), it is finally solved by K\"{u}hn and Osthus up to an additive constant.
Pham, Sah, Sawhney and Simkin~\cite{Pham2} proved a random sparsification version of Theorem~\ref{thm:HSz}.  
More precisely, they showed that given an $n$-vertex graph $G$ with $\delta(G)\ge \left(1-\frac{1}{r}\right)n$, 
if $p\ge Cn^{-2/r}(\log n)^{1/{r \choose 2}}$, then w.h.p.~$G(p)$ contains a $K_r$-factor.  
The case $r=3$ was established earlier by Allen, Böttcher, Corsten, Davies, Jenssen, Morris, Roberts and Skokan~\cite{Allen}. 
{Kelly, M\"uyesser and Pokrovskiy~\cite{Kelly2024}, and independently Joos, Lang and Sanhueza-Matamala~\cite{Joos-Lang2023},  extended this result to general $F$-factors in hypergraphs, determining the  {asymptotically} optimal threshold $p$ for which  $G(p)$ contains an $F$-factor with high probability. For further results on transversal robust versions of some classical theorems, we refer the reader to~\cite{Anastos-Chakraborti, Ferber2022, Han2025}. 
%There are many results about the transversal robust version of Dirac' theorem~\cite{Ferber2022,Anastos-Chakraborti} and Theorem~\ref{thm:HSz}~\cite{Ferber2022,Han2025}. 
}%{(add more references on $G[p]$)}

By analogy with Conjecture~\ref{conj:Hamilton cycle}, one may ask how many vertex subsets $S\subseteq V(G)$ have the property that $G[S]$ contains a $K_r$-factor.   Dragani\'{c}, Keevash and M\"{u}yesser~\cite{Keevash2025} proposed the following {conjecture regarding the Hajnal-Szemer\'{e}di theorem.} %{robust version of the Hajnal–Szemer\'{e}di theorem.}
\begin{conjecture}[\cite{Keevash2025}]\label{conj:K_r factor}
For any $r\ge 2$ there is some constant $c>0$ so that if $G$ is an $((r-1)n+1)$-regular graph on $rn$ vertices, then at least $c2^{rn}$ subsets of $V(G)$ induce a $K_r$-factor.
\end{conjecture}

%\sout{They also proposed that a plausible class of extremal constructions for Conjecture~\ref{conj:K_r factor} may be  slightly unbalanced complete $r$-partite graphs perturbed with a suitable factor in the largest part.}
In the following, we make some remarks on this conjecture. %also remark that:
\begin{enumerate}[label=(\arabic*)]
    \item The regularity condition $((r-1)n+1)$ cannot be lowered to $(r-1)n$.  
Indeed, consider the balanced complete $r$-partite graph $G:=K_{n,\ldots,n}$: an induced subgraph $G[S]$ contains a $K_r$-factor exactly when $S$ contains the same number of vertices from each part. Hence the number of such subsets $S$ is %$\sum_{k=1}^n \binom{n}{k}^r\le n \big(\frac{\sum_{k=1}^n\binom{n}{k}}{n}\big)^r=\frac{2^{rn}}{n^{r-1}}$;

$$
\sum_{k=1}^n \binom{n}{k}^r \le n \cdot \binom{n}{\lfloor n/2 \rfloor}^r \leq n \left( \frac{2^n}{\sqrt{\pi n/2}} \right)^r =n\frac{2^{rn}}{(\pi n/2)^{r/2}}.
 $$
 \item The regularity assumption cannot be reduced to a minimum degree of $(r-1)n+1$. 
Consider the graph $G$ obtained from  the balanced complete $r$-partite graph $K_{n,\dots,n}$ (with vertex partition $A_1\cup \ldots \cup A_r$) by adding a spanning star  in each part.  
In $G$, a subset $S\subseteq V(G)$ induces a subgraph that contains a $K_r$-factor only when the size of %the intersection of 
each $|A_i\cap S|$ %with each part 
deviates from $\frac{|S|}{r}$ by less than  $r$\footnote{Note that each $G[A_i\cap S]$ has matching number at most one. Thus $|A_i\cap S|-\frac{|S|}{r}\leq 1$ for each $i\in [r]$, which means $\min\limits_{i\in[r]}|A_i\cap S|>\frac{|S|}{r}-r$.}. {Then the number of such subsets $S$ is at most
$$
\sum_{s_1=0}^n\sum_{\substack{|s_j-s_1|<2r\\j\in[2,r] }}\prod_{i=1}^r\binom{n}{s_i}\leq \sum_{s_1=0}^n\sum_{\substack{|s_j-s_1|<2r\\j\in[2,r] }}\binom{n}{\lfloor n/2 \rfloor}^r\leq (n+1)(4r)^{r-1}\Big( \frac{2^n}{\sqrt{\pi n/2}} \Big)^r, 
$$
where $s_i=|A_i\cap S|$ for each  $i\in [r]$. }

%{compute the number of S?}%The regularity assumption cannot be weakened to the minimum degree $(r-1)n+1$ as a balanced complete multipartite graph $G=K_{n,\cdots,n}$ with $r$ spanning stars added in each part. In this case, $G[S]$ has a $K_r$-factor only if the size of each part in $G[S]$ deviates from $\frac{|S|}{r}$ by at most $r$.
\end{enumerate}

Our main contribution is to resolve Conjecture~\ref{conj:K_r factor} for all sufficiently large $n$.

\begin{theorem}\label{thm:main thm}
For any $r\ge 2$, there is a constant $c>0$ such that the following holds for sufficiently large $n$. Let $G$ be an $((r-1)n+1)$-regular graph on $rn$ vertices. Then at least $c2^{rn}$ subsets of $V(G)$ induce a $K_r$-factor.
\end{theorem}


In fact, we prove that the constant {$c$ can be taken as  $\frac{1}{(40r^2)^r}$}. It is worth noting that the case $r=2$ (perfect matchings) of Conjecture~\ref{conj:K_r factor} follows from the result of Dragani\'{c}, Keevash and M\"{u}yesser~\cite{Keevash2025} concerning Hamilton cycles, together with the fact that a required subset must have even size, which occurs with probability $\frac{1}{2}$.
Let $G[p]$ be a random induced subgraph of $G$ with each vertex kept independently with probability $p$, where $p\in (0,1)$. 
Consequently, Theorem~\ref{thm:main thm} can be reduced to studying the probability that \(G[\frac{1}{2}]\) contains a \(K_r\)-factor. %and the key ingredient in our proof of Theorem~\ref{thm:main thm} is the following.

\begin{theorem}\label{thm:the random version of main thm}
For any $r\ge 3$, let $G$ be an $((r-1)n+1)$-regular graph on $rn$ vertices with $n$ sufficiently large. 
Then $\mathbb{P}\big[G[\frac{1}{2}]~\text{admits a}~K_r\text{-factor}\big]\ge \frac{1}{(40r^2)^r}$.
\end{theorem}
{Indeed, the result of Theorem~\ref{thm:the random version of main thm} can be extended to a more general probability $p\in (0,1)$. 
%Let $G[p]$ be a random induced subgraph with each vertex kept independently with probability $p$.
Following the same  proof of Theorem~\ref{thm:the random version of main thm}, one may show that $G[p]$ contains a $K_r$-factor with probability at least $(\frac{p^2}{20r^2})^r$.
%Since the argument is nearly identical to the proof of Theorem~\ref{thm:the random version of main thm}, we skip the details.
}


\subsection{Notation}
For a graph $G$, let $V(G)$ and $E(G)$ denote the vertex set and edge set of $G$, respectively, and denote $|G|=|V(G)|$ and $e(G)=|E(G)|$. For a vertex $v\in V(G)$, let $N_G(v)$ be the neighborhood of $v$ in $G$. For two vertex subsets $U,U'\subseteq V(G)$, let $N_G(v,U)=N_G(v)\cap U$ and $N_G(U',U)=\bigcap_{v\in U'}N_G(v,U)$. We let \(d_{G}(v)=|N_{G}(v)|\) denote the degree of \(v\) and write \(d_{G}(v,U)=|N_{G}(v,U)|\). % for the degree of \(v\) into a subset \(U\subset V(G)\). 
As we do elsewhere, we omit \(G\) from the subscript whenever there is no risk of confusion. The minimum and maximum degree of \(G\) is denoted by \(\delta(G)\) and \(\Delta(G)\) respectively. 


For a subset \(U\subseteq V(G)\), let \(G[U]\) denote the graph induced by \(U\). Let \(G-U\) be the graph induced by \(V(G)\setminus U\). % {and \(G+U\) be the graph induced by \(V(G)\cup U\)}. 
Given a subset {\(U'\subset V(G)\setminus U\)}, let \(G[U,U']\) denote the bipartite graph with bipartition \(U\cup U'\) and edges of the form \(uu'\in E(G)\) with \(u\in U\) and \(u'\in U'\). Denote $e(U,U')=|E(G[U,U'])|$. %Given a subset \(U\subset V(G)\), let \(G-U\) be the graph induced by \(V(G)\setminus U\).%, and use similar natural and common other notation. 

%\sout{Given a set $U$ and an integer $k$, we write $\binom{U}{k}$ the collection of all $k$-element subsets of $U$.}
For any integers $a\leq b$, define $[a,b]:=\{i\in \mathbb{N}:a\leq i\leq b\}$ and $[b]:=[1,b]$. Given real numbers $a,b,c$, we write $a=b\pm c$ if it holds that $b-c\leq a\leq b+c$. When we write  $\alpha\ll \beta\ll \gamma$, we always mean that $\alpha, \beta, \gamma$  are constants in $(0,1)$; and $\alpha\ll \beta$ means that there exists $\alpha_0=\alpha_0(\beta)$ such that the subsequent arguments hold for all $0<\alpha\leq \alpha_0$. Hierarchies of other lengths are defined analogously. 

%\sout{Let $G$ be a graph on $rn$ vertices and  $\mathcal{P}:=\{V_1,\ldots,V_t\}$ be a vertex partition of $G$. We say $V_i$ is a \textit{large} part of $\mathcal{P}$ is $|V_i|>\frac{|G|}{r}$, and a \textit{small} part otherwise.}%  matching, vertex cover, 


\section{Overview}
Let $G$ be an $((r-1)n+1)$-regular graph on $rn$ vertices. Given $\gamma>0$ and a subset $S\subseteq V(G)$, we say $S$ is a $\gamma$-\textit{independent set} in $G$ if $e(G[S])\leq \gamma n^2$.  We divide the proof of Theorem~\ref{thm:the random version of main thm} into two parts according to whether $G$ contains a $\gamma$-independent set of size $n$ (extremal case) or not (non-extremal case). 

\begin{theorem}[Non-extremal case]\label{thm:Non-extremal case}
Suppose that $r\geq 3, n\in \mathbb{N}$ and $\frac{1}{n}\ll \gamma \ll \frac{1}{r}$. Let $G$ be an $((r-1)n+1)$-regular graph on $rn$ vertices. If $G$ contains no $\gamma$-independent set of size $n$, then $G[\frac{1}{2}]$ admits a $K_r$-factor with probability at least $\frac{1}{r}-o(1)$.
\end{theorem}

Note that an essentially necessary condition for the existence of a $K_r$-factor in a graph is that the order of the graph is divisible by $r$. % is the divisibility of $|G[\frac{1}{2}]|$ by $r$, which means 
It follows that $\mathbb{P}\big[G[\frac{1}{2}]~\text{admits a}~K_r\text{-factor}\big]\le \frac{1}{r}$. A result of Gan, Han and Hu~\cite{gan2024} (see Theorem~\ref{thm:Non-extremal case r>3}) shows that under a weaker degree condition $\delta(G)\geq (r-1)n - o(n)$, if $G$ contains no $\gamma$-independent set, then $G$ admits a $K_r$-factor. The non-extremal case essentially follows from this theorem because, with high probability, the required minimum degree and the non-existence of sparse sets would be inherited in $G[\frac{1}{2}]$. The major task is to deal with the extremal case as follows.


\begin{theorem}[Extremal case]\label{lem:extremal-case}
Suppose that $r\geq 3$, $n\in \mathbb{N}$   and $\frac{1}{n}\ll \gamma \ll \frac{1}{r}$. Let $G$ be an $((r-1)n+1)$-regular graph on $rn$ vertices which contains a $\gamma$-independent set of size $n$. Then  $G[\frac{1}{2}]$ admits a $K_r$-factor with probability at least $\frac{1}{(40r^2)^r}$.
\end{theorem}
In the extremal case, the existence of a $\gamma$-independent $n$-set in $G$ allows us to construct a partition $\mathcal{P}^0 = \{A_1^0, \dotsc, A_s^0, B^0\}$ of $V(G)$ for some $s\in [r]$, which is simply obtained by iteratively choosing disjoint $\gamma$-independent $n$-sets $A_i^0$ so that $B^0$ contains no $\gamma$-independent set {of size $n$}. 
By suitably reassigning some vertices, $\mathcal{P}^0$ can be refined into a \textit{good partition} of $G$ (see Definition~\ref{def:good partition}).  We first define the following three types of vertices with respect to a partition of $V(G)$.


\begin{definition}\label{def:bad vertex for Ai}
Let $r,t,n\in \mathbb{N}$ and $\alpha>0$. Suppose that $G$ is a graph on $rn$ vertices and   $\mathcal{P}=\{V_1,\ldots,V_t\}$ is a partition of $V(G)$. For each $i\in [t]$, we say that 
\begin{itemize}
    \item a vertex $v\in V_i$ is $(\alpha, V_i)$-good if $d(v,V_i)\le \alpha n$;
    \item a vertex $v\in V(G)\setminus V_i$ is $(\alpha, V_i)$-bad if $d(v,V_i)\le \alpha n$;
    \item a vertex $v\in V_i$ is $(\alpha,V_i)$-exceptional if $d(v,V_j)\le |V_j|-\alpha n$ for some $j\in [t]\setminus \{i\}$.
\end{itemize} 
\end{definition}
%For this, we introduce 
We now introduce the following definition of our desired partition.
\begin{definition}[Good partition]\label{def:good partition}
For $r,s,n\in \mathbb{N}$ with $s\le r$, and positive constants $\alpha, \beta,\beta', \gamma$, suppose that $G$ is a graph on $rn$ vertices and $\mathcal{P}=\{A_1,\ldots,A_s,B\}$ is a partition of $V(G)$. We say that $\mathcal{P}$ is an \text{$(\alpha,\beta,\beta',\gamma)$-good partition} of $G$ if all of the following hold: 
    \begin{enumerate}
[label =\rm  (A\arabic{enumi})]
    \item\label{B1} $ |A_i|\leq n+\alpha n$ for each $i\in [s]$, and $(r-s)n-r\alpha n\leq|B|\leq (r-s)n+r\alpha n$; % if $s\leq r-2$ and $n-r\alpha n\leq |B|\leq n+r\alpha n$ if $s=r-1$ and $B=\emptyset$ if $s=r$;
    \item\label{B2}  if $|A_i|>n$, then $\Delta(G[A_i])\leq \beta' n$ and the matching number of $G[A_i]$ is at least ${|A_i|-n+r}$; if $|A_i|=n$, then $G[A_i]$ is not empty;
    \item\label{B03} the number of $(\alpha^{1/5},A_i)$-good vertices is at least $|A_i|-2\alpha n$ for each $i\in[s]$; 
    \item\label{B04} for each $i\in [s]$ and each vertex $v\in V(G)\setminus A_i$, $d(v,A_i)\geq \beta n$;
    \item\label{B05} $d(v,B)\geq \beta n$ for each vertex $v\in V(G)\setminus B$ and $\de(G[B])\ge (r-s-1)n-r\alpha n$;
    \item\label{B06} the number of $(\alpha^{1/5},B)$-exceptional vertices is at most $r\alpha n$;

    %\item\label{B6} $d(v,D)\geq |D|-2\alpha^{1/5} n$ for every $(\alpha^{1/5},A_i)$-good vertex $v\in A_i$ with $A_i\neq D$;
    \item\label{B07} %\sout{$A_i$ is a $\gamma$-independent set for each $i\in [s]$, and} 
    $B$ has no $\gamma$-independent set of size at least $\frac{|B|}{r-s}$.
    %\item\label{B07} if $s=r-2$, then  for any two disjoint vertex sets of size at least $n-\beta n$ in $B$, there is at least one edge between them.
\end{enumerate}
\end{definition}


%is much easier to handle. 
Lemma~\ref{lem:construct good-partition} tells us that $G$ admits a good partition, say  $\mathcal{P}=\{A_1,\ldots,A_s,B\}$. Let $S\subseteq V(G)$ be the random subset formed by including each vertex of $V(G)$ independently with probability $\frac{1}{2}$, and let $\mathcal{P}'=\{A_1\cap S,\ldots,A_s\cap S,B\cap S\}$.
{We call $A_i\cap S$ a \textit{large} part of $\mathcal{P}'$ if $|A_i\cap S|>\frac{|S|}{r}$ and $B\cap S$ a \textit{large} part if $|B\cap S|>\frac{r-s}{r}|S|$. Otherwise the part is \textit{small}.}
Our proof proceeds in the following two steps.

\medskip
{\bf Step 1.} Prove that $\mathcal{P}'$ is a good partition of $G[S]$ with probability at least $\frac{1}{(40r^2)^r}$.
\medskip


%Based on Lemma \ref{lem:covering bad vertices for balance}, it suffices to show that $\mathcal{P}'$ is a good partition of  $G[S]$ with positive  probability. 

%We guarantee that $\mathcal{P}'$ is a good partition by employing the following two steps.

It suffices to show that $\mathcal{P}'$ satisfies~\ref{B1}-\ref{B07}. By Chernoff bound and concentration, conditions~\ref{B03}-\ref{B07} hold with high probability.
We therefore turn to the verification of~\ref{B2}, that is, every large part $A_i\cap S$ of $\mathcal{P}'$ contains a matching of size at least $|A_i\cap S|-\frac{|S|}{r}+r$. The strategy here is to restrict the size of each part  to a certain interval (make sure this happens with a constant probability), so that {$\mathcal{P}'$ satisfies~\ref{B1} and} any possible \( A_i \) with insufficient matching edges in \( G \) would inevitably turn into a small part of \(\mathcal{P}'\).
%The key idea is to estimate the probability that each part lies within a certain size interval, so that any part containing few matching edges being a small part in $\mathcal{P}'$. %is to estimate  the probability corresponding to the size of each part in $\mathcal{P}'$, where this size ensures that any part with few matching is \textit{small} in $\mathcal{P}'$.
% Thus, our main goal is to show that every \textit{large} part $A_i\cap S$ of $\mathcal{P}'$ contains a  matching of size $|A_i\cap S|-\frac{|S|}{r}$. The key idea is to control the size of parts with a large matching, thereby ensuring that any part with relatively few matching edges to be a \textit{small} part of $\mathcal{P}'$. 
%Since there are only a few low-degree vertices, we can reassign each of them to a designated part such that every vertex in the resulting partition has a high degree to the other parts. For a \textit{large} part, we can bound its maximum degree by moving some high-degree vertices to \textit{small} parts.

%Our primary objective is to find a $K_r$-factor in $G[\frac{1}{2}]$ where $G$ has a $\gamma$-independent set of size $n$. In order to obtain a clear picture of $G$, our first step is to find as many vertex-disjoint sparse sets as possible. Let $s$ be the maximum number of sparse sets with $1\le s\le r$ and $A_1^0,A_2^0\cdots,A_s^0$ are pairwise disjoint $\gamma$-independent sets each of size $n$ such that $B^0:=G\setminus \bigcup_{j=1}^{s}A_j^0$ has no sparse set. Let $\mathcal{P}_0=\{A_1^0,\cdots,A_s^0,B^0\}$ be an obtained partition of $V(G)$ which can only capture the rough structure of $G$. We conclude that there are very few vertices $v\in A_i^0$ with low degree to some part $A_j^0$ or $B^0$ where $j\neq i$.
%%we categorize the vertices of each part by their degrees to other parts and their internal degree (see Definition~\ref{def:bad vertex for Ai}), 
%We deal with such vertices by moving them to $A_j^0$ such that every vertex has high degree to all other parts and obtain a \emph{\textit{good}} partition (see Definition~\ref{def:good partition} and Lemma~\ref{lem:construct good-partition}), denoted as $\mathcal{P}=\{A_1,\cdots,A_s,B\}$.% by moving only a small number of vertices in $\mathcal{P}_0$.
%{(Yang: mainly highlight the idea, use less notations if possible)} 
%Note that there exists a small number of $(\beta, B)$-exceptional vertices in $B$.


%Let $G[\frac{1}{2}]:=G[S]$, where the random subset $S\subseteq V(G)$ is formed by including each vertex of $V(G)$ independently with probability $\frac{1}{2}$. Then we define a partition $\mathcal{P}'$ of $S$ by restricting the original partition to $S$, that is, $\mathcal{P}'=\{A_1\cap S,\cdots,A_s\cap S,B\cap S\}$.
% By chernoff bound, $\mathcal{P}'$ inherits the degree distributions of $\mathcal{P}$ with high probability and we infer that $\mathcal{P}'$ satisfies the conditions of Definition~\ref{def:good partition} except for~\ref{B2}.
% However, $\mathcal{P}'$ may be unbalanced.

\medskip
{\bf Step 2.} Prove that $G':=G[S]$ has a  $K_r$-factor if $G'$ admits a good partition (see Lemma~\ref{lem:covering bad vertices for balance}).\medskip

In doing this, we employ Lemma~\ref{lem:balance} and for that we shall do some preliminary cleaning and contraction to leave behind an auxiliary graph together with an almost balanced partition satisfying~\ref{A1}-\ref{A2}, to which we then apply Lemma~\ref{lem:balance}. The proof of Lemma~\ref{lem:balance} is, in turn, based on a multipartite Hajnal-Szemer\'{e}di theorem (see Lemma~\ref{lem:HS thm in multipartite graph} and we will elaborate more details in Section~5).
%For the proof of Lemma~\ref{lem:covering bad vertices for balance}, we will make use of Lemma \ref{lem:balance}, which forces us to  construct an auxilary graph satisfies all of \ref{A1}-\ref{A3}. %  which  a partition satisfies all  we will make use of a result %by Gan, Han, and Hu 
%concerning $K_r$-factors in balanced $r$-partite graphs (see Lemma~\ref{lem:HS thm in multipartite graph}). 
%Under the given good partition, we give a sufficient condition to construct a balanced multi-partite graph using the matching in \textit{large} part.
%In fact, the cliques within large parts can all contribute to making the partition $\mathcal{P}'$ balanced.
% The application of Lemma~\ref{lem:balance} proceeds with the following three preliminary steps.
{To be more specific, building on a good partition $\mathcal{P}'$, we construct an auxiliary graph satisfying~\ref{A1}-\ref{A2} in  the following three steps.}
\begin{itemize}
    \item {\bf  Cleaning bad or exceptional vertices:} We first clean all vertices failing the degree condition of~\ref{A3} (i.e., bad or exceptional vertices) by using vertex-disjoint copies of $K_r$. % in $G'$. % (i.e., those failing the degree condition of \ref{A3}). 
    Note that there are only $o(n)$ such vertices, far fewer than the number of good neighbors that each such vertex has in every $A_i\cap S$. Moreover,~\ref{B05} ensures the existence of linearly many vertex-disjoint copies of $K_{r-s}$ in $G'[B\cap S]$. % within the common neighborhood of all previously selected vertices and that each vertex in $B\cap S$ lies in linearly many disjoint copies of $K_{r-s}$.
    Thus, one may find a  $K_r$-tiling $\mathcal{K}$ covering all bad or exceptional vertices.
    %This can be easily achieved by \ref{B04}-\ref{B05} as there are only $o(n)$ many such vertices, combined with the fact that each vertex is contained in linearly many disjoint copies of $K_r$. % that is orders of magnitude larger-specifically, linear in $n$.
%Denote by $\mathcal{T}$ the resulting $K_r$-tiling obtained in this step. 
Furthermore, this tiling $\mathcal{K}$ can be chosen such that $(r-s)|V(\mathcal{K})\cap A_i|=|V(\mathcal{K})\cap B|$ for each $i\in [s]$.

\item {\bf  Fixing divisibility and balancing:} Let $G'':=G'-V(\mathcal{K})$, $A_i':=(A_i\cap S)\setminus V(\mathcal{K})$ for each $i\in [s]$ and $B':=(B\cap S)\setminus V(\mathcal{K})$.
The worst-case scenario is when the size of $B'$ is not divisible by $r-s$ or too small (i.e $\frac{|B'|}{r-s}<\frac{|G''|}{r}$). To overcome this, we will separately pick up two distinct collections of $r$-cliques, denoted as  $\mathcal{H}$ and $\mathcal{R}$, %\todo{Rename $\mathcal{H}_1$ in the last proof},
such that each $K_r$ in $\mathcal{H}$  uses $r-s+1$ vertices from \(B'\); every \(K_r\) in $\mathcal{R}$ instead uses \(r-s-1\) vertices from \(B'\) and a matching edge from some large part $A_i'$. In this way, we can ensure that $(r-s)\mid |B'\setminus V(\mathcal{H}\cup\mathcal{R})|$ and moreover $b:=\frac{|B'\setminus V(\mathcal{H}\cup\mathcal{R})|}{r-s}- \frac{|G''-(\mathcal{H}\cup\mathcal{R})|}{r}\ge 0$. In doing this, we need to find a large number of disjoint copies of $K_{r-s+1}$ or $K_{r-s-1}$, and this follows from~\ref{B05} and the fact that $G'[B\cap S]$ contains no $o(1)$-independent set of size at least $\frac{|B\cap S|}{r-s}$. In fact, a supersaturation version of stability  result (see Lemma~\ref{SuperT}) implies that $G'[B\cap S]$ actually contains {$\Omega(n^{r-s+1})$ distinct} copies of $K_{r-s+1}$. Moreover, \ref{B2} ensures that each ${G'}[A_i\cap S]$ contains a matching of the required size.


    \item {\bf Contraction:} {%Recall that, for every $i \in [s]$, we have $|B\cap S| \approx (r-s)\,|A_i\cap S|$.
    To obtain an almost balanced partition as in~\ref{A1} and~\ref{A2},}  we will find a perfect tiling with $K_{r-s}$ (possibly need extra copies of $K_{r-s+1}$, denoted as $\mathcal{F}$) in $G'[B'\setminus V(\mathcal{R}\cup\mathcal{H})]$ and contract each copy of $K_{r-s}$ (resp. $K_{r-s+1}$ in $\mathcal{F}$) into a vertex (resp. an edge), yielding a new vertex set, denoted as $B^*$ whose size is close to $|A_i'|$ for every $i\in[s]$. In this case, ~\ref{A2} can be easily achieved by~\ref{B2}, and in particular, the matching number in $B^*$ can be achieved by the construction of $B^*$ and an additional constraint on $|\mathcal{F}|=(r-s)b$. Indeed, such a perfect tiling can be constructed by first taking $\mathcal{F}$ (a few number of $K_{r-s+1}$) from $B'\setminus V(\mathcal{H}\cup \mathcal{R})$ and then finding a $K_{r-s}$-factor in the remaining set in two cases depending on whether $r-s=2$ or not. For $r-s\geq 3$, a desired $K_{r-s}$-factor can be obtained from Theorem~\ref{thm:Non-extremal case r>3}; for $r-s=2$, a perfect matching can also be found by Lemma~\ref{thm:r-s=2} after a minor adjustment towards $\mathcal{K}\cup \mathcal{H}\cup \mathcal{F}$. 
    
    
    
    
    
    
    
    %Find a $K_{r-s}$-factor in $G'[B\cap S]-V(\mathcal{T})]$:}
    %\todo{How to achieve a $1:1:1:\cdots (r-s)$-balanced partition??} %One may assume that $(r-s)$ divides $|B\cap S|-|V(\mathcal{T})|$. Indeed, a
%For this, we would like to find a large  matching in every large part.  By~\ref{B2}, each large part of $\mathcal{P}'$ contains a matching of size at least $|A_i\cap S| - \frac{|G'|}{r}+r$. 


 %, which can be extended to a $K_r$-tiling of $G'-V(\mathcal{T})$, say $\mathcal{K}$. 
 % of which can be extend %It follows that $(r-s)\mid |(B\cap S)\setminus V(\mathcal{T}\cup \mathcal{K})|$. %{By removing a constant number of disjoint copies of $K_{r-s+1}$ (back into $\mathcal{T}$ ??)}, we can achieve the divisibility. 

 
%Now suppose $\frac{|B'\setminus V(\mathcal{R})|}{r-s}>\frac{|G''-\mathcal{R}|}{r}$.  We can select at least \(|B'\setminus V(\mathcal{R})| - \frac{(r-s)|G''-\mathcal{R}|}{r}\) vertex‑disjoint copies of \(K_{r-s+1}\) in \(G''[B']\), say $\mathcal{H}$,  in $G''[B']$  (vertex disjoint from $\mathcal{K}\cup \mathcal{R}$), so that %Suppose $\frac{|B'|}{r-s} > \frac{|G''|}{r}$. Then we can pick at least $|B'| - \frac{(r-s)|G''|}{r}$ vertex-disjoint copies of $K_{r-s+1}$, say $\mathcal{H}$,  in $G''[B']$ % and then extend them to a $K_r$-tiling $\mathcal{H}$ (vertex disjoint from $\mathcal{K}$), so that  {$(r-s)\mid |B'\setminus V(\mathcal{R}\cup\mathcal{H})|$ and $\frac{|B'\setminus V(\mathcal{R}\cup \mathcal{H})|}{r-s}=\frac{|G''-(\mathcal{R}\cup\mathcal{H})|}{r}$.}  We then contract each such copy in $\mathcal{H}$ into an edge.  For the subgraph induced by the remaining vertices of $B'$, we claim that it contains a $K_{r-s}$-factor, where each copy of $K_{r-s}$ will be contracted to a single vertex.  Indeed, for $r-s\geq 3$, a desired $K_{r-s}$-factor can be obtained from Theorem~\ref{thm:Non-extremal case r>3}; for $r-s=2$, a perfect matching can also be found by Lemma~\ref{thm:r-s=2} after a minor adjustment towards $\mathcal{K}\cup \mathcal{H}$. 
% to $V(\mathcal{T})$. 
%Note that we already obtain an auxiliary graph with $s+1$ parts and a minimum degree condition.

%\medskip
%Let $\mathcal{P}'' := \{A_1', \ldots, A_s', B'\}$ be the vertex partition of $G'' := G - V(\mathcal{T} \cup \mathcal{K})$ obtained from $\mathcal{P}'$ by removing all vertices of $\mathcal{T} \cup \mathcal{K}$ from their corresponding parts. By Definition~\ref{def:good partition}~\ref{B2}, each large part of $\mathcal{P}''$ contains a matching of size at least $|A_i'| - \frac{|G''|}{r}$, and $G''[B]$ contains $|B'| - \frac{(r-s)|G''|}{r}$ disjoint copies of $K_{r-s+1}$ (since $G[B]$ contains no $o(1)$-independent set of size $n$).

%\item {\bf Extra matching edges~\ref{A2}}. It remains to verify that every large part of the final partition has a large matching as required. Let $G^*$ be a graph with vertex set $(V(G'')\setminus B') \cup B^*$, and $uv\in E(G^*)$ if and only if $u,v\in V(G'')\setminus B'$ or $u\in V(G'')\setminus B'$, $v\in B^*$ and $u$ is adjacent to all vertices of the clique that was contracted to form $v$ in $G''$. Now, in $G^*$,~\ref{A2} can be easily achieved by \ref{B2} and by the construction of $B^*$. Consequently, Lemma~\ref{lem:balance} implies that $G^*$ has a $K_{s+1}$-factor, and hence $G'$ admits a $K_r$-factor. 
\end{itemize}





\medskip

\noindent{\bf Organization:} {In Section~\ref{sec:non-extremal} we consider the non-extremal case, establishing Theorem~\ref{thm:Non-extremal case}.
Section~\ref{sec:extremal} deals with the extremal case, i.e., Theorem~\ref{lem:extremal-case}.
The proof of Theorem~\ref{lem:extremal-case} splits into two parts: first we show that \(G\) admits a good partition (Lemma~\ref{lem:construct good-partition}), and then we prove that every graph possessing such a partition contains a \(K_r\)-factor (Lemma~\ref{lem:covering bad vertices for balance}).
The proofs of Lemma~\ref{lem:construct good-partition} and Lemma~\ref{lem:covering bad vertices for balance} are provided in Section~\ref{sec:lemma}. Some concluding remarks are given in the last section.}


\section{Non-extremal case}\label{sec:non-extremal}

This section is devoted to the proof of Theorem~\ref{thm:Non-extremal case}, which concerns the case when the host graph $G$ contains no large sparse set. The existence of a $K_r$-factor in such a graph $G$ was established by Gan, Han, and Hu~\cite{gan2024}, under a slightly weaker minimum degree condition.
%In this section, we prove Theorem \ref{lem:extremal-case}, which considers the case when the host graph $G$ has no large independent set. 
%\begin{definition}[$\gamma$-independent set]
%Given an $n$-vertex graph $G$ and a constant $\gamma>0$, a vertex subset $U\subseteq V(G)$ is called a $\gamma$-independent $k$-set if $e(G[U])<\gamma n^2$ and $|U|=k$.
%\end{definition}
%We divide the proof into two cases based on whether $G$ contains a $\gamma$-independent set of size $\frac{|G|}{r}$ (the \textit{extremal case}) or not (the \textit{non-extremal case}). 
%For the non-extremal case, the existence of a $K_r$-factor in $G$ was established by Gan, Han and Hu~\cite{gan2024}.

\begin{theorem}[\cite{gan2024}]\label{thm:Non-extremal case r>3}
Suppose $r\ge 3$ and $\frac{1}{n}\ll\alpha\ll\gamma\ll\frac{1}{r}$. Let $G$ be an $rn$-vertex graph with $\de(G)\ge (r-1-\alpha)n$. If $G$ has no $\gamma$-independent set of size $n$, then there is a $K_r$-factor in $G$.
\end{theorem}

Next, we introduce the Frieze-Kannan Regularity Lemma~\cite{Frieze}, which will be used in the proof of Theorem~\ref{thm:Non-extremal case}.%, which will be used to show that  w.h.p., the  bidensenessis inherited by the random induced subgraph $G[\frac{1}{2}]$. 
\begin{lemma}[\cite{Frieze}]\label{partition}
    For any $\eps>0$, there are $T,n_0>0$ such that the following holds for all $n\geq n_0$. Every $n$-vertex graph $G$ admits a vertex partition $\{V_1,\ldots,V_t\}$  with $t\leq T$ and $|V_1|\leq \cdots\leq |V_t|\leq |V_1|+1$ satisfying for any $A,B\subseteq V(G)$ we have 
    $$
    \Big|e(G[A,B])-\sum_{1\le i,j\le t}\frac{|A\cap V_i||B\cap V_j|}{|V_i||V_j|}e(G[V_i,V_j])\Big|<\eps n^2.
    $$
\end{lemma}
\medskip


Now we are ready to present a short proof of Theorem~\ref{thm:Non-extremal case}.%Now we prove the following result to end this section.
%\begin{lemma}[Non-extremal case]
%Let $\gamma>0$ and $G$ be an $((r-1)n+1)$-regular graph on $rn$ vertices. Suppose that $G$ contains no $\gamma$-independent $n$-set. Then $G[\frac{1}{2}]$ admits a $K_r$-factor with probability at least $\frac{1}{r}-o(1)$.
%\end{lemma}

\begin{proof}[{\bf Proof of Theorem~\ref{thm:Non-extremal case}}]
Choose $r\geq 3,n\in \mathbb{N}$ and a constant $\gamma$ satisfying $\frac{1}{n}\ll\gamma\ll \frac{1}{r}$. Let $G$ be an $((r-1)n+1)$-regular graph on $rn$ vertices such that $G$ contains no $\gamma$-independent set of size $n$. Let $\{V_1,\ldots,V_t\}$ be a partition of $V(G)$ obtained from Lemma~\ref{partition} applied with $\eps=\frac{1}{10}\gamma$. Consider $G[\frac{1}{2}]=G[S]$, where $S$ is a random subset of $V(G)$ with $r\mid |S|$. Note that $r \mid |S|$ is a necessary condition for $G[S]$ to contain a $K_r$-factor, which occurs with probability $\frac{1}{r}$. By Chernoff bound, {w.h.p.}~$|S|=\frac{r}{2}n\pm n^{0.6}$, $|S\cap V_i|= \frac{rn}{2t}\pm n^{0.6}$ for each $i\in [t]$,  and every vertex has degree $\frac{r-1}{2}n\pm n^{0.6}$ in $G[S]$. 
%Our goal is to prove that 
Conditioned on these events,  we prove the following claim. 
\begin{claim}\label{no-independent}
    $G[S]$ contains no $\frac{\gamma}{25}$-independent set of size $\frac{|S|}{r}$. 
\end{claim}

%By Chernoff bound, $|S\cap V_i|= \frac{rn}{2t}\pm n^{0.6}$. 
\begin{proof}
Given any $A'\subseteq S$ with $|A'|= \frac{|S|}{r}$, we are to show that $e(G[A'])\ge \frac{\gamma}{25} n^2$. 
Construct $A\subseteq V(G)$ %with $|A|\sim n$ 
so that {$|A\cap V_i|=2|A'\cap V_i|\pm 2n^{0.6}$} for all $i\in [t]$. Notice that $|A|=n\pm 3tn^{0.6}$ and $e(G[A])\geq \frac{4}{5}\gamma n^2$ as $G$ itself contains no $\gamma$-independent set of size $n$. Lemma~\ref{partition} implies that
\begin{itemize}
    \item $e(G[A])$ differs by at most $\frac{\gamma}{10}n^2$ from $\sum:=\sum_{i,j}\frac{|A\cap V_i||A\cap V_j|}{|V_i||V_j|}e(G[V_i,V_j])$;
    \item $e(G[A'])$ differs by at most $\frac{\gamma}{10}n^2$ from $\sum_{i,j}\frac{|A'\cap V_i||A'\cap V_j|}{|V_i||V_j|}e(G[V_i,V_j])\geq {\frac{1}{5}}\sum$. %{(make it more precise!)}.
\end{itemize}
{{It follows that 
$$
e(G[A'])\geq \frac{1}{5}\sum-\frac{1}{10}\gamma n^2\geq \frac{1}{5}\Big(e(G[A])-\frac{\gamma}{10} n^2\Big)-\frac{\gamma}{10}n^2\geq \frac{7}{50}\gamma n^2-\frac{\gamma}{10} n^2= \frac{\gamma}{25}n^2. 
$$}} %{(How? add more details)}. 
That is, $G[S]$ contains no $\frac{\gamma}{25}$-independent set of size $\frac{|S|}{r}$. 
\end{proof}
Combined with the condition $\de(G[S])\geq \frac{r-1}{2}n- n^{0.6}$, Theorem~\ref{thm:Non-extremal case r>3} implies that $G[S]$ contains a $K_r$-factor, which happens with probability at least $\frac{1}{r}-o(1)$. 
%By Chernoff bound, $|A\cap S|= \frac{n}{2}\pm n^{0.6}$ and $e(G[A\cap S])=\frac{\gamma}{4}n^2\pm n$. 
\end{proof}

\section{Extremal case}\label{sec:extremal}

For the extremal case, our first task is to establish a good partition of the graph $G$, provided that it contains a large sparse set. 
\begin{lemma}\label{lem:construct good-partition}
    Let $r,n$ be positive integers with $r\ge 3$, and choose $\frac{1}{n}\ll \gamma\ll\alpha\ll\beta\ll\gamma'\ll \frac{1}{r}.$ {Let $G$ be an $rn$-vertex graph with $\delta(G)\geq (r-1)n+1$.} Suppose that $G$ contains a $\gamma$-independent set of size $n$. Then there exists an $(\alpha,\beta,\beta,\gamma')$-good partition $\mathcal{P}=\{A_1,\cdots,A_s,B\}$ of $G$ for some $s\in [r]$. %Furthermore, we have that

    %\begin{enumerate}[label =\rm  (B2')]
        %\item\label{B2'}  $\Delta(G[A_i])\leq \beta n$ and the matching number of $G[A_i]$ is at least $\frac{|A_i|-n}{4\beta}$ if $|A_i|>n$; 
        %\end{enumerate}
        %\begin{enumerate}[label =\rm  (B6')]
        %\item\label{B6'}  the number of $(\frac{\beta}{2},B)$-exceptional vertices is at most $r\alpha n$;
        %\end{enumerate}
        %\begin{enumerate}[label =\rm  (B7')]
%\item\label{b7'}   $d(v,D)\geq |D|-\alpha^{1/5} n$ for every $(\alpha^{1/5},A_i)$-good vertex $v\in A_i$ with $A_i\neq D$.
    %\end{enumerate}
%\begin{enumerate}[label =\rm  (B7')]
    %\item\label{B7'} $A_i$ is a $\gamma'$-independent set for each $i\in [s]$, and $B$ has no $\gamma'$-independent set of size $n$.
    %\end{enumerate}
\end{lemma}

Given a good partition of $G$, one can expect that $G[S]$ also admits a good partition. If so, then the following lemma would ensure the existence of a $K_r$-factor in $G[S]$.

\begin{lemma}\label{lem:covering bad vertices for balance}
Let $r,n$ be positive integers with $r\ge 3$ and choose $\frac{1}{n}\ll\alpha\ll\beta',  \beta\ll\gamma\ll \frac{1}{r}$. Let $G$ be an $rn$-vertex graph with {$\de(G)\geq (r-1)n-\alpha n$.}  %\sout{$(r-1)n-\alpha n\leq \de(G)\leq \Delta(G) \leq (r-1)n+\alpha n$}. 
If $G$ admits an $(\alpha,\beta,{\beta'},
\gamma)$-good partition, then $G$ has a $K_r$-factor. % partition $\mathcal{P}$ can be transformed into a balanced one by removing a $K_r$-tiling of seize at most ?? vertices.
\end{lemma}
We will prove Theorem~\ref{lem:extremal-case} using Lemmas~\ref{lem:construct good-partition} and~\ref{lem:covering bad vertices for balance}, whose proofs are postponed to the next section.
\subsection{Proof of Theorem~\ref{lem:extremal-case}}
In this subsection we prove Theorem~\ref{lem:extremal-case}, which establishes Theorem~\ref{thm:the random version of main thm} for graphs containing a sparse set of size $n$. 
%In this section we prove our main result Lemma~\ref{lem:extremal-case} using the sufficient conditions for $K_r$-factors in Lemma~\ref{lem:covering bad vertices for balance}. 
As the size of a random subset of $A$ is distributed as the binomial $B(|A|, 1/2)$, we shall use the following normal approximation, which is a consequence of the Central Limit Theorem.
%As the size of a random subset of a set $A$ has the binomial $B(|A|, 1/2)$ distribution, we need the following observation on the difference of independent binomials.



\begin{lemma}[\cite{Keevash2025}]\label{lem-int}
If $X\sim B(n,1/2)$, then $\mathbb{P}(a\sqrt{n}/2\leq X-n/2\leq b\sqrt{n}/2)=\int_a^b \frac{1}{\sqrt{2\pi}}e^{-t^2/2}\,dt+o(1)$.
\end{lemma}
We will also use the following version of Chernoff bound.
\begin{lemma}[Chernoff bound]
    Let $X$ be a binomial random variable and $\delta>0$. Then 
    $$
    \mathbb{P}[|X-\mathbb{E}X|\geq \delta \mathbb{E}X]\leq 2\exp(-\delta^2\mathbb{E}X/(2+\delta)).
    $$
\end{lemma}
%Motivated by a problem from statistical mechanics, Fortuin, Kasteleyn and Ginibre~\cite{} prove the following useful inequality, which has become known as the FKG Inequality. We introduce a version formulated for independent binomial distributions.
%\begin{lemma}[The FKG Inequality~\cite{}]\label{The FKG Inequality}
%For any two increasing events $A$ and $B$ where $A$ and $B$ follow binomial distributions, we have $\mathbb{P}[A\cap B]\ge \mathbb{P}[A]\cdot \mathbb{P}[B]$.
%\end{lemma}

%We also require the following standard fact concerning the difference of independent binomial random variables. 
%\begin{lemma}[\cite{Keevash2025}]\label{lem:binomial distribution}
%If random sets $X\sim B(n,1/2)$ and $Y\sim B(m,1/2)$ are independent, then $X+m-Y\sim B(n+m,1/2)$.
%\end{lemma}


%The following obvious fact illustrates that minimum degrees can be preserved with high probability in the random vertex samples.
%\begin{fact}
%Given a bipartite graph $G=(A_1,A_2)$ on $2n$ vertices, if $\de(G)\ge \beta n$, then $\de(G[\frac{1}{2}])\ge \frac{\beta n}{4}$ with probability $1-o(1)$.
%\end{fact}
%\begin{proof}
%Consider $G[\frac{1}{2}]=G[S]$ where $S$ is a random subset of $V(G)$. For every vertex $v \in S$, since $d_G(v) \ge \beta n$, we have $\mathbb{E}[d_{G[S]}(v)]\ge \frac{1}{2}\cdot \beta n$. By chernoff bound, w.h.p.~we have $d_{G[S]}(v)\ge \frac{\beta n}{4}$ for every $v \in S$.
%\end{proof}

%The required matchings in Lemma~\ref{good-partition} will be provided by the following two lemmas. The first claim gives that the matching in each large part can be preserved with high probability after the random vertex sampling, which plays an essential role in this paper.
%\sout{Our next lemma establishes that, with high probability, the degree condition and the number of matching edges in a good partition are both inherited by a random subset, which plays an essential role in our proof.}

%The required matching in Lemma~\ref{lem:extremal-case} will be provided by the following two lemmas.
%\begin{lemma}\label{lem:inherit matching}
%\sout{For any part $A_i$ with $|A_i|=\frac{n}{r}+k_i$, if $k_i>0$ then the random set $A_i \cap S$ admits a matching of size at least $\frac{k_i}{32\beta}$ with probability $1-o(1)$.}
%\end{lemma}
%\begin{proof}
%Note that $G[A_i]$ contains a matching $M_i$ of size at least $\frac{k_i}{4\beta}$.
%Let $X=\sum_{e=\{uv\}\in M_i}X_e$ where $$X_e= 
%\begin{cases} 
%1, & u\in A_i\cap S~\text{and}~v\in A_i\cap S \\
%0, &~\text{otherwise}
%\end{cases}$$ is an indicator variable.
%Note that $X\sim B(|M_i|,1/2)$.
%As $\mathbb{P}[X_e=1]=\mathbb{P}[u\in A_i\cap S]\cdot \mathbb{P}[v\in A_i\cap S]=\frac{1}{4}$, we have
%\[
%\mathbb{E}[X]=\sum_{e\in M_i}\mathbb{E}[X_e]=\frac{1}{4}|M_i|\ge \frac{k_i}{16\beta}.
%\]
%By chernoff bound, as $\beta\ll\frac{1}{r}$, we have
%\[
%\mathbb{P}[X\le \frac{1}{2}\mathbb{E}[X]]\le2\exp\left(-\frac{k_i}{160\beta}\right)=o(1),
%\]
%which implies that there exists a matching of size at least $\frac{k_i}{32\beta}$ in $A_i\cap S$ with probability $1-o(1)$.
%\end{proof}
%We note that a good partition necessitates a large matching in each large part. The next lemma guarantees the existence of a subset with size $n$ inducing a large matching in an $((r-1)n+1)$-regular graph on $rn$ vertices, whose proof will be given in the next section.  %
%{Given a partition $\mathcal{P}:=\{V_1,\ldots,V_t\}$ of $G$, we call $\mathcal{P}$ balanced if $|V_1|=\cdots=|V_t|$.} 
The following lemma ensures that for any  $rn$-vertex $((r-1)n+1)$-regular graph $G$ with a balanced partition  $\{A_1,A_2,\ldots,A_r\}$ (i.e., $|A_1|=\cdots=|A_r|=n$), there exists a large matching within some $G[A_i]$. The proof will be given in the next section.  % a subgraph induced by a set of size $n$. %  gives that for a balanced partition, there is a matching in some part, which uses regularity.

\begin{lemma}\label{vertex-cover}
Let $G$ be an $((r-1)n+1)$-regular graph on $rn$ vertices and $\{A_1,A_2,\ldots,A_r\}$ be a balanced partition of $V(G)$. Suppose that $C_i$ is a minimum vertex cover of $G[A_i]$ for each $i\in [r]$. Then $\max\{|C_i|:i\in [r]\}\geq \frac{1}{r}\sqrt{n}$.
\end{lemma}
%The following simple remark  will be used in our proof.

\begin{rmk}\label{remark:vertex-cover}
Let $M$ be a maximal matching in a graph $H$. Then $V(M)$ is a vertex cover in $H$. In particular, $|M| \geq |C|/2$ where $C$ is a minimum vertex cover {in $H$}. 
If moreover $H$ has maximum degree $\Delta$,  %\sout{ and $C$ is a vertex cover in $H$}. 
then $e(H) \leq \Delta |C|$.
%Let $M_i$ be a maximum matching of $A_i$. By Vizing's theorem, we have
%\begin{align*}
%|M_i|\ge \frac{e(G[A_i])}{\Delta(G[A_i])+1}\ge \frac{|A_i|\cdot (k_i+1)}{2(\beta n+1)}\ge \frac{k_i}{2\beta}.
%\end{align*}
\end{rmk}

We conclude this section by giving the proof of Theorem~\ref{lem:extremal-case}.
\begin{proof}[{\bf Proof of Theorem~\ref{lem:extremal-case}}]
Choose $r\geq 3,n\in \mathbb{N}$ and constants satisfying $$\frac{1}{
n}\ll\gamma\ll \alpha\ll \beta\ll \delta\ll\gamma'\ll \frac{1}{r}.
$$
Assume that $G$ is an $((r-1)n+1)$-regular graph on $rn$ vertices and $G$ contains a $\gamma$-independent set of size $n$.
By Lemma~\ref{lem:construct good-partition}, $G$ admits  an $(\alpha,\beta,\beta,\gamma')$-good partition $\mathcal{P}:=\{A_1,\ldots,A_s,B\}$ for some $s\in [r]$, that is, $\mathcal{P}$ satisfies~\ref{B1}-\ref{B07}. %partition $\mathcal{P}:=\{A_1,\cdots,A_s,B\}$ of $V(G)$ satisfying~\ref{B1}-\ref{B07}.

Consider the random subset $S\subseteq V(G)$ obtained by including each vertex of $V(G)$ independently with probability $\frac{1}{2}$. Then, w.h.p.~we have {$\de(G[S])\geq \frac{(r-1)n}{2}-n^{0.6}$}. %By Chernoff bound, {w.h.p.}~$|S|=\frac{r}{2}n\pm n^{0.6}$, $|S\cap A_i|= \frac{n}{2}\pm r\alpha n$ for each $i\in [s]$ and $|S\cap B|= \frac{(r-s)n}{2}\pm r\alpha n$. 
Let $\mathcal{P}'=\{A_1\cap S,\ldots,A_s\cap S,B\cap S\}$. By Lemma~\ref{lem:covering bad vertices for balance}, in order to obtain a $K_r$-factor in $G':=G[S]$, it suffices to determine the probability that  $\mathcal{P}'$ is a  $(3\alpha,\frac{\beta}{4},3\beta,\frac{\gamma'}{25})$-good partition of $G'$ under the condition that $r\mid |S|$.  
%In what follows, we condition on the event that  \( r \mid |S| \), which occurs with probability \( \frac{1}{r} \) and is necessary for a \(K_r\)-factor in \(G[S]\).
%Note that $S\sim B(rn,1/2)$ and $|S|=\frac{rn}{2}\pm n^{0.6}$ with high probability.
%By Chernoff  bound,  %~$|S|=\frac{rn}{2}\pm n^{0.1}$, $d_{G[S]}(v)=(r-1)\frac{|S|}{r}\pm n^{0.1}$ for each $v\in S$. %, and %$|S\cap A_i|=\frac{|A_i|}{2}\pm n^{0.1}$ for each $i\in [s]$, and $|S\cap B|=\frac{|B|}{2}\pm n^{0.1}$. 
%Moreover, by chernoff bound, 
%we have the following: 
%\begin{claim}\label{claim:B1-B06 in S}
%Given constants $\frac{1}{n}\ll\gamma\ll\alpha\ll\beta\ll1$, there exists a constant $\gamma'\ll\frac{1}{r}$ such that the random 

\begin{claim}\label{claim-prob}
With high probability, $\mathcal{P}'$ satisfies the following properties:
\begin{enumerate}[label=\rmfamily (A\arabic*), start=3]
    %\item[{\rm (B1)}]\label{D1} $\frac{n}{2}-2r\alpha n\leq |A_i|\leq \frac{n}{2}+2r\alpha n$ for each $i\in [s]$ and $(r-s)\frac{n}{2}-2r\alpha n\leq |B|\leq (r-s)\frac{n}{2}+2r\alpha n$; %\item\label{D1}
    %$|A_i\cap S|=\frac{|A_i|}{2}\pm n^{0.1}$ for each $i\in [s]$ and $|B\cap S|=\frac{|B|}{2}\pm n^{0.1}$;%$|A_i\cap S|=\frac{n}{2}+ n^{0.1}$ if $|A_i|\le n$ and $|A_i\cap S|=\frac{|A_i|}{2}|\pm n^{0.1}$ if $|A_i|=n+k_i$ where $k_i>0$;
    %\item[{\rm (B2)}]\label{D2} $\Delta(G[A_i\cap S])\leq \beta n$ for each $i\in [s]$;%\todo{relabel $D_i$ as $B_i$}
    \item\label{D1} the number of $((3\alpha)^{1/5},A_i\cap S)$-good vertices is at least $|A_i\cap S|-{6\alpha \frac{|S|}{r}}$ for each $i\in [s]$; 
    \item\label{D3} $d(v,A_i\cap S)\geq {\frac{\beta}{4}\frac{|S|}{r}}$ for each vertex $v\in S\setminus A_i$ with $i\in [s]$;
    \item\label{D4} $d(v,B\cap S)\geq {\frac{\beta}{4}\frac{|S|}{r}}$ for each vertex $v\in S\setminus B$ and $\de(G[B\cap S])\ge {\frac{(r-s-1)|S|}{r}-3\alpha |S|}$;
    \item\label{D5} the number of $((3\alpha)^{1/5},B\cap S)$-exceptional vertices is at most ${3\alpha |S|}$;
    %\item\label{D6} for any two sets of size $\frac{|S|}{r}$ in $B\cap S$ there are at least $\frac{\beta n}{4}$ edges between them;
    %\item\label{D7} $d(v,D\cap S)\ge \frac{|D\cap S|}{2}-2\alpha^{1/5}n$ for every $(\alpha^{1/5},A_i)$-good vertex $v\in A_i\cap S$ where $A_i\cap S\neq D$ and $D\in \mathcal{P}'$;
    \item\label{D7} %\sout{$A_i\cap S$ is a $5\gamma_1$-independent set for each $i\in [s]$ and} 
    $B\cap S$ has no $\frac{\gamma'}{25}$-independent set of size  at least $\frac{|B\cap S|}{r-s}$.
    %\item[{\rm (B06)}]\label{D8} if $s=r-2$, then  for any two sets of size $n-\beta n$ in $B$, there is at least one edge  between them.
\end{enumerate}
\end{claim}

\begin{proof}%[Proof of Claim \ref{claim-prob}]
Note that $|S|\sim B(rn,1/2)$, $|A_i\cap S|\sim B(|A_i|,1/2)$ for all $i\in[s]$ {and $|B\cap S|\sim B(|B|,1/2)$}. By Chernoff bound, w.h.p.~$|S|=\frac{rn}{2}\pm n^{0.6}$, $|A_i\cap S|=\frac{|A_i|}{2}\pm n^{0.6}$ for each $i\in [s]$ and $|B\cap S|=\frac{|B|}{2}\pm n^{0.6}$.

(A3) %{Simplify this. It seems very trivial}. 
Let $v\in A_i$ be an $(\alpha^{1/5},A_i)$-good vertex in $G$ for some $i\in [s]$.  Recall that $d(v,A_i)\leq \alpha^{1/5} n$. %If $d(v,A_i)\leq (3\alpha)^{1/5}\frac{n}{3}$, then w.h.p. we have $$d(v,A_i\cap S)\leq d(v,A_i)\leq (3\alpha)^{1/5}\frac{|S|}{r}.$$  If $d(v,A_i)\geq (3\alpha)^{1/5}\frac{n}{3}$, t
Then by Chernoff bound, w.h.p.~we have $$d(v,A_i\cap S)\leq \frac{1}{2}\alpha^{1/5}n+ n^{0.6}\leq (3\alpha)^{1/5}\frac{|S|}{r}.$$ %Hence $d(v,A_i\cap S)\leq d(v,A_i)\leq \alpha^{1/5} n\leq (2^6\alpha)^{1/5}\frac{|S|}{r}$. 
Together with~\ref{B03}, there are at least $|A_i\cap S|-{2}\alpha n{\ge |A_i\cap S|-6\alpha\frac{|S|}{r}}$ vertices in $G'$ that are $((3\alpha)^{1/5},A_i\cap S)$-good.\medskip

(A4)-(A5) trivially follow from the Chernoff bound. \medskip

(A6) Let $v\in B$ be a vertex that is not $(\alpha^{1/5},B)$-exceptional in $G$. 
Then $d(v,A_j) > |A_j| - \alpha^{1/5} n$ for every $j \in [s]$. 
By Chernoff bound, w.h.p. we have
\[
d(v,A_j \cap S)\geq  \frac{|A_j| - \alpha^{1/5} n}{2} - n^{0.6} \ge |A_j \cap S| - (3\alpha)^{1/5}\frac{|S|}{r} \ \text{for all } j \in [s].
\]
It follows from~\ref{B06} that the number of $((3\alpha)^{1/5},B \cap S)$-exceptional vertices in $G'$ is at most $r\alpha n\le {3\alpha |S|}$. \medskip%\todo{maybe we can replace this $2\alpha$ with $3\alpha$}. 

(A7) %\sout{By \ref{B07}, w.h.p. we have that for each $i\in [s]$,}  $$e(G[A_i\cap S])\leq e(G[A])\leq \gamma_1 n^2\leq 5\gamma_1 \Big(\frac{|S|}{r}\Big)^2.$$ Thus, $A_i\cap S$ is a $5\gamma_1$-independent set of $G'$. 
By applying Claim~\ref{no-independent} with $(G[B],\gamma')$ in place of $(G,\gamma)$, we obtain that $G[B\cap S]$ contains no $\frac{\gamma'}{25}$-independent set of size at least $\frac{|B\cap S|}{r-s}$, as desired. 
%Notice that  each edge in $G$ is selected into $G[S]$ with probability $\frac{1}{4}$. 
%and for any set 
%Let $B'$ be an arbitrary subset of $B\cap S$ with size  $\frac{|B\cap S|}{r-s}$. Then there exists a vertex subset $B''$ of $B$ with size $\frac{|B|}{r-s}\pm n^{0.6}$ such that $B'\subseteq B''$.  there exists an \(n\)-vertex set \(B''\supseteq B'\) such that \todo{How to choose $B''$?, The following is too informal to believe.}
%$$e(G[B']){\geq ??}{\Big(\frac{1}{4}-{o(1)??}\Big)e(G'[B''])\geq \Big(\frac{1}{4}-o(1)\Big)\gamma_2n^2\geq \frac{\gamma_2}{2}\Big(\frac{|S|}{r}\Big)^2.$$ 
%Thus, w.h.p.~(B7) holds. {(Here we refer to the proof Thm2.1)}. 
\end{proof}
%Since $A_i$ is a $\gamma'$-independent set in $G$, we have $e(G[A_i\cap S])\leq (\frac{1}{4}+o(1))e(G[A])\leq (\frac{1}{4}+o(1))\gamma'n^2\leq 2\gamma'(|S|/r)^2$. Since $B$ contains no is a $\gamma'$-independent $n$-set in $G$, for any set $B'\subseteq B\cap S$ with size $\frac{n}{2}$, we have $e(G'[B'])\geq (\frac{1}{4}-o(1))\gamma'n^2\geq \frac{1}{2}\gamma'(|S|/r)^2$. %together with \ref{B7'}, it is easy to check that $\mathcal{P}'$ w.h.p.~satisfies~\ref{B07}.

%Since $\mathcal{P}$ satisfies~\ref{B03}, that is, there are at least $|A_i|-2\alpha n$ $(\alpha^{1/5},A_i)$-good vertices in each $A_i$, it follows that w.h.p.~the number of $(\frac{\alpha^{1/5}}{2},A_i\cap S)$-good vertices is at least $|A_i\cap S|-4\alpha n$. By the similar argument, we deduce that $\mathcal{P}'$ w.h.p.~satisfies~\ref{B03}-\ref{B07}. As each edge in $G$ is selected into $G[S]$ with probability $1/4$, it is easy to check that $\mathcal{P}'$ w.h.p.~satisfies~\ref{B07}.

Let $\mathbf{E}_0$ be the event that~\ref{D1}-\ref{D7} hold, and let $\mathbf{E}$ be the event that $\mathcal{P}'$ is a $(3\alpha,\frac{\beta}{4},3\beta,\frac{\gamma'}{25})$-good partition of $G'$. 
By Claim~\ref{claim-prob}, we have $\mathbb{P}[\mathbf{E}_0]\geq 1-\frac{1}{n}$. %\footnote{Make it more precise, for example $1-1/n$}.
To compute $\mathbb{P}[\mathbf{E}]$, it remains to estimate {the probability that~\ref{B1} and~\ref{B2} hold}. % conditioned on $\mathbf{E}_0$. 


Next, we will restrict the size of each part so that they fall into a prescribed interval with a positive probability.
This ensures that every large part \( A_i \cap S \) of \(\mathcal{P}'\) has a matching of size at least \( |A_i \cap S| - \frac{|S|}{r} + r \). Define $k_i:=|A_i|-n$ for each $i\in [s]$.  
% We decompose the event $\mathbf{E}$ into $s+1$ sub-events $\mathbf{E_i}$\todo{put this to the where it is defined?}, each of which measures the size of $A_i\cap S$ for each $i\in [s]$ or $B\cap S$ for $i=s+1$. Note that $\mathbb{P}[\mathbf{E}]=\prod_{i=1}^{s+1} \mathbb{P}[\mathbf{E_i}]$.
We proceed by considering the following  three cases based on the value of $\max_{i\in [s]}\{k_i\}$.



%For each $A_i\in \mathcal{P}$ with $|A_i|\le n$, we have $\mathbb{P}[|A_i\cap S|\le\frac{|A_i|}{2}]\ge \frac{1}{2}-o(1)$. Together with $|S|\ge \frac{rn}{2}-n^{0.6}$, we obtain
%\[
%|A_i\cap S|\le\frac{|A_i|}{2}\le\frac{n}{2}\le,
%\]
%which implies that $A_i\cap S$ is a small part with respect to $\mathcal{P}'$.
%Let $k_i := |A_i|-n$ for each $i \in [s]$, and assume without loss of generality that $k_1 \leq \cdots \leq k_s$.

% We proceed by considering the following  two cases based on the value of $k_s$.

% \medskip
% {\bf Case 1. $k_s>0$.}
% \medskip
%Note that $|A_i\cap S|\le\frac{|S|}{r}$ for each $A_i$ with $|A_i|\le n$ with probability $\frac{1}{2}-o(1)$. For each $A_i$ with $k_i>0$, Lemma~\ref{lem:inherit matching} gives that there is a matching of size at least $\frac{k_i}{32\beta}$ in $A_i\cap S$.
%By Lemma~\ref{lem:binomial distribution}, we have
%\[
%\mathbb{P}[|A_i\cap S|-\frac{|S|}{r}\le\frac{k_i}{32\beta}]=\mathbb{P}[B(n+|A_i|,1/2)\le n+\frac{k_i}{32\beta}]\ge \frac{1}{2}-o(1).
%\]

%Together with Lemma \ref{lem:covering bad vertices for balance}, we know that $G[S]$ has a $K_r$-factor with probability at least $o(1)$, as desired. 
%Let $E$ be the event that $\mathcal{P}'$ is a $(3\alpha,\alpha^{1/5},\frac{\beta}{4},\frac{\gamma'}{2})$-good partition.

% \begin{proof}
% Given $A_i$ with $|A_i|=\frac{n}{r}+k_i$, we know that $\de(G[A_i])\ge k_i+1$ and $e(G[A_i])\ge \frac{1}{2}|A_i|(k_i+1)$.
% Let $X=\sum_{e\in E(G[A_i])}X_e$ where $X_e=1$ if $e\in E(G[A_i])$, otherwise $X_e=0$.
% Note that $X\sim B(e(G[A_i]),1/2)$ and $E[X]=\frac{1}{4}e(G[A_i])$.
% Let $S_i$ be an indicator variable which denotes whether $v_i$ is selected, that is, 
% \[
% S_i = 
% \begin{cases} 
% 1, & v_i \in S; \\
% 0, & v_i \notin S .
% \end{cases}
% \]
% Note that $\{S_i\}_{i=1}^{|A_i|}$ are independent and identically distributed random variables defined on a complete probability space.
% Let $\mathcal{F}_k=\sigma(S_1,S_2,\cdots,S_k)$ and $Z_k=E[X\mid S_1,\cdots,S_k]$.
% Since
% \[
% E[Z_k\mid \mathcal{F}_{k-1}]=E[E[X\mid \mathcal{F}_k]\mid \mathcal{F}_{k-1}]=E[X\mid \mathcal{F}_{k-1}]=Z_{k-1},
% \]
% we obtain that $\{Z_k\}_{k\ge 0}$ is a martingale.
% Note that $Z_0=E[X]$, $Z_{|A_i|}=X$ the martingale difference is $$|Z_k-Z_{k-1}|\le\Delta(G[A_i])\le \beta n.$$
% Applying with Azuma-Hoeffding inequality, we have
% \[
% P[|X-E[X]|\ge t]\le 2\exp\left(-\frac{2t^2}{|A_i|\cdot(\beta n)^2}\right)=o(1),
% \]
% which implies that there are at least $\frac{1}{4}e(G[A_i])-n$ edges in $A_i\cap S$ with probability $1-o(1)$.
% Therefore, we have
% \[
% \frac{1}{4}e(G[A_i])-\beta n^{2/3}-2ck_i\cdot\Delta(G[A_i])\ge k_i\cdot(\frac{1}{4}|A_i|-2c\beta n)+\frac{1}{4}|A_i|-\beta n^{2/3}
% \]


% \end{proof}


% Since the partition $\mathcal{P}'$ inherits the conditions on degree and independent sets with high probability, we can also find a $K_r$-tiling of size at most $10r\alpha n$ that cover all bad vertices, just as in subsection 4.2.
% Removing such $K_r$-tiling, we obtain a new partition $\mathcal{P}=\{A_1\cap S,\cdots,A_s\cap S,B\cap S\}$ (adhere to the previous notation for convenience) satisfying
% \begin{enumerate}
% [label =\rm  (E\arabic{enumi})]
%     %\item\label{E1} $||A_i\cap S|-\frac{|S|}{r}|\le |S|^{0.1}$ if $|A_i|\le n$ and $||A_i\cap S|-\frac{|S|}{r}|\le k_i+$ if $|A_i|>n$ where $k_i=|A_i|-\frac{|S|}{r}$;
%     \item\label{E1} $\Delta(G[A_i\cap S])\leq \beta |A_i\cap S|+n^{0.1}$ for each $i\in [s]$;
%     \item\label{E2} $d(v,A_i\cap S)\geq |A_i\cap S|-\beta n-n^{0.1}$ for each vertex $v\in S\setminus A_i$;
%     \item\label{E3} $d(v,B\cap S)\geq |B\cap S|-\beta n-n^{0.1}$ for each vertex $v\in S\setminus B$;
%     \item\label{E4} $(r-s)\mid |B|$ and $B\cap S$ has no $\frac{\gamma_{s+1}}{2}$-independent set of size $n/2$.
% \end{enumerate}

% %Without loss of generality, assume that $|A_i\cap S|=\frac{|S|}{r}+k_i$ for each $i\in [s]$ and $k_1\leq \cdots\leq k_s$ where $\sum_{i\in [s]}k_i=0$.
% To complete the proof via Lemma~\ref{lem:balance} and Claim~\ref{lem:inherit matching}, it remains to show that there exists a large matching in each large part.


\medskip
{\bf Case 1. $\max_{i\in [s]}\{k_i\}>\delta \sqrt{n}$.}% and $r-s\neq 1$.}
\medskip


Without loss of generality, assume that $k_1:=\max_{i\in [s]}\{k_i\}$. By Lemma~\ref{lem-int}, %and Lemma~\ref{lem:binomial distribution}, 
we have 
\[
\mathbb{P}\Big[\mathbf{E_1}:\frac{|A_1|}{2}+{\sqrt{|A_1|}}\leq |A_1\cap S|\leq \frac{|A_1|}{2}+5{\sqrt{|A_1|}}\Big]
\footnote{which differs by at most $\delta$ from $$
\int_{2}^{10}\frac{1}{\sqrt{2\pi}}e^{-t^2/2}\,dt=\Phi(10)-\Phi(2)\geq \frac{1}{50},
$$ 
where we take Taylor expansion {$\Phi(z)=\frac{1}{2}+\frac{1}{\sqrt{2\pi}}\sum_{k=0}^{\infty}\frac{(-1)^kz^{2k+1}}{k!\cdot 2^k\cdot(2k+1)}$}.}\ge \frac{1}{100}.\]
%Thus, $\mathbb{P}(E)\geq \frac{1}{2^r}-o(1)$. 
Similarly, for each $i\in [2,s]$ we 
have 
$$
\mathbb{P}\Big[\mathbf{E}_i:\frac{|A_i|}{2}-\frac{1}{2r}{\sqrt{|A_i|}}\leq |A_i\cap S|\leq \frac{|A_i|}{2}+\frac{1}{2r}{\sqrt{|A_i|}}\Big]\geq \frac{1}{2r},
$$ 
and
\[
\mathbb{P}\Big[\mathbf{E}_{s+1}:\frac{|B|}{2}-\frac{1}{2r}{\sqrt{|B|}}\leq |B\cap S|\leq \frac{|B|}{2}+\frac{1}{2r}{\sqrt{|B|}}\Big]\geq \frac{1}{2r}.
\]
Therefore, conditioned on all events $\mathbf{E}_i$ for $i\in [s+1]$, % when all these events occur simultaneously, 
we apply the Cauchy-Schwarz inequality to obtain   
\begin{align*}
|S|&=\sum_{i\in [s]}|A_i\cap S|+|B\cap S|\geq \frac{rn}{2}+{\sqrt{|A_1|}-\frac{1}{2r}\sum_{i\in [2,s]}\sqrt{|A_i|}-\frac{1}{2r}\sqrt{|B|}}\\
&=\frac{rn}{2}+{\frac{1}{2r}\Big((2r+1)\sqrt{|A_1|}-\sum_{i\in [s]}\sqrt{|A_i|}-\sqrt{|B|}\Big)}\\
&\geq 
\frac{rn}{2}+{\frac{1}{2r}\Big((2r+1)\sqrt{|A_1|}-r\sqrt{n}\Big)}\\
&\geq\frac{rn}{2}+\frac{r+1}{2r}\sqrt{|A_1|},  
\end{align*}
and 
$$
|S|=\sum_{i\in [s]}|A_i\cap S|+|B\cap S|\leq \frac{rn}{2}+5\sqrt{|A_1|}+\frac{r-1}{2r}\sqrt{|A_1|}\le\frac{rn}{2}+6\sqrt{|A_1|}.
$$
Hence, 
\begin{align}\label{A1capS}
|A_1\cap S|-\frac{|S|}{r}\le\frac{|A_1|}{2}+5\sqrt{|A_1|}-\frac{n}{2}-\frac{r+1}{2r^2}\sqrt{|A_1|}\le\frac{|A_1|-n}{2}+5\sqrt{|A_1|}\leq 3\alpha \frac{|S|}{r};
\end{align}
for each $i\in [2,s]$, we have 
\begin{align}\label{AcapS}
|A_i\cap S|-\frac{|S|}{r}\le\frac{|A_i|}{2}+\frac{1}{2r}\sqrt{|A_i|}-\frac{n}{2}-\frac{r+1}{2r^2}\sqrt{|A_1|}\leq \frac{|A_i|-n}{2}-\frac{1}{2r^2}\sqrt{|A_1|}\leq 3\alpha \frac{|S|}{r}.
\end{align}
%and 
%\begin{align*}
%{\sout{|A_1\cap S|-\frac{|S|}{r}\geq |A_i\cap S|-\frac{|S|}{r}\ge\frac{|A_i|}{2}-\frac{1}{2r}\sqrt{n}-\frac{n}{2}-\frac{6}{r}\sqrt{n}=\frac{|A_i|-n}{2}-\frac{13}{2r}\sqrt{n}\geq -3\alpha |S|.}}
%\end{align*}
%\begin{align}\label{AcapS}
%|A_i\cap S|-\frac{|S|}{r}\le\frac{|A_i|}{2}+\frac{1}{2r}\sqrt{n}-\frac{n}{2}-\frac{r+1}{2r^2}\sqrt{n}=\frac{|A_i|-n}{2}-\frac{1}{2r^2}\sqrt{n}\leq 3\alpha |S|,
%\end{align}
%and 
%\[
%|A_i\cap S|-\frac{|S|}{r}\ge\frac{|A_i|}{2}-\frac{1}{2r}\sqrt{n}-\frac{n}{2}-\frac{6}{r}\sqrt{n}=\frac{|A_i|-n}{2}-\frac{13}{2r}\sqrt{n}\geq -3\alpha |S|.
%\]
Similarly, 
\[
-3\alpha \frac{|S|}{r}\leq |B\cap S|-{(r-s)\frac{|S|}{r}}<3\alpha \frac{|S|}{r}.
\]
That is,~\ref{B1} holds. Fix an $i\in [s]$ with $|A_i\cap S|\geq \frac{|S|}{r}$. If \(i=1\), then the assumption \(\max_{i\in[s]}\{k_i\} > \delta\sqrt{n}\) yields \(|A_1|\ge n+\delta\sqrt{n}\);  if \(i\in[2,s]\), then by~\eqref{AcapS} we have \(|A_i|\ge n+\frac{1}{r^2}\sqrt{n}\). %$A_i\cap S$ is a large part of $\mathcal{P}'$, then by~\eqref{AcapS} one has $|A_i|\geq n+\frac{1}{r^2}\sqrt{n}$. 
Recall that $\Delta(G[A_i])\leq \beta n$. Then, w.h.p.~we have $\Delta(G[A_i\cap S])\leq \Delta(G[A_i])\leq \beta n\leq 3\beta \frac{|S|}{r}$.  By Remark~\ref{remark:vertex-cover} and {the condition that} $G$ is $((r-1)n+1)$-regular, we conclude that the matching number of $G[A_i]$ is at least $$\frac{e(G[A_i])}{2\beta n}\geq \frac{\delta(G[A_i])|A_i|}{4\beta n}\geq \frac{(|A_i|-n+1)|A_i|}{4\beta n}\geq \frac{|A_i|-n+1}{5\beta}.$$ 
%Moreover, \ref{B1} implies that $|B\cap S|\leq \frac{(r-s)|S|}{r}$ if $r-s\geq 2$. 
%Otherwise, $|A_i\cap S|\le \frac{|S|}{r}$.
%For the former case, b
By Chernoff bound, {\eqref{A1capS}-\eqref{AcapS} and $\beta\ll \delta$}, w.h.p.~we have that $G[A_i\cap S]$ contains a matching of size at least $\frac{|A_i|-n}{21\beta}>|A_i\cap S|-\frac{|S|}{r}+r$, thus~\ref{B2} holds.
It follows that  $${\mathbb{P}[\mathbf{E}]\geq \mathbb{P}\big[\mathbf{E}_0\mathbf{E}_1\mathbf{E}_2\cdots\mathbf{E}_{s+1} \big]\geq\mathbb{P}\Big[{\bigcap_{i\in [s+1]}\mathbf{E}_i }\Big]+%\footnote{are they independent events under $E_0$?}
\mathbb{P}[\mathbf{E}_0 ]-1\ge \frac{1}{100(2r)^{r-1}}-\frac{1}{n}.}
$$ 
%{more computational details}.  {what about $B04-B06$?} 
% and $G[B\cap S]$ contains a matching of size at least $|B\cap S|-\frac{|S|}{r}$ if $r-s=1$.
%Notice that $\mathbb{P}(r\mid |S|)=\frac{1}{r}$. Together with Lemma~\ref{lem:covering bad vertices for balance}, we obtain that $G[S]$ has a $K_r$-factor with probability at least $\frac{1}{100r(2r)^{r-1}}$. 


% {\color{red}Note that $k_1\le k_2\le\cdots \le k_s$ and $k_s\ge \de n^{0.2}$. Together with $|S|=\frac{rn}{2}\pm n^{0.1}$. Then $\frac{|S|}{r}\ge \frac{n}{2}+\de n^{0.1}$. Thus, Claim~\ref{claim:B1-B06 in S}\ref{D1} gives that if $A_i$ is a small part or $k_i\le\de n^{0.2}$ with respect to $\mathcal{P}$, then w.h.p.~$A_i\cap S$ is also a small part with respect to $\mathcal{P}'$. If $k_i>\de n^{0.2}$, Lemma~\ref{lem:inherit matching} implies that there is a matching of size at least $\frac{k_i}{32\beta}$ in $A_i\cap S$ with probability $1-o(1)$. Therefore, we verify that condition~\ref{B2} in Lemma~\ref{good-partition} holds. Thus, we have
% \[
% P[G[\frac{1}{2}]~\text{contains a $K_r$-factor}]\ge P[r\mid |S|]-o(1)\ge \frac{1}{r}-o(1).
% \]}

\medskip
{\bf Case 2.} 
$\max_{i\in [s]}\{k_i\}\leq \de \sqrt{n}$ and $s=r$.
\medskip

%For convenience, we denote $B$ as $A_r$ when $r-s=1$.
We consider any balanced partition $\{A_1^*,A_2^*,\ldots,A_r^*\}$ obtained from $\{A_1,A_2,\ldots,A_r\}$ by moving less than $\de \sqrt{n}$ vertices from each large part to small parts. Let $C_i$ be a minimum vertex cover of $G[A_i^*]$ for each $i\in [r]$. By Lemma~\ref{vertex-cover}, we have $\max\{|C_1|,\ldots,|C_r|\}\geq \frac{1}{r}\sqrt{n}$. Without loss of generality, assume that $|C_1|\geq \frac{1}{r}\sqrt{n}$. 
As in Remark~\ref{remark:vertex-cover}, $G[A_1^*]$ has a matching $M$ of size at least $\frac{\sqrt{n}}{2r}$. By Chernoff bound, w.h.p.~$M[S]$ contains at least $\frac{2\sqrt{n}}{17r}$ edges. Observe that at most $\de \sqrt{n}$ edges in $M[S]$ contain a {vertex that is moved from a large part to $A_1$}, so this gives a matching of size at least $\frac{\sqrt{n}}{9r}$ in ${G'}[A_1\cap S]$. By Lemma~\ref{lem-int}, % and Lemma~\ref{lem:binomial distribution},
we have 
\begin{align*}\label{eq:the size of A1}
&\ \mathbb{P}\Big[\mathbf{E}_1:\frac{|A_1|}{2}+ \frac{1}{20r}\sqrt{{|A_1|}}\leq |A_1\cap S|\leq \frac{|A_1|}{2}+\frac{1}{10r}\sqrt{{|A_1|}}\Big]\\
=&\ \mathbb{P}\Big[\frac{|A_1|}{2}+\frac{1}{20r}\sqrt{{|A_1|}}\leq B(|A_1|,1/2)\leq \frac{|A_1|}{2}+\frac{1}{10r}\sqrt{{|A_1|}}\Big]\\
%&\ge \int_{\frac{1}{10r}}^{\frac{1}{5r}}\frac{1}{\sqrt{2\pi}}e^{-t^2/2}\,dt-\de\\
\geq&\  \Phi\Big(\frac{1}{5r}\Big)-\Phi\Big(\frac{1}{10r}\Big)-\de\geq \frac{1}{30r}.
\end{align*}
Similarly, for each $i\in [2,r]$ we 
have 
\begin{align*}
\mathbb{P}\Big[\mathbf{E}_i:\frac{|A_i|}{2}-\frac{1}{40r^2}\sqrt{{|A_i|}}\leq |A_i\cap S|\leq \frac{|A_i|}{2}+\frac{1}{40r^2}\sqrt{{|A_i|}}\Big]\geq \frac{1}{40r^2}.
\end{align*}
Therefore, conditioned on all events $\mathbf{E}_i$ for $i\in [r]$, by the Cauchy-Schwarz inequality we have  
\begin{align*}
|S|&=\sum_{i\in [r]}|A_i\cap S|\ge \frac{rn}{2}+\frac{1}{20r}\sqrt{{|A_1|}}-\frac{1}{40r^2}\sum_{i\in [2,r]}\sqrt{|A_i|}\\
&= \frac{rn}{2}+\frac{1}{40r^2}\Big((2r+1)\sqrt{|A_1|}-\sum_{i\in [r]}\sqrt{|A_i|}\Big)\\
&\ge \frac{rn}{2}+\frac{1}{40r^2}\left((2r+1)\sqrt{|A_1|}-r\sqrt{n}\right).
\end{align*}
% \leq \frac{rn}{2}+\frac{5r-1}{40r^2}\sqrt{n}
Note that $\max_{i\in [s]}\{k_i\}\leq \de \sqrt{n}$, {which implies that $|A_i|=n\pm r\de\sqrt{n}$ for each $i\in [r]$. As $\de\ll\frac{1}{r}$,} for each $i\in [2,r]$, we have 
\begin{align*}
|A_i\cap S|-\frac{|S|}{r}&\le\frac{|A_i|}{2}+\frac{1}{40r^2}\sqrt{|A_i|}-\frac{n}{2}-\frac{1}{40r^3}\left((2r+1)\sqrt{|A_1|}-r\sqrt{n}\right)\\
&\leq \frac{\delta\sqrt{n}}{2}-\frac{1}{40r^3}\left(2r\sqrt{|A_1|}-r\sqrt{n}\right)<0.
\end{align*}
%and 
%\[{\sout{|A_i\cap S|-\frac{|S|}{r}\ge\frac{|A_i|}{2}-\frac{1}{40r^2}\sqrt{n}-\frac{n}{2}-\frac{5r-1}{40r^3}\sqrt{n}=\frac{|A_i|-n}{2}-\frac{6r-1}{40r^3}\sqrt{n}\geq -3\alpha |S|.}}
%\]
Similarly, 
$$
%-3\alpha |S|\le{\frac{|A_1|-n}{2}+\frac{2r-5}{40r^2}\sqrt{n}}\leq 
|A_1\cap S|-\frac{|S|}{r}<\frac{\sqrt{|A_1|}}{10r}<\frac{2\sqrt{n}}{19r}<3\alpha \frac{|S|}{r}. 
$$
This implies that~\ref{B1} holds and $A_i\cap S$ is a small part of $\mathcal{P}'$ for each $i\in [2,r]$. 
Recall that $G[A_1\cap S]$ contains a matching of size at least $\frac{\sqrt{n}}{9r}$, 
which is larger than $|A_1\cap S|-\frac{|S|}{r}+r$. Thus~\ref{B2} holds. It follows that  $${\mathbb{P}[\mathbf{E}]\geq \mathbb{P}\big[\mathbf{E}_0\mathbf{E}_1\mathbf{E}_2\cdots\mathbf{E}_r \big]\geq \mathbb{P}\Big[{\bigcap_{i\in [r]}\mathbf{E}_i}\Big]+
\mathbb{P}[\mathbf{E}_0 ]-1\ge  \frac{1}{30r(40r^2)^{r-1}}-\frac{1}{n}.}$$
 %Together with Lemma~\ref{lem:covering bad vertices for balance}, we know that $G[S]$ has a $K_r$-factor with probability at least $\frac{1}{(40r^2)^r}$, as desired.

\medskip
{\bf Case 3.} $\max_{i\in [s]}\{k_i\}\leq \de\sqrt{n}$ and $s\leq r-1$.
\medskip

%In this case, $B$ is a set of size $(r-s)n$ containing no $\gamma'$-independent $n$-set. Choose $\eps'\ll\eps\ll \gamma'$.  If $r-s=1$, then there are  $\eps n$ vertex disjoint edges in $G[B]$ as $e(G[B])\ge \gamma'n^2$ and $\gamma'\gg \eps$. If $r-s\ge 2$, then \ref{B05} implies  
%\[
%\de(G[B])\ge (r-1)n+1-sn-r\alpha n\ge \left(1-\frac{1}{r-s}-\beta\right)\cdot|B|.
%\]
%By Theorem \ref{SuperT}, there are at least $\eps n$ vertex-disjoint copies of $K_{r-s+1}$ in $G[B]$. By Chernoff bound, we know that w.h.p.~$G[B\cap S]$ has at least $\frac{\eps}{2^{r+1}}n$ vertex disjoint  copies of $K_{r-s+1}$. 
By a similar discussion as in Case~2, for each $i\in [s]$ we have 
\[
\mathbb{P}\Big[\mathbf{E}_i:\frac{|A_i|}{2}-\frac{1}{40r^2}\sqrt{|A_i|}\leq |A_i\cap S|\leq \frac{|A_i|}{2}+\frac{1}{40r^2}\sqrt{|A_i|}\Big]\geq \frac{1}{40r^2},
\]
and
$$
\mathbb{P}\Big[\mathbf{E}_{s+1}:\frac{|B|}{2}+\frac{1}{20r}\sqrt{|B|}\leq |B\cap S|\leq \frac{|B|}{2}+\frac{1}{10r}\sqrt{|B|}\Big]\ge \frac{1}{30r}.
$$
Thus, conditioned on all events $\mathbf{E}_i$ for $i\in [s+1]$, by a direct calculation we have 
$$
%-3\alpha |S|\leq 
|A_i\cap S|-\frac{|S|}{r}<0\ \ \text{for each}\ i\in [s]$$
and
$$-3\alpha \frac{|S|}{r}\leq |B\cap S|-{(r-s)\frac{|S|}{r}}<3\alpha \frac{|S|}{r}. 
$$
It follows that~\ref{B1} holds and~\ref{B2} becomes trivial, while 
$${\mathbb{P}[\mathbf{E}]\geq \mathbb{P}\big[\mathbf{E}_0\mathbf{E}_1\mathbf{E}_2\cdots\mathbf{E}_{s+1} \big]\geq \mathbb{P}\Big[{\bigcap_{i\in [s+1]}\mathbf{E}_i }\Big]+
\mathbb{P}[\mathbf{E}_0 ]-1\ge  \frac{1}{30r(40r^2)^{r-1}}-\frac{1}{n}.}$$
%$$
%\frac{|S|}{r}\geq \frac{\sum_{i\in[r]}|A_i\cap S|+|B\cap S|}{r}\geq \frac{n}{2}-\frac{r-1}{40r^3}\sqrt{n}+\frac{1}{20r^2}\sqrt{ n}>|A_i\cap S|\ \text{for each}\ i\in [s].
%$$
%Moreover,
%$$
%|B\cap S|-\frac{|S|}{r}\leq \frac{|B|}{2}+\frac{\sqrt{n}}{10r}-\Big(\frac{n}{2}+\frac{1}{40r^2}\sqrt{n}\Big)\leq \frac{\sqrt{n}}{10r}.
%$$

%Together with Lemma \ref{lem:covering bad vertices for balance}, we obtain that $G[S]$ has a $K_r$-factor with probability at least $\frac{1}{(40r^2)^r}$.}}

\medskip
Combining Cases~1-3, we derive $\mathbb{P}[\mathbf{E}]\geq \frac{1}{35r(40r^2)^{r-1}}$.  Notice that $\mathbb{P}[|S|\equiv0~(\text{mod}~r)]=\frac{1}{r}$. Together with Lemma~\ref{lem:covering bad vertices for balance}, we obtain that $G[S]$ has a $K_r$-factor with probability at least $\frac{1}{(40r^2)^r}$. 
\end{proof}


\section{Proof of Lemmas}\label{sec:lemma}
In this section, we give the proofs of Lemma~\ref{lem:construct good-partition}, Lemma~\ref{lem:covering bad vertices for balance} and Lemma~\ref{vertex-cover}, respectively. 
\subsection{Proof of Lemma~\ref{lem:construct good-partition}}




\begin{proof}[{\bf Proof of Lemma~\ref{lem:construct good-partition}}]
Choose constants $r\geq 3,n\in \mathbb{N}$ and $\gamma_1,\gamma_2,\dots,\gamma_r$ such that 
\begin{align*}%\label{eq:p2}
    \frac{1}{n}\ll \gamma \ll \gamma_1\ll \dots \ll \gamma_r\ll\frac{1}{r}.
\end{align*}
Suppose that $G$ is an $rn$-vertex graph with $\delta(G)\geq (r-1)n+1$ and $G$ contains a $\gamma$-independent $n$-set. We proceed to greedily choose a maximal family of vertex-disjoint $n$-sets $A_1^0, \dotsc, A_s^0$ such that each $A_i^0$ is $\gamma_i$-independent.
Let $B^0 := V(G) \setminus \bigcup_{i \in [s]} A_i^0$ and denote $\mathcal{P}_0:=\{A_1^0,\ldots,A_s^0,B^0\}$.
Note that $G[B^0]$ has no $\gamma_{s+1}$-independent $n$-set and $|B^0|=(r-s)n$.

We additionally choose 
\begin{align*}%\label{eq:p3}
    \gamma_s\ll \alpha \ll\beta\ll \gamma_{s+1}.
\end{align*}
It is easy to see that for distinct $i,j\in [s]$, any vertex $v\in V(G)\setminus (A_i^0\cup A_j^0)$ cannot be both $(2\beta,A_i^0)$-bad and $(2\beta,A_j^0)$-bad. Next, we estimate the number of bad or exceptional vertices with respect to the partition $\mathcal{P}_0$. % and there is no~\ref{bad vtx3} and~\ref{bad vtx4} if $s=r$ in the following claim.
\begin{claim}\label{fact:bad vertices for Ai}
For each $i\in [s]$, the following hold:
\begin{enumerate}[label=(\arabic*)]
    \item\label{large-degree vtx} the number of $(\alpha^{1/4},A_i^0)$-good vertices in $A_i^0$ is at least $|A_i^0|-\alpha n$;
    \item\label{bad vtx2} the number of {$(2\beta,A_i^0)$}-bad vertices in $V(G)\setminus A_i^0$ is at most ${\alpha}n$;
    \item\label{bad vtx3} the number of {$(2\beta,B^0)$}-bad vertices in $V(G)\setminus B^0$ is at most $(r-1){\alpha} n$ if $s=r-1$, and there is no $({2\beta},B^0)$-bad vertex in $V(G)\setminus B^0$ if $s<r-1$;
    %\item\label{bad vtx3} any vertex $v\in V(G)$ cannot be both $(\beta,A_i)$-bad and $(\beta,A_j)$-bad for distinct $i,j\in [s]$;
    \item\label{bad vtx4} the number of $(\alpha^{1/4},B^0)$-exceptional vertices in $B^0$ is at most $\alpha n$.
\end{enumerate}
\end{claim}

\begin{proof}
(1) Recall that $A_i^0$ is a $\gamma_i$-independent set of size $n$ for each $i\in [s]$. As $\gamma_i\ll \alpha$, we obtain that the number of vertices in $A_i^0$ that are not $(\alpha^{1/4},A_i^0)$-good is at most 
$
\frac{2\gamma_i n^2}{\alpha^{1/4}n}\leq \alpha n$,
as desired.\medskip

(2) Since {$\delta(G)\geq (r-1)n+1$} and $A_i^0$ is a $\gamma_i$-independent set, the number of edges in $G[A_i^0,V(G)\setminus A_i^0]$ is at least
$$\underset{v\in A_i^0}{\sum}d(v)-2e(G[A_i^0])\geq 
|A_i^0| \cdot ((r-1)n+1)-2 e(G[A_i^0]) \ge |A_i^0||V(G)\setminus A_i^0|-2\gamma_i n^2.
$$
This implies that there are at most $2\gamma_i n^2$ non-edges between $A_i^0$ and $V(G)\setminus A_i^0$.
Note that every $(2\beta,A_i^0)$-bad vertex in $V(G)\setminus A_i^0$ contributes at least $(1-2\beta) n$ non-edges between $A_i^0$ and $V(G)\setminus A_i^0$.
As $\gamma_i \ll \alpha \ll \beta$,  we have that the number of $(2\beta,A_i^0)$-bad vertices in $V(G)\setminus A_i^0$ is at most 
$\frac{2\gamma_i n^2}{(1-2\beta) n}\le{\alpha}n$. \medskip


(3) If $s< r-1$, then $d(v,B^0)\ge |B^0|-n+1\ge n+1$ for every $v\in V(G)$. Hence, $G$ contains no $(2\beta,B^0)$-bad vertex. If $s=r-1$ and $v\in A_i^0$ is $(2\beta,B^0)$-bad, then $v$ cannot be an $(\alpha^{1/4},A_i^0)$-good vertex. Together with (1), we conclude that the number of $(2\beta,B^0)$-bad vertices is at most $(r-1){\alpha} n.$\medskip
% $$|V_{ex}(\mathcal{P}_0,\beta,i)|\leq |V_{ne}(\mathcal{P}_0,\alpha^{1/5},i)|\le \frac{2\gamma_i n^2}{\alpha^{1/5}n}\leq  \alpha n.$$

%(3) Let $u\in A_i^0$ be a $(\beta,A_j)$-bad vertex for distinct $i,j\in [s]$. Then $d(u,A_{\ell}^0)\geq n-\beta n$ for all $\ell\in [s]\setminus \{j\}$ due to the regularity of $G$. 

(4) For each $i\in [s]$, let $B_i\<B^0$ be the set of all $(\alpha^{1/4},B^0)$-exceptional vertices $v$ such that $d(v,A_i^0)\le |A_i^0|-\alpha^{1/4}n$ and $B_e=\bigcup_{i\in [s]}B_i$. Now we consider the size of $B_i$.
Note that $|B^0|=(r-s)n$.
Thus 
\[
e(G[A_i^0,B^0])\le |B_i|\cdot \big(|A_i^0|-{\alpha^{1/4} n}\big)+\big((r-s)n-|B_i|\big)\cdot|A_i^0|=|A_i^0|(r-s)n-|B_i|\alpha^{1/4}n.
\]
Recall that $\delta(G)\geq (r-1)n+1$ and $A_i^0$ is a $\gamma_i$-independent set. Then % Combing with $d(v)=(r-1)n+1$ for each $v\in A_i$, we obtain that
\[
2\gamma_in^2\geq 2e(G[A_i^0])\ge|A_i^0|\cdot\big((r-1)n+1-(s-1)n\big)-e(G[A_i^0,B^0])\ge |A_i^0|+|B_i|\alpha^{1/4}n=n+|B_i|\alpha^{1/4}n,
\]
which implies $|B_i|\leq \frac{\alpha}{r} n$ as $\gamma_i\ll \alpha\ll \frac{1}{r}$. Thus 
%Recall that $A_i^0$ is a $\gamma_i$-independent set of $G$ and at most $r\alpha n$ vertices are moved from other part to $A_i^0$.
%As $\gamma_i\ll\alpha\ll\frac{1}{r}$, we have
%$$
%e(G[A_i])\leq (\gamma_i+2r\alpha) n^2<\sqrt{\alpha}n^2.
%$$
%Thus $|A_i|+|B_i|\cdot\beta n\le e(G[A_i])\le \sqrt{\alpha}n^2$, which implies that
$
|B_e|\leq \sum_{i\in [s]}|B_i|\le \alpha n.$ 
%as $\alpha\ll\beta$.
\end{proof}

Now, we move all $(2\beta, D^0)$-bad vertices from $V(G)\setminus D^0$ to $D^0$ for each $D^0\in \mathcal{P}_0$, and denote the resulting partition by $\mathcal{P}_1:=\{A_1^1,\ldots,A_s^1,B^1\}$. {As $\alpha\ll\beta$, Claim~\ref{fact:bad vertices for Ai} (2) implies} that every vertex in $D^1\in \mathcal{P}_1$ has at least {$\frac{3}{2}\beta n$} neighbors in each $D'$ with $D'\in \mathcal{P}_1\setminus \{D^1\}$. By Claim~\ref{fact:bad vertices for Ai} (2)-(3), one has $|A_i^1|\leq n+\alpha n$ for each $i\in [s]$ and $|B^1|\leq (r-s)n+(r-1)\alpha n$. %Based on Claim \ref{fact:bad vertices for Ai}, we obtain that \ref{B1}, \ref{B04} and \ref{B05} hold. 
%Next, we will proceed with the following two-step operations of moving at most $\alpha n$ vertices. The first step aims to obtain a new partition satisfying the degree condition in Lemma~\ref{lem:balance}. It is obvious that the new partition may be unbalanced. So the purpose of Step 2 is to identify a collection of matchings in the \textit{large} parts, whereby we adjust the partition to be balanced.

%\medskip
%{\bf Step 1. Move all $(\beta, D)$-bad vertices from $V(G)\setminus D$ to $D$ for each $D\in \mathcal{P}_0$.}
%\medskip

%For each $D\in \mathcal{P}$, if $v\in V(G)\setminus D$ is $(\beta,D)$-bad, then we move $v$ to $D$. 
%In the process of moving vertices, there is a tiny difference between the cases $s< r-1$ and the others in whether to move vertices to $B^0$. To be precise, i
%Notice that if $s< r-1$, then $d(v,B^0)\ge |B^0|-n+1\ge n+1$ for every $v\in V(G)$. In this case, we do not need to move any vertices to $B^0$. If $s\ge r-1$, then there may exists $v\in A_i$ such that $d(v,B^0)\le \beta n$ for some $i\in [s]$ (such a vertex cannot be $(\beta,A_i)$-good). In this case, we move such vertex $v$ into $B^0$. Note that the resulting partition after the above moving procedure may not be balanced and we denote it as $\mathcal{P}_1:=\{A_1^1,\ldots,A_s^1,B^1\}$, in which every vertex in $A_i^1$ has large degree to the outside of $A_i^1$. 
Recall that for each $i \in [s]$, a part $A_i^1$ is called {large} if $|A_i^1| > n$, and {small} otherwise. 
Next, we will move vertices in large parts to satisfy~\ref{B2}.  
%\medskip
%{\bf Step 2. Move not $(\beta, A_i)$-bad vertices from large part $A_i^1$ to small one $A_j^1$.}
%\medskip
Let $A_i^1$ be a large part. % set with size greater than $n$. 
If $\Delta(G[A_i^1])\leq \beta n$, then we are done. Otherwise, move a vertex in \(A_i^1\) which has at least \( \beta n \) neighbors in \(A_i^1\) to any part \(A_j^1\) with $|A_j^1|<n$, or $B^1$ if $|B^1|<(r-s)n$, and update \(\mathcal{P}^1\) accordingly. %Assume that $|A_i^1|>n$ and $\Delta(G[A_i^1])\leq \beta n$ for some $i\in [s]$.  arbitrarily move each vertex in $A_i^1$ which has at least $\beta n$ neighbors in $A_i^1$ to other side $A_j^1$, where $j\in [s]$ with $|A_j^1|<n$, and update $\mathcal{P}^1$ accordingly. 
Repeat this until we reach the point {where} either $|A_j^1|=n$ for each $j\in [s]$ and $|B^1|=(r-s)n$, or $G[A_i^1]$ has maximum degree at most $\beta n$ whenever $|A_i^1|>n$ for $i\in [s]$; denote the resulting partition by $\mathcal{P}:=\{A_1,\ldots,A_s,B\}$. 
Obviously,~\ref{B1} holds. Moreover,  each vertex has  at least $(r-s-1)n-r\alpha n$ neighbors in $B$ and every vertex in $D$ has at least $\beta n$ neighbors in each $D'$ with $D'\in \mathcal{P}\setminus \{D\}$, that is,~\ref{B04}-\ref{B05} hold.

Recall that $\de(G)\geq (r-1)n+1$. Then $\de(G[A_i])\geq |A_i|-n+1$ when $|A_i|\geq n$ with $i\in [s]$.  %Let $|A_i|=n+k_i$ for each $i\in [s]$. 
If $|A_i|>n$, then together with $\Delta(G[A_i])\leq \beta n$, Remark~\ref{remark:vertex-cover} implies that $G[A_i]$ has a matching of size at least
$$
\frac{\de(G[A_i])|A_i|}{4\Delta(G[A_i])}\geq \frac{(|A_i|-n+1)|A_i|}{4\beta n}\geq \frac{|A_i|-n+1}{5\beta}\ge |A_i|-n+r.
$$
That is,~\ref{B2} holds. 

{By Claim~\ref{fact:bad vertices for Ai} (1)-(2) we obtain~\ref{B03}, and by (3)-(4) we obtain~\ref{B06}. Recall that %each $A_i^0$ is a $\gamma_i$-independent set, 
$B^0$ contains no $\gamma_{s+1}$-independent $n$-set and $\alpha\ll\gamma_{s+1}$. %Then %for each $i\in [s]$ we have 
%$$
%e(G[A_i])\leq e(G[A_i^0])+|A_i\setminus A_i^0||A_i|\leq \gamma_in^2+2\alpha n\cdot (n+\alpha n)\leq 3\alpha n^2. 
%$$
%Moreover, 
Let $B'$ be a subset of $B$ with size $\lceil\frac{|B|}{r-s}\rceil$. If $|B'\cap B^0|\geq n$, then $e(G[B'])\geq e(G[B'\cap B^0])\geq \gamma_{s+1}n^2$.  If $|B'\cap B^0|< n$, then there is a subset $B''$ such that  $B'\cap B^0\subseteq B''\subseteq B^0$ with $|B''|=n$. Therefore,  
$$
e(G[B'])\geq e(G[B''])-|B''\setminus B'||B''|\geq \gamma_{s+1}n^2-r\alpha n\cdot n\geq \frac{\gamma_{s+1}}{2}n^2.
$$
That is, $G[B]$ contains no  $\frac{\gamma_{s+1}}{2}$-independent set of size at least $\frac{|B|}{r-s}$.  Hence,~\ref{B07} holds by taking $\gamma':=\frac{\gamma_{s+1}}{2}$. Thus, $\mathcal{P}$ is an  $(\alpha,\beta,\beta,\gamma')$-good partition of $G$, as desired.} 
%Since Together with Claim \ref{fact:bad vertices for Ai}, it is straightforward to check that~\ref{B1}-\ref{B06} hold and~\ref{B07} holds with parameters $2\gamma_s$ and $\frac{\gamma_{s+1}}{2}$.  %(you chose $\gamma_{s+1}\ll \gamma'$ at the beginning!)}. 
%\medskip
%After the two-step moving operation, we state the following fact that a large matching can be found in a large part.
\end{proof}
%Let $r,n,s$ be positive integers with $r\ge 3,\ 1\le s\le r$ and let $\gamma$ be a positive constant. Define constants $\gamma, \gamma_1,\gamma_2,\dots,\gamma_r, \gamma_{r+1}$ satisfying 
%\begin{align}\label{eq:p2}
%    0<\frac{1}{n}\ll \gamma \ll \gamma_1\ll \dots \ll \gamma_r\ll\gamma_{r+1}=\frac{1}{r}.
%\end{align}

%Since $G$ contains a $\gamma$-independent $n$-set, we may greedily choose vertex-disjoint $\gamma_i$-independent $n$-sets $A_1^0, \dotsc, A_s^0$, until no $\gamma_{s+1}$-independent $n$-set remains. 
%Let $B^0 := V(G) \setminus \bigcup_{i \in [s]} A_i^0$ and we obtain a partition $\mathcal{P}_0:=\{A_1^0,\ldots,A_s^0,B^0\}$ of $V(G)$.
%Note that $B^0$ has no $\gamma_{s+1}$-independent $n$-set and $|B^0|=(r-s)n$.

%When the constant $s$ is fixed, we define 
%\begin{align}\label{eq:p3}
%    \gamma_s\ll \alpha \ll\beta\ll \gamma_{s+1}.
%\end{align}





%In the rest of the proof, our main challenge is to adjust the partition $\mathcal{P}$ slightly to satisfy the following degree condition.



\subsection{Proof of Lemma~\ref{lem:covering bad vertices for balance}}

Keevash and Mycroft \cite{KM-multi} proved the following multipartite version of Hajnal-Szemer\'edi theorem, which provides the best possible degree condition that guarantees a transversal $K_r$-factor in a balanced $r$-partite graph.%The following multi-partite version of Hajnal-Szemer\'edi theorem provides a degree condition that guarantees the existence of $K_r$-factors in a balanced $r$-partite graph. %, which follows from \cite[Lemma 6.3]{gan2024} by using Edmonds algorithm.{(why don't we use the exact degree condition of Keevash and Mycroft?)}

\begin{lemma}[\cite{KM-multi}]\label{lem:HS thm in multipartite graph} 
Given $r \in \mathbb{N}$ there exists an $n_0 \in \mathbb{N}$ such that the following holds. Suppose $G$ is an $r$-partite graph with vertex classes $V_1, \dots, V_r$ where $|V_i| = n \geq n_0$ for all $1 \leq i \leq r$. If
\[
\min_{i\in [r]}\min_{v\in V_i}\min_{j\in [r]\setminus \{i\}} |N(v)\cap V_j| \geq \Big(1 - \frac{1}{r}\Big)n + 1,
\]
then $G$ contains a $K_r$-factor.
\end{lemma}
We will use this to prove that a $K_r$-factor can also be found in an almost balanced partition,  provided that every large part {contains a matching with enough edges.} In fact, this enables us to first pick a few vertex-disjoint copies of $K_r$ each containing exactly one `extra' vertex in a large part (compared to the balanced copies) so as to obtain a balanced $r$-partite subgraph, to which Lemma~\ref{lem:HS thm in multipartite graph} would be applied. 

For any \( i \in [r] \), let \( \mathbf{u}_i\in \mathbb{R}^r \) be the \( i \)-th unit vector, i.e.\ \( \mathbf{u}_i \) has $1$ on the \( i \)-th coordinate and $0$ on the other coordinates.  Let \( \mathbf{1}_r\in \mathbb{R}^r \) be an $r$-vector {with all entries equal to $1$.}
% all of whose entries are $1$. 
Given a vertex partition $\mathcal{P}=\{V_1,\ldots,V_r\}$ of $G$ and $F\<G$,  the \textit{index vector} of \( F \)  (with respect to $\mathcal{P}$) is \( \mathbf{i}(F) = (x_1, x_2, \ldots, x_{r}) \) where \( x_i = |V(F) \cap V_i| \) for each \( i \in [r] \). %In the following, we give a sufficient condition that guarantees the existence of $K_r$-factors in $G$.
%Applying with Lemma~\ref{lem:HS thm in multipartite graph}, we have the following result to guarantee the existence of $K_r$-factors in $G$. 


%Without loss of generality, assume that $|A_i|=n+k_i$ for each $i\in [s]$ and $k_1\leq \cdots\leq k_s$. Note that $\sum_{i\in [s]}k_i=0$. We have the following result for unbalanced partition.
\begin{lemma}\label{lem:balance}
Let $\frac{1}{n}\ll \alpha \ll \beta \ll\frac{1}{r}$. Suppose that $G$ is a graph on $rn$ vertices with a partition $V(G) =\{A_1, A_2,\ldots, A_r\}$ satisfying 
\begin{enumerate}
[label =\rm  (B\arabic{enumi})]
    \item\label{A1} $ |A_1|\leq \cdots\leq |A_r|\leq n+\alpha n$;
    
    \item\label{A3} for any distinct $i,j \in [r]$ and any $v \in A_i$, we have $d(v, A_j) \ge (1-\beta)|A_j|$;
    \item\label{A2} for each $i \in [r]$, $G[A_i]$ contains a matching $M_i$ of size $k_i:=|A_i|-n$.
\end{enumerate}
{Then $G$ contains a $K_r$-factor $\mathcal{F}$ such that $\bigcup_{i\in [r]}E(M_i)\subseteq E(\mathcal{F})$,  and each $K_r$ in $\mathcal{F}$ satisfies    
\begin{itemize}
    \item either $|V(K_r)\cap A_i|=1$ for all $i\in [r]$;
    \item  or $|E(K_r)\cap E(M_i)|=1$ for exactly one $i\in [r]$ and $|V(K_r)\cap A_j|=1$ for every $j\in [r]\setminus \{i\}$ with $|A_j|\geq n$.
\end{itemize}}
 %{Moreover, each edge in the  matching of $G[A_i]$ lies in distinct copies of $K_r$ }
\end{lemma}


%We only consider the case that the following operations yield an unbalanced partition, otherwise we can greedily find a $K_r$-factor. The following section details how to adjust the partition satisfying Lemma~\ref{lem:balance} and how to find a $K_r$-factor in an unbalanced partition. For the current $(r,s)$-partition $\mathcal{P}_0$, we define the \textit{bad} vertices motivated by Lemma~\ref{lem:balance} and then move them to appropriate parts.


% we define the following types of vertices. 
% \begin{itemize}
%     \item For a vertex $x\in A_i$ with $i\in [s]$, we say $x$ is $(\beta, A_i)$-bad (or simply {$(\beta, A_i)$-bad}) if $d(x,A_i)\ge \beta n$. %Otherwise, we say $x$ is \textit{$(\beta,A_i)$-good}. 
%     Let $V_b(\mathcal{P},\beta,i)$ be the set of $(\beta,A_i)$-bad vertices.
%     \item For a vertex $x\in V(G)\setminus A_i$ with $i\in [s]$, we say that $x$ is $(\beta, A_i)$-exceptional (or simply {$(\beta, i)$-exceptional}) if $d(x,A_i)\le \beta n$. %Otherwise we say $x$ is \textit{$(\beta,i)$-acceptable}. 
%     Let $V_{ex}(\mathcal{P},\beta,i)$ be the set of $(\beta,i)$-exceptional vertices. 
%     \item For a vertex $x\in V(G)\setminus A_i$ with $i\in [s]$, we say that $x$ is $(\beta, A_i)$-excellent (or simply {$(\beta, i)$-excellent}) if $d(x,A_i)\ge |A_i|-\beta n$. Let $V_e(\mathcal{P}, \beta,i)$ be the set of $(\beta ,i)$-excellent vertices, and let $V_{ne}(\mathcal{P}, \beta,i):=V(G)\setminus (A_i\cup V_e(\mathcal{P}, \beta,i))$. % be the set of vertices which are not $(\beta, i)$-excellent vertices and not belonged to $A_i$. 
    
%     \item A vertex $x\in V(G)\setminus B$ is called $(\beta,B)$-excellent if $d(x,B)\ge |B|-\beta n$.  Denote by $V_e(\mathcal{P},\beta, B)$ the set of all such  $(\beta,B)$-excellent vertices, and let $V_{ne}(\mathcal{P},\beta, B):=V(G)\setminus (B\cup V_e(\mathcal{P}, \beta,B))$.
% \end{itemize}

%\end{definition}

%Since $G$ is $((r-1)n+1)$-regular, we estimate the number of bad vertices under the partition $\mathcal{P}^0$.




% For each $i\in [s]$, since $A_i^0$ is a $\gamma_i$-independent set in $G$ and $\gamma_i\ll \alpha \ll \beta$, we have 
% $$|V_b(\mathcal{P}_0,\beta,i)|\le\frac{2e(G[A_i^0])}{\beta n} \le  \frac{2\gamma_i n^2}{\beta n}\le \alpha n.$$ %Observe that at least $\binom{|A_i^0|}{2}-\gamma_i n^2$ edges are missing in  $G[A_i^0]$. 


%\begin{fact}
%For each $i\in [s]$, if $|A_i|>n$ then there exists a matching of size at least $\frac{k_i}{2\beta}$ where $k_i:=|A_i|-n$.
%\end{fact}

%\begin{proof}
%Note that $\Delta(G[A_i])\leq \beta n$ and $\de(G[A_i])\ge k_i+1$ as $G$ is $((r-1)n+1)$-regular.
%Let $M_i$ be a maximum matching of $A_i$. By Vizing's theorem, we have
%\begin{align*}
%|M_i|\ge \frac{e(G[A_i])}{\Delta(G[A_i])+1}\ge \frac{|A_i|\cdot (k_i+1)}{2(\beta n+1)}\ge \frac{k_i}{2\beta}.
%\end{align*}
%\end{proof}




%Now we are ready to prove Lemma~\ref{lem:balance}, which establishes the existence of $K_r$-factors.

\begin{proof}%[Proof of Lemma~\ref{lem:balance}]
Suppose that $|A_j|=|A_1|+a_j$ for each $j \in [2, r]$ and $|A_1|=a_1$.
Without loss of generality, {we may assume $A_1,A_2,\ldots,A_p$ are all the parts of size at most $n$ for  some $p\in[r]$}.
If $a_j=0$ for each $j \in [2, r]$, that is, $\{A_1,A_2, \ldots, A_r\}$ is a balanced partition, then Lemma~\ref{lem:HS thm in multipartite graph} implies that $G$ contains a desired $K_r$-factor.
Otherwise, recall that for each $i \in [p+1,r]$, $G[A_i]$ contains a matching $M_i$ of size $k_i:=|A_i|-n$, we proceed in two steps. 
\begin{itemize}
    \item[$(1)$]  For each $i \in [2, p]$, iteratively select $n - |A_i|$ vertex-disjoint copies of $K_r$ in $G$, {each of which} has an index vector $\mathbf{1}_r - \mathbf{u}_i + \mathbf{u}_\ell$ for some $\ell\in[p+1, r]$,  
provided that $M_{\ell}$ has an edge not appearing in previously selected copies of $K_r$.

    \item[$(2)$]  For each $i \in [p+1, r]$, pick $k_i - k_i'$ vertex-disjoint copies of $K_r$ with index vector $\mathbf{1}_r - \mathbf{u}_1 + \mathbf{u}_i$, where $k_i'$ is the total number of edges of $M_i$ used in step~$(1)$.
\end{itemize}

{Indeed, in step~$(1)$, it necessarily holds that $$\sum_{i\in [2,p]}(n-|A_i|)\leq \sum_{i\in [p]}(n-|A_i|)= \sum_{i\in [p+1,r]}(|A_i|-n)=\sum_{i\in [p+1,r]}k_i.$$
Meanwhile, the degree condition in~\ref{A3} allows a greedy vertex-selection of each $K_r$ starting with an edge (in both steps). The resulting $K_r$-tiling from above is denoted as $\mathcal{K}$. Note that  every edge in $M_i$ lies in exactly one  clique of $\mathcal{K}$, and such a clique  contains exactly  one vertex from each $A_j$ with $j\in [r]\setminus \{i\}$ and $|A_j|\geq n$.}  
% {(explain in details! Are there enough matching edges?)}
We now show that the resulting partition $\mathcal{P}'=\{A_1', \ldots, A_r'\}$ is balanced, where $A_i':=A_i\setminus V(\mathcal{K}).$
\begin{claim}\label{result-balance}
    For each $\ell\in [r]$, we have $|A_{\ell}'|=a_1 + \sum_{i\in [2,p]} k_i$.
\end{claim}
\begin{proof}
Notice that $\sum_{i\in [r]}k_i=0$ and $\sum_{i\in [p+1,r]}k_i'=\sum_{i\in [2,p]}(n-|A_i|)$. Hence %$-\sum_{i\in [p]}k_i=\sum_{i\in [p+1,r]}k_i$. We verify the number of remaining vertices in each part by   direct calculation. 
\begin{itemize}
    \item $|A_1'|= a_1-\sum_{i=2}^p(n-|A_i|)=a_1+\sum_{i=2}^pk_i$;
    \item for each $\ell\in [2,p]$, 
        \begin{align*}
        |A_{\ell}'|=&a_1+a_\ell-\sum_{i\in [2,p]\setminus\{\ell\}}(n-|A_i|)-\sum_{i\in [p+1,r]}(k_{i}-k_i')\\
        =&n-\sum_{i\in [p+1,r]}k_{i}=a_1-k_1+\sum_{i\in [p]}k_{i}=a_1+\sum_{i\in [2,p]}k_i;
        \end{align*}
        
    \item for each $\ell\in [p+1,r]$,                  \begin{align*}
        |A_{\ell}'|=&a_1+a_\ell-\sum_{i\in [2,p]}(n-|A_i|)-k_{\ell}'-\sum_{i\in [p+1,r]}(k_{i}-k_i')-(k_{\ell}-k_{\ell}')\\
        =&|A_{\ell}|-\sum_{i\in [p+1,r]}k_{i}-k_{\ell}=n-\sum_{i\in [p+1,r]}k_{i}=a_1+\sum_{i\in [2,p]}k_i,
        \end{align*}
\end{itemize}
as desired. 
\end{proof}

By~\ref{A1}, for each $i\in [r]$ one has 
$$|V(\mathcal{K})\cap A_i|\leq 2\sum_{i\in [2,p]}(n-|A_i|)+2\sum_{i\in [p+1,r]}(k_i-k_i')\leq 4\sum_{i\in [p+1,r]}k_i\leq 4r\alpha n.$$ % we remove at most $4\alpha n{??}$ vertices in each part $A_i$.
Hence $|A_i'|\geq |A_i|-4r\alpha n$. Together with~\ref{A3} and the choice  $\alpha\ll\beta\ll\frac{1}{r}$, for any distinct $i,j \in [r]$ and any $v \in A_i'$, we have
\[
d(v, A_i') \ge |A_i'|-{2\beta}n\ge \Big(1-\frac{1}{r}\Big)|A_i'|+1.
\]
Applying Theorem~\ref{lem:HS thm in multipartite graph} to $G - V(\mathcal{K})$ yields a $K_r$-factor {such that each copy of $K_r$ contains exactly one vertex from each $A_i$ with $i\in [r]$}. This together with $\mathcal{K}$ gives a desired  $K_r$-factor of $G$. 
\end{proof}
%We call a part $A_i\in \mathcal{P}$ a \textit{large} part if $|A_i|> \frac{|G|}{r}$, and a \textit{small} part otherwise. 
%In order to use Lemma~\ref{lem:balance} to find a \(K_r\)-factor in \(G\), we need to construct a graph satisfying all of \ref{A1}-\ref{A2}.  
To apply Lemma~\ref{lem:balance} in constructing a \(K_r\)-factor of \(G\), 
we shall construct a graph satisfying conditions~\ref{A1}-\ref{A2}. Let $\mathcal{P}=\{A_1,\ldots,A_s,B\}$ be a good partition of $G$. %{\sout{We call a copy of $K_r$ \textit{balanced} if it has index vector $(1,\cdots,1,r-s)$}}{(Yang:move it to where it is used)}. 
As outlined in Section~2, we proceed as follows. 

\begin{itemize}
    \item For~\ref{A3}, we construct a $K_r$-tiling $\mathcal{K}$ that covers all bad or exceptional vertices of $G$.
    \item For~\ref{A1}, we find a $\{K_{r-s},K_{r-s+1}\}$-factor in $G[B\setminus V(\mathcal{K})]$ and contract each copy to either an edge or a single vertex.
    \item For~\ref{A2}, the matching number is essentially derived from~\ref{B2} and the contraction as above.
\end{itemize}

%Note that in a good partition \(\mathcal{P} := \{A_1, \dotsc, A_s, B\}\) of $G$ as defined in Definition~\ref{def:good partition}, we have \(|B| \sim (r-s)|A_i|\) for each \(i \in [s]\), and the size of \(|A_i|\) may differ for distinct \(i \in [s]\).  
%As stated in Section 2, for \ref{A3}, we will find a $K_r$-tiling $\mathcal{K}$ that covers all bad or exceptional vertices in $G$; for \ref{A1}, we will find a $\{K_{r-s},K_{r-s+1}\}$-factor in $G[B]-V(\mathcal{T})$ and contract each of them into an edge or a single vertex; for \ref{A2}, it can be verified by using \ref{B2}.  %cover bad and exceptional vertices and then balance this good partition by removing a $K_r$-tiling of small size and contracting a $\{K_{r-s},K_{r-s+1}\}$-factor in the remaining vertices in $B$.}
% We balance this good partition and construct a \(K_r\)-factor in \(G\) using Lemma~\ref{lem:HS thm in multipartite graph}.
% the construction of a \(K_r\)-factor in \(G\) using Lemma~\ref{lem:HS thm in multipartite graph} consists of the following three steps:

% \begin{itemize}
%     \item Cover all bad or exceptional vertices with a $K_r$-tiling $\mathcal{T}$ to ensure the remaining graph meets the minimum degree condition of Lemma~\ref{lem:HS thm in multipartite graph}.
%     \item {Find a $\{K_{r-s},K_{r-s+1}\}$-factor in \(G[B \setminus V(\mathcal{T})]\) provided that $|B \setminus V(\mathcal{T})|$ is divisible by $r-s$. Then contract each copy of \(K_{r-s}\) into a single vertex and each copy of \(K_{r-s+1}\) into a new edge, obtaining a new vertex set \(B^*\).
%     We obtain an auxiliary graph with part $A_1, \dotsc, A_s$ and $B^*$ where each part has roughly equal size.}
%     % thereby reducing the problem to an $r$-partite setting where each part has roughly equal size.

% \item {Balance the sizes of all parts (including $A_1, \dotsc, A_s$ and $B^*$) via matchings, that is, removing a $K_r$-tiling with types $(1,\cdots,0\cdots2,\cdots,1,r-s)$ and $(1,\cdots,0,\cdots,1,r-s+1)$ such that the resulting $(s+1)$-partite graph induced on the remaining vertices fulfilling all requirements of Lemma~\ref{lem:HS thm in multipartite graph}.}
% \end{itemize}

%$|A_1^0|=\cdots=|A_s^0|=n$ and $|B_0|=(r-s)n$. In order to obtain a balance $(s+1)$-partite graph (may be minor adjustments based on the current partition $\mathcal{P}_0$), a natural idea is to contract each copy of $K_{r-s}$ in $B_0$ into a single vertex. In other words, we need to find a $K_{r-s}$-factor in $B_0$ provided the given conditions of degree and independent set. When $r-s\ge 3$, it follows from Theorem~\ref{thm:Non-extremal case r>3}. 

%In order to satisfy  \ref{A1}, we are to find a $K_{r-s}$-factor in $G[B]$.  
%{\sout{Notice that \ref{B05} and \ref{B07} hold. Thus, }
{For $r-s\geq 3$, Theorem~\ref{thm:Non-extremal case r>3} gives a $K_{r-s}$-factor in $G[B']$, where $B'$ denotes the largest subset of $B\setminus V(\mathcal{K})$ with $(r-s)\mid |B'|$.  %desired  When $r-s\ge 3$, by~\ref{B05} and~\ref{B07}, we obtain that Theorem~\ref{thm:Non-extremal case r>3}  implies that there is a $K_{r-s}$-factor in $G[B]-V(\mathcal{K})$.
For $r-s=2$, by applying the subsequent  result of~\cite{gan2024} to $G[B']$, we obtain that either $G[B']$ contains a perfect matching, or it is disconnected with  two odd components. 
%In the latter case, by applying some minor modification, % we construct two copies of $K_r$ with index vectors $(1,\dots,1,2,1)$ and $(1,\dots,1,0,3)$, or do some minor modification. Applying Lemma~\ref{thm:r-s=2} again 
In the latter case, after some minor adjustments, we can ensure that the remaining graph contains a perfect matching.} 
%Indeed, on the one hand,~\ref{B2} or Lemma~\ref{vertex-cover} gives us that there is an edge in some $A_i$ and we extend it to $H_1$ by using one vertex from one odd component in $G[B]$; on the other hand, we can find a copy of $K_3$ in the other odd component in $G[B]$ and we extend it to $H_2$.}
% Let $\epsilon > 0$. Suppose that $G$ is a graph on $2n$ vertices. We say that $G$ is $\epsilon$-\textit{close} to $2K_n$ if there exists a partition $A, B$ of $V(G)$ such that $|A| = n$, $|B| = n$ and $|E(G[A,B])|\le \epsilon n^2$.

\begin{lemma}[\cite{gan2024}]\label{thm:r-s=2}
Let $\beta>0$ and $n\in \mathbb{N}$ such that $\frac{1}{n}\ll\beta$. Suppose that $G$ is a graph on $2n$ vertices with $$\de(G)\ge (1-\beta)n.$$ Then at least one of the following holds:
\begin{itemize}
    \item $G$ contains a $2\beta$-independent set of size at least $n$;
    \item {$G$ is disconnected and consists of two odd components;}
    % $G$ is $3\beta$-close to $2K_n$, {where $n$ is odd};
    \item $G$ contains a perfect matching.
\end{itemize}
\end{lemma}
The following lemma allows us to find many vertex-disjoint copies of $K_r$ under a slightly weaker minimum degree condition. 
\begin{lemma}[\cite{Treg2015}]\label{SuperT}
    Let $\frac{1}{n} \ll \eps \ll \alpha \ll \frac{1}{r}$ where $r\in \mathbb{N}$ and $r\ge 2$. Suppose that $G$ is an $n$-vertex graph that contains at most $\eps n^r$ copies of $K_r$ and 
$$\delta(G)\ge \Big(1-\frac{1}{r-1}-\alpha\Big)n.$$
%and so that $G$ contains at most $\eps n^r$ copies of $K_r$. 
Then $G$ admits a $\sqrt{\alpha}$-independent set of size at least $\frac{n}{r-1}$.
\end{lemma}


%In the following sections, we are ready to construct a partition of independent sets firstly (in section 4.1) and then transform it into a balanced one through a series of appropriate operations (in section 4.2).


In the following, we will use Lemma~\ref{lem:balance} to prove Lemma~\ref{lem:covering bad vertices for balance}. 
%Suppose that $G$ admits a good partition.  In order to derive the model in Lemma~\ref{lem:balance}, we need to cover bad vertices that violate the degree condition~\ref{A3} by $K_r$-tiling. %Since $G$ is $((r-1)n+1)$-regular and $\Delta(G[A_i])\leq \beta n$ for each large part $A_i$, every vertex belonging to $A_i$ already meets condition~\ref{A3}. Thus, we only need to address two types of bad vertices:
%\begin{itemize}
%    \item vertices of large degree within small parts. Claim~\ref{fact:bad vertices for Ai}~\ref{large-degree vtx} implies that there are few bad vertices of the first type.
%    \item vertices in $B$ with many non-neighbors in some $A_i$, i.e., $(\beta,B)$-bad vertices. %By Lemma~\ref{good-partition}~\ref{B04}, every vertex $v\in B$ has at least $\beta n$ neighbors in $A_i$ for each $i\in[s]$. However, there may exist some vertices in $B$ which have many non-neighbors in some $A_i$, i.e., .
%\end{itemize}
%For such bad vertices, we attempt to cover them via a $K_r$-tiling so that the remaining partition satisfies a degree condition in Lemma~\ref{lem:balance}. A key requirement of this covering process is to ensure that the number of vertices that used in each part is balanced. 
%We first prove that the number of $(\beta,B)$-bad vertices is small. 


%Note that each vertex in large part satisfies condition~\ref{A3} as $G$ is $((r-1)n+1)$-regular and $\Delta(G[A_i])\leq \beta n$. Thus, we are left to handle only the two types of bad vertices: one is large-degree vertices in small parts and the other is vertices in $B$ that do not have a sufficiently high degree. Crucially, this covering process should guarantee that the resulting partition is balanced again. We first show that the number of those bad vertices is also small.

% 

%For such bad vertices, we attempt to cover them via a $K_r$-tiling so that the remaining partition satisfies a degree condition analogous to that in Lemma~\ref{lem:balance}. To achieve this goal, we define bad vertices in $B$ and show that their number is small.

%\begin{definition}\label{def:bad vertex for B}

%We say a vertex $v\in B$ is bad if $d_{A_i}(v)\le |A_i|-\beta n$ for some $i\in [s]$.
%\end{definition}
%counting bad vertices


% Suppose that the number of bad vertices in $B$ is at least $\alpha^{1/4} n$.
% By the Pigeonhole Principle, there exists some $i\in [s]$ such that there are at least $\frac{|B_b|}{s}$ vertices $v\in B$ satisfying $d_{A_i}(v)\le |A_i|-\beta n$.
% Since $d(v)=(r-1)n+1$ for each $v\in A_i$, we obtain that 
% $$|E(G[A_i])|\geq \frac{1}{2}\cdot\frac{\alpha^{1/4} n}{s} \cdot 2\alpha^{1/5}n\ge \sqrt{\alpha}n^2.$$ 
% Recall that $A_i^0$ is a $\gamma_i$-independent set of $G$ and at most $r\alpha n$ vertices are moved from other part to $A_i^0$. Hence 
% $$
% |E(G[A_i])|\leq (\gamma_i+2r\alpha) n^2<\sqrt{\alpha}n,
% $$
% a contradiction. 
% Therefore, 
% \begin{align}\label{size-of-ne}
%    |V_{ne}(3\alpha)^{1/5}^{1/5},B)|\le \alpha^{1/4} n\ \text{and}\ |V_{ne}(3\alpha)^{1/5}^{1/5})|\leq 2\alpha^{1/4}n.
% \end{align}






%The following lemma provides a method for covering the two types of bad vertices with a $K_r$-tiling and adjusting the current partition into a balanced one.

\begin{proof}[{\bf Proof of Lemma~\ref{lem:covering bad vertices for balance}}]
Choose $r\geq 3,n\in \mathbb{N}$ and constants satisfying %{replace $\gamma_1,\gamma_2$ with $\gamma,\gamma'$? } 
$$\frac{1}{n}\ll \alpha\ll\eps\ll\beta',  \beta\ll\gamma\ll \frac{1}{r}.$$ 
Assume that $G$ is an $rn$-vertex graph with $\de(G)\geq (r-1)n-\alpha n$ and $G$ admits an $(\alpha,\beta,\beta',\gamma)$-good partition $\mathcal{P}=\{A_1,\ldots,A_s,B\}$ {for some  $s\in [r]$}. Without loss of generality, assume that $|A_1|\leq |A_2|\leq \cdots\leq |A_s|$  and assume $A_1,\ldots,A_p$ are all the parts of size at most $n$ for some $p\in[s]$. 
For each $i \in [s]$, if $|A_i| > n$, then let $M_i$ be a matching of size $|A_i| - n + r$ in $G[A_i]$ (its existence follows from~\ref{B2}); {if $|A_s| = n$, then let $M_s$ be an edge in $G[A_s]$ containing at most one vertex that is not $(\alpha^{1/5},A_s)$-good (this also follows from~\ref{B2} together with the definition of $(\alpha^{1/5},A_s)$-good vertex and~\ref{B03});} %\sout{and $G[A_s]$ contains an edge whose both endpoints are $(\alpha^{1/5}, A_s)$-good}%\todo{This is not stated in $B_2$}, then let $M_s$ to be that edge;}
otherwise we let $M_i = \emptyset$. Define $U = \bigcup_{i \in [s]} V(M_i)$ as the union of the vertex sets of these matchings. By~\ref{B1}, one has $|U| \leq 3r\alpha n$. %Choose additional 
Our goal is to construct an auxiliary graph and a corresponding partition as required in Lemma~\ref{lem:balance}. 

\medskip
{\bf Step 1.~Greedily cover all bad or exceptional vertices using vertex-disjoint copies of $K_r$ whilst avoiding any vertex in $U$.}
\medskip%

Notice that $\Delta(G[A_i])\leq \beta' n$ if $|A_i|>n$. This together with the minimum degree condition gives that $d(v,D)\geq |D|-2{\beta'} n$ for any $v\in A_i$ with $|A_i|> n$ and any $D\in \mathcal{P}\setminus \{A_i\}$, {which satisfies~\ref{A3}}. It remains to consider the following two types of vertices:  %{in $B$ and $A_i$ with $|A_i|\le n$}:
\begin{enumerate}[label =\rm  {\bf Type \arabic{enumi}.}, ref = \rm {\bf Type \arabic{enumi}},  leftmargin=6em]
    \item\label{v1} $v\in A_i\setminus U$ and $v$ is not $(\alpha^{1/5},A_i)$-good  for some $i\in [s]$ with $|A_i|\leq n$;
    \item\label{v2} %$v\in B$ and $d(v,A_i)\leq |A_i|-\alpha^{1/5} n$ for some $i\in [s]$, namely, 
    $v\in B$ and $v$ is $(\alpha^{1/5},B)$-exceptional. 
\end{enumerate}




%Suppose for contradiction that there exists $v_i\in A_i\setminus U$ be of \ref{v1} that is not covered by $\mathcal{Q}$.
The following   claim guarantees the existence of a copy of $K_r$ that covers any fixed vertex while avoiding a prescribed vertex set.   
\begin{claim}\label{covering-v}
   Let $W\subseteq V(G)$ be a set of size at most $\frac{\beta}{3}n$ and $v_i$ be a vertex in $A_i\setminus W$ for some $i\in [s]$. Then there is a copy of $K_r$ in $G-W$ with index vector $(1,\ldots,1,r-s)$ containing $v_i$.  
\end{claim}
\begin{proof}
Let $v_i\in A_i\setminus W$ for some $i\in [s]$ {and choose \( j \in [s] \setminus \{i\} \). It follows from~\ref{B04} that} \( d(v_i, A_j) \geq \beta n \), and from \ref{B03} that at most \( 2\alpha n \) vertices in \( A_j \) are not \( (\alpha^{1/5}, A_j) \)-good. As $\alpha\ll\beta$, there are at least \( (\beta - 2\alpha -\frac{\beta}{3})n\ge \frac{\beta}{2}n \) vertices in \( N(v_i, A_j)\setminus W \) that are \( (\alpha^{1/5}, A_j) \)-good; we  fix one such vertex as \( v_j \). {Given} $\ell\in [s] \setminus \{i, j\}$, the minimum degree condition \( \delta(G) \geq (r-1)n - \alpha n \) implies \( d(v_j, A_\ell) \geq |A_\ell| - 2\alpha^{1/5} n \).
%{We pick exactly one \( (\alpha^{1/5}, A_\ell) \)-good vertex $v_{\ell}$ greedily for each \( \ell \in [s] \setminus \{i, j\} \) in the common neighbor of chosen vertices in $A_{\ell}$, while avoiding vertices in $W$. Since the number of previously chosen vertices is at most $s$, we have at least \( (\beta - 2s\alpha^{1/5} - 2\alpha-\frac{\beta}{3})n\ge \frac{\beta}{2}n \) choices for such good vertex in each step, which is adjacent to all previously fixed vertices.}
{Thus} at least \( (\beta - 2\alpha^{1/5} - 2\alpha-\frac{\beta}{3})n\ge \frac{\beta}{2}n \) vertices in \( N(\{v_i, v_j\}, A_\ell\setminus W) \) are \( (\alpha^{1/5}, A_\ell) \)-good; we fix one such vertex as \( v_{\ell} \). 
Iterating this process for all \( \ell \in [s] \setminus \{i,j\} \), one can take an \((\alpha^{1/5}, A_\ell) \)-good vertex \( v_{\ell}\in A_\ell\setminus W \), which is adjacent to all previously fixed vertices. This follows {from the fact that \( \alpha \ll \beta \), which guarantees at least \( \frac{\beta}{2}n \) available choices  at each step.}
Hence we obtain a clique on vertices, say $\{v_1,v_2,\ldots,v_s\}$.


%Choose \( j \in [s] \setminus \{i\} \) and from~\ref{B04} it holds that \( d(v_i, A_j) \geq \beta n \), and at most \( 2\alpha n \) vertices in \( A_j \) are not \( (\alpha^{1/5}, A_j) \)-good. Therefore, \( N(v_i, A_j)\setminus (U{\cup V(\mathcal{Q}))} \) contains at least \( (\beta - 2\alpha -5r^2\alpha)n\ge \frac{\beta}{2}n \) vertices that are \( (\alpha^{1/5}, A_j) \)-good, and fix one such vertex \( v_j \). For any \( \ell \in [s] \setminus \{i, j\} \), the minimum degree condition \( \delta(G) \geq (r-1)n - n^{0.6} \) gives \( d(v_j, A_\ell) \geq |A_\ell| - 2\alpha^{1/5} n \), so at least {\( (\beta - 2\alpha^{1/5} - 2\alpha-5r^2\alpha)n\ge \frac{\beta}{2}n \)} vertices in \( N(\{v_i, v_j\}, A_\ell\setminus (U\cup V(\mathcal{Q}))) \) are \( (\alpha^{1/5}, A_\ell) \)-good and we fix such a vertex as \( v_{\ell} \). Iterating this process for all \( \ell \in [s] \setminus \{i,j\} \), one can take a \((\alpha^{1/5}, A_\ell) \)-good vertex \( v_{\ell}\in A_\ell\setminus (U\cup V(\mathcal{Q})) \), which is adjacent to all previously fixed vertices. This follows as \( \alpha \ll \beta \), with at least \( \frac{\beta}{2}n \) choices available at each step. This yields a clique on vertices say $\{v_1,v_2,\ldots,v_s\}$.
%It follows from \ref{B04}-\ref{B05} that $d(v,D)\geq \beta n$ for each $D\in \mathcal{P}\setminus \{A_i\}$. 
%Choose $j\in [s]\setminus \{i\}$. It follows from  \ref{B2}-\ref{B04}  that $d(v,A_j)\geq \beta n$ and $A_j$ has at most $2\alpha n$ vertices that are not $(\alpha^{1/5},A_i)$-good. Hence there exists at least $(\beta-2\alpha)n$ vertices  in $N(v,A_j)$ that are  $(\alpha^{1/5},A_i)$-good. Let $v_j$ be such a vertex and $\ell\in [s]\setminus \{i,j\}$. Notice that $d(v_j,A_\ell)\geq |A_\ell|-2\alpha^{1/5}n$ as $\de(G)\geq (r-1)n-n^{0.6}$. Hence there are $(\beta-\alpha^{1/5}-4\alpha)n$ vertices in $N(\{v,v_j\},A_\ell)$ that are  $(\alpha^{1/5},A_\ell)$-good. Together with  $\alpha\ll \beta$, one may apply the above process iteratively to  get a copy of  $K_s$ that covers $v$ and $V(K_s)\cap A_j$ consists of a unique $(\alpha^{1/5},j)$-good vertex for each $j\in [s]\setminus \{i\}$. Moreover, each of such $(\alpha^{1/5},j)$-good vertex has at least $(\beta-(r-1)(\alpha^{1/5}+2\alpha))n\geq \frac{\beta}{2}n$ choices. 
Next, we are to find a copy of $K_{r-s}$ in $N(\{v_1,\ldots,v_s\},B) \setminus W$, and thus form a copy of $K_r$ covering the vertex $v_i$. Let $\Tilde{B}=N(\{v_1,\cdots,v_s\},B)\setminus W$. Similarly, 
$
|\Tilde{B}|\geq \frac{\beta}{2}n
.$ 
If $r-s\le 1$, then we are done. In what follows, we assume $r-s\ge 2$. Recall that $(r-s-r\alpha)n\leq |B|\leq (r-s+r\alpha)n$. % {and $|\Tilde{B}|\ge (r-s-1)n-2(s-1)\alpha^{1/5}n$}. 
Then~\ref{B05} implies that 
\begin{align*}
\de(G[B])&\ge (r-s-1)n-r\alpha n=\Big(1-\frac{1}{r-s}\Big)\cdot\big((r-s)n+r\alpha n\big)-\Big(2-\frac{1}{r-s}\Big)\cdot r\alpha n\\
&\ge \Big(1-\frac{1}{r-s}-\alpha^{1/5}\Big)\cdot|B|.
\end{align*}
%\medskip
%{\bf Covering bad vertices in $A_i$ with $|A_i|\le n$.}
%\medskip
%We cover each worse vertex $v_i$ in $A_i$ with a copy of vertex-disjoint $K_r$ if $d_{A_i}(v_i)\ge\alpha^{1/5} n$.
%For each bad vertex $v_i\in A_i$, the condition~\ref{B04} and~\ref{B05} in Lemma~\ref{lem:partition after adjust vertices} imply that $v_i$ has at least $\alpha^{1/5} n$ neighbors in each of the other parts of the partition. By Fact~\ref{fact:bad vertices for Ai}~\ref{large-degree vtx}-\ref{bad vtx2}, one can choose an $(\alpha^{1/5}, j)$-bad vertex $v_j$ in $A_j$ for each $j\neq i$ such that $v_1,v_2,\cdots,v_s$ induces a $K_s$ as $\alpha\ll\alpha^{1/5}$. Note that such vertex $v_j$ is adjacent to all but at most $\alpha^{1/5} n$ vertices in $V(G)\setminus A_j$. We now try to find a copy of $K_{r-s}$ in the common neighborhood of $v_1,v_2,\cdots,v_s$ in $B$. If $s=r$, then we are done since we have identified such $r$ vertices. Let $G'=N_B(\{v_1,v_2,\cdots,v_s\})$. Since $G$ is $((r-1)n+1)$-regular and $\alpha\ll\alpha^{1/5}$, we have
%\[
%\de(G[B])\ge (r-1)n+1-sn\ge \left(1-\frac{1}{r-s}-\alpha^{1/5}\right)\cdot|B|.
%\]
Let $x:=|B\setminus \Tilde{B}|$. {Since $v_j$ is $(\alpha^{1/5},A_j)$-good for each $j\in [s]\setminus \{i\}$, one has  
\begin{align}\label{eq:x}
    x\leq \Big|\bigcup_{i\in [s]}(B\setminus N(v_i))\Big|+|W|\leq  n+2(s-1)\alpha^{1/5} n+\frac{\beta}{3} n.
\end{align}
It follows that 
$|\Tilde{B}|\geq (r-s-1)n-\frac{\beta}{2}n$ as $\alpha\ll\beta$.} % , then $d(v_i,B)\ge |B|-n+1$ for each $i\in [s]$. Let $x$ be the number of non-neighbors of $v_1,v_2,\cdots,v_s$. As $v_j$ is a $(\beta,A_j)$-bad vertex in $A_j$ for each $i\neq j$, we have
%\[
%|B|-x=|G'|\ge |B|-n+1-(s-1)\alpha^{1/5} n,
%\]
%which means that $x\le n+(s-1)\beta n$.
Thus, \allowdisplaybreaks
%as $\alpha\ll\alpha^{1/5}\ll\eps\ll 1$ we have
\begin{align*}
\de(G[\Tilde{B}])&\ge \de(G[B])-x\ge \Big(1-\frac{1}{r-s}-\alpha^{1/5}\Big)\cdot|B|-x\\
&=|B|-x-\frac{|B|}{r-s}-\alpha^{1/5}|B|\\
&\ge|B|-x-\frac{(1+\beta)(|B|-x)}{r-s-1}-(|B|-x)\beta \\
&\geq\Big(1-\frac{1}{r-s-1}-2\beta\Big)\cdot(|B|-x)\\
&=\Big(1-\frac{1}{r-s-1}-2\beta\Big)\cdot|\Tilde{B}|,
\end{align*}
where the last inequality follows as % $x\le n+{2}(s-1)\alpha^{1/5} n$ and therefore 
$\frac{|B|}{r-s}\le\frac{(1+{\beta})(|B|-x)}{r-s-1}$ by~\eqref{eq:x}.
Based on~\ref{B07} and $\beta\ll \gamma$, we conclude that $G[\Tilde{B}]$ has no $\sqrt{\beta}$-independent set of size at least $\frac{|\Tilde{B}|}{r-s-1}$.  Applying Theorem~\ref{SuperT} with $(G[\Tilde{B}],\beta)$ in place of $(G,\alpha)$ yields a copy of $K_{r-s}$ in $G[\Tilde{B}]$. {This yields a copy of \(K_r\) with index vector $(1,\ldots,1,r-s)$ covering \(v_i\) in \(G-W\), as desired.} % contradicting the maximality of \(\mathcal{Q}\). Therefore, \(\mathcal{Q}\) is a \(K_r\)-tiling in $G-U$ that covers all vertices of \ref{v1}, and each copy of \(K_r\) in \(\mathcal{Q}\) has index vector \((1,\dots,1,r-s)\).} 
\end{proof}

By~\ref{B03}, there are at most $2r\alpha n$ vertices of~\ref{v1}. 
Combining Claim~\ref{covering-v} with the fact that $|U|\leq 3r\alpha n$, {we conclude that} there exists a \(K_r\)-tiling \(\mathcal{Q}\) in \(G-U\) covering all vertices of~\ref{v1}, such that each copy of \(K_r\) has index vector \((1,\dots,1,r-s)\) and contains exactly one vertex of~\ref{v1}.  
Then $|V(\mathcal{Q})|\le2r^2\alpha n$.
%{Let $\mathcal{Q}$ be a such $K_r$-tiling that covers all vertices of~\ref{v1} and  }
% as desired.%独立集版本的Turan thm

%If $s=r-1$, that is, we already have the vertices $v_1,v_2,\cdots,v_{r-1}$, then we only need to choose a vertex in the common neighborhood of $v_1,v_2,\cdots,v_{r-1}$ in $B$. Note that $B$ has no $\gamma_{r}$-independent set and $d(v,B)\geq \beta n$ for each vertex $v\in V(G)$. As $v_j\in A_j$ is a $(\beta,A_j)$-bad vertex for each $j\neq i$ and $\beta$, we have \[
%|G'|=d_B(v_1,v_2,\cdots,v_{r-1})\ge\beta n-(r-1) \alpha^{1/5}n>\frac{\beta n}{2}>0,
%\]
%which implies that we can find $v_r\in B$ such that $\{v_1,v_2,\cdots,v_r\}$ induced a $K_r$ in $G$.

 %{For any $v\in A_i$ of~\ref{v1}, since there are at least $\frac{\beta}{2}n-2r\alpha n$ available vertices in each $A_{\ell}$ with $\ell\in [s]\setminus \{i\}$,}
%As $\alpha\ll\beta$, the above process yields a $K_r$-tiling $\mathcal{Q}$ that covers all vertices of~\ref{v1} and avoids vertices in $U$. Furthermore, each $K_r$ in $\mathcal{Q}$ has index vector $(1,\ldots,1,r-s)$.

%\medskip
%{\bf Covering bad vertices in $B$.}
%\medskip
Now, we consider vertices of~\ref{v2}. %Recall that in \ref{B06} there are at most $r\alpha n$ $(\alpha^{1/5},B)$-exceptional vertices in $B$.  %$(\alpha^{1/5},B)$-exceptional vertices. 
%The subsequent claim ensures that any vertex in $B$ lies in  a copy of  $K_r$ in $G[B\setminus (U\cup V(\mathcal{Q}))]$.

\begin{claim}\label{claim:good Kr-s in B}
   Let $W\subseteq V(G)$ be a set of size at most $\frac{\beta}{3}n$. For any $w\in B\setminus W$, there is a copy of $K_r$ in $G-W$ with index vector $(1,\ldots,1,r-s)$ containing $w$.  
\end{claim}

%\begin{claim}\label{claim:good Kr-s in B}
%For any $W\subseteq B$ with $|W|\le\beta n$, every vertex $w \in B \setminus W$ lies in a copy of $K_{r-s}$ that is disjoint from $W$.
%\end{claim}
\begin{proof}
Let $w\in B\setminus W$. We first show that the vertex $w$ lies in a copy of $K_{r-s}$ in $G[B\setminus W]$. 
The case for $r-s\le 1$ is trivial, so we only consider $r-s\geq 2$. Let $G_1$ be a graph obtained from $G[B\setminus W]$ by deleting all $(\alpha^{1/5},B)$-exceptional vertices in $G$. 
It suffices to find a copy of $K_{r-s-1}$ in $G_1[N_{G_1}(w)]$. Recall that in~\ref{B06} there are at most $r\alpha n$ $(\alpha^{1/5},B)$-exceptional vertices in $B$. Hence  
$$|N_{G_1}(w)|\ge \de(G[B])-\frac{\beta}{3} n-r\alpha n\ge \Big(1-\frac{1}{r-s}-2\beta\Big)\cdot{|B\setminus W|}.$$
For any vertex $u\in N_{G_1}(w)$, as $\beta\ll\frac{1}{r}$, we have 
\begin{align*}
|N_{G_1}(w)\cap N_{G_1}(u)|
&\ge |N_{G_1}(w)|+|N_{G_1}(u)|-|G_1|\\
&\ge |N_{G_1}(w)|+\Big(1-\frac{1}{r-s}-2\beta\Big)\cdot{|B\setminus W|}-|B\setminus W|\\
&= |N_{G_1}(w)|-\Big(\frac{1}{r-s}+2\beta\Big)\cdot{|B\setminus W|}\\
&\geq \Big(1-\frac{1}{r-s-2}-\beta\Big)\cdot|N_{G_1}(w)|,
\end{align*}
which implies $\de(G_1[N_{G_1}(w)])\ge \big(1-\frac{1}{r-s-2}-\beta\big)|N_{G_1}(w)|$. {Based on~\ref{B07} and $\beta\ll \gamma$, we obtain that $G_1[N_{G_1}(w)]$ contains no $\sqrt{\beta}$-independent set of size $\frac{|N_{G_1}(w)|}{r-s-2}$.  
Therefore, by applying Theorem~\ref{SuperT}  with $( G_1[N_{G_1}(w)], \beta)$ in place of $(G, \alpha)$,} there is a copy of $K_{r-s-1}$ in $G_1[N_{G_1}(w)]$, which together with $w$ forms a {copy of} $K_{r-s}$ in $G_1$. Let $T$ be {one such} $K_{r-s}$ with vertex set $\{w,u_1,\ldots,u_{r-s-1}\}$. 


Observe that $d(u_i,A_j)\geq |A_j|-\alpha^{1/5}n$ for each $i\in [r-s-1]$ and each $j\in [s]$, and $d(w,A_j)\geq \beta n$ for each $j\in [s]$ by~\ref{B04}. By a similar construction with $\mathcal{Q}$ and the fact that $\alpha\ll \beta$, we can extend \( T\) into a copy of \( K_r \) with index vector $(1,\ldots,1,r-s)$ while avoiding $W$. %, denoted as $\mathcal{T}'$, such that for each $i\in [s]$, each \(K\in \mathcal{T}'\) consists of a single \( (\alpha^{1/5}, A_i) \)-good vertex and $V(K)\cap (U\cup V(\mathcal{Q}))=\emptyset$.
\end{proof}

{Recall that $|V(\mathcal{Q})|+|U|\leq 2r^2\alpha n+3r\alpha n$ and there are at most $r\alpha n$ $(\alpha^{1/5},B)$-exceptional vertices in $B$.  Hence, Claim~\ref{claim:good Kr-s in B} yields a $K_r$-tiling $\mathcal{T}$ in $G$ that avoids $V(\mathcal{Q})\cup U$ and covers all vertices of~\ref{v2}, where each copy of $K_r$ has index vector $(1,\ldots,1,r-s)$.} %Moreover, each clique in $\mathcal{T}$ contains exactly one $(\alpha^{1/5},B)$-exceptional vertex. 
%Let $T\in\mathcal{T}$ with vertex set $\{v,u_1,\ldots,u_{r-s-1}\}$,  where $v$ is the $(\alpha^{1/5},B)$-exceptional vertex in $T$.  %For the $(\alpha^{1/5},B)$-exceptional vertex $v\in B$, let $K^v$ be the $K_{r-s}$ copy that covers $v$, with $V(K^v)=\{v,u_1,\ldots,u_{r-s-1}\}$. 
%Observe that $d(u_i,A_j)\geq |A_j|-\alpha^{1/5}n$ for each $i\in [r-s-1]$ and each $j\in [s]$ and from \ref{B04} that $d(v,A_j)\geq \beta n$ for each $j\in [s]$. By a similar construction with $\mathcal{Q}$ and the fact $\alpha\ll \beta$, we can extend all \( T\in \mathcal{T}\) into pairwise vertex-disjoint copies of \( K_r \), denoted as $\mathcal{T}'$, such that for each $i\in [s]$, each \(K\in \mathcal{T}'\) consists of a single \( (\alpha^{1/5}, A_i) \)-good vertex and $V(K)\cap (U\cup V(\mathcal{Q}))=\emptyset$. %By \ref{B07}, the number of $(\alpha^{1/5},B)$-exceptional vertex is at most $r\alpha n$. 
%Consequently, there is a \(K_r\)-tiling \(\mathcal{T}\) in \(G - V(\mathcal{Q}) - U\) that covers \(T\), and every copy of \(K_r\) in \(\mathcal{T}\) has index vector \((1,\dots,1,r-s)\). %K_{r-s}$ copies in $H$  there is a $K_r$-tiling that covers  Moreover, the \( K_r \)-copies corresponding to distinct \( (\alpha^{1/5}, B) \)-exceptional vertices are pairwise disjoint.
%For each $(\alpha^{1/5},B)$-exceptional vertex $v\in B$, Claim~\ref{claim:good Kr-s in B} implies that there exists a copy of $K_{r-s}$ covering $v$ in $B$ external to the vertices already used in the covering. Moreover, the other $r-s-1$ vertices in this $K_{r-s}$ are all good and we denote them by $u_1,\cdots,u_{r-s-1}$. Note that each $u_i$ has at least $|A_i|-\alpha^{1/5}n$ vertices in $A_i$. Combing Lemma~\ref{lem:partition after adjust vertices}~\ref{B04} and Fact~\ref{fact:bad vertices for Ai}~\ref{large-degree vtx}\ref{bad vtx3}, for each $j\in[s]$ we can choose a $(\beta,A_j)$-bad vertex $v_j\in A_j$ in the common neighbor of $v,u_1,\cdots,u_{r-s-1}$ such that $\{v_1,v_2,\cdots,v_s\}$ induces a $K_s$ as $\alpha\ll\beta$. Therefore, $\{v_1,\cdots,v_s,v,u_1,\cdots,u_{r-s-1}\}$ is a copy of $K_r$ covering $v$. 
Let $\mathcal{K}=\mathcal{T}\cup\mathcal{Q}$. Then $\mathcal{K}$ has index vector $|\mathcal{K}|(1,1,\ldots, r-s)$ and  $|V(\mathcal{K})\cup U|\leq 4r^2\alpha n$. %Let $\mathcal{P}':=\{A_1',\ldots,A_s',B'\}$ be a partition of $G':=G-V(\mathcal{K})$, which is obtained from $\mathcal{P}$ by deleting all vertices in $V(\mathcal{K})$.


\medskip
{\bf Step 2. Contracting disjoint copies of $K_{r-s}$ and $K_{r-s+1}$ in $G[B\setminus V(\mathcal{K})]$.}% so as to achieve~\ref{A1} in the resulting auxilary graph.}
\medskip

\noindent \textbf {Step 2.1. Fixing divisibility and balancing.
}

Note that it may happen that $(r-s) \nmid |B \setminus V(\mathcal{K})|$. Assume that $r-s\ge2$\footnote{The case $r-s=1$ is much simpler.} and $|B\setminus V(\mathcal{K})|\equiv q \pmod{(r-s)}$. %The final step is to deal with the case that the number of vertices in $B'$ is not divisible by $r-s$.  %\medskip
%{\bf Covering non-divisible vertices in $B$.}
%\medskip
%Fact~\ref{fact:bad vertices for Ai} and Fact~\ref{fact:bad vertices in B} imply that the covering process in the first two steps requires no more than $10r\alpha n$ vertices.
By~\ref{B05} and $\alpha\ll\beta$, we have
\begin{align}\label{mini-degree-K}
\de(G[B\setminus V(\mathcal{K})])\ge (r-s-1)n-r\alpha n-4r^2\alpha n\ge \Big(1-\frac{1}{r-s}-\beta\Big)\cdot|B\setminus V(\mathcal{K})|.
\end{align}
As $\eps\ll\beta$, applying Theorem~\ref{SuperT} with $(G[B \setminus V(\mathcal{K})],\beta)$ in place of $(G,\alpha)$ yields $\eps n$ vertex-disjoint copies of $K_{r-s+1}$, denoted as $\mathcal{H}_0$, in $G[B \setminus V(\mathcal{K})]$.  Let $H$ be a copy of $K_{r-s+1}$ in $\mathcal{H}_0$. %Denote $B''=B'\setminus V(H)$.  
Note that each vertex $v$ in $B\setminus V(\mathcal{K})$ satisfies $d(v,A_i)\ge |A_i|-\alpha^{1/5} n$ for all $i\in [s]$. 
{Thus, for any $j\in [s]$, we can extend $H$ into $K_r$ by iteratively select an $(\alpha^{1/5}, A_{i})$-good vertex $v_i\in A_i$ for all $i\in [s]\setminus\{j\}$, whilst avoiding the vertices in $V(\mathcal{K})\cup U$.}
%\sout{Thus, by avoiding vertices in $V(\mathcal{K})\cup U$, for any $j\in [s]$, we can iteratively select vertices \(v_1 \in A_1, \dotsc,v_{j-1}\in A_{j-1},v_{j+1}\in A_{j+1},\ldots, v_{s} \in A_{s}\) such that each \(v_i\) is in the common neighborhood of $V(K')$ and all previously chosen vertices.}
{This is possible because each chosen vertex \(v_i\) is \((\alpha^{1/5},A_i)\)-good, and the number of candidates for \(v_i\) is at least
\[
|A_i| - (r-s+1)\alpha^{1/5}n - s\alpha^{1/5}n-|V(\mathcal{K})\cup U| \ge \frac{n}{2}
.\]
}
%\sout{This can be achieved as every vertex $v_i$ is $(\alpha^{1/5}, A_{i})$-good for each $i\in [s]\setminus \{j\}$.}
%Then \(V(K') \cup \{v_i:i\in[s]\setminus \{j\}\}\) induces a \(K_r\) in \(G\). 
As $q < r-s$, we extend $q$ copies of $K_{r-s+1}$ from $\mathcal{H}_0$ as above into a $K_r$-tiling $\mathcal{H}$ and moreover%\footnote{We remark here that the choice of $\mathcal{H}$ would ensure that every large part from the final partition actually corresponds to a large part from the original $\mathcal{P}$.},
\begin{enumerate}
[label =\rm  (C\arabic{enumi})]
    \item\label{E1} if $|A_s| > n$, then $V(\mathcal{H}) \cap A_s = \emptyset$;
    \item\label{E2} if $|A_s|\leq n$, % and $p\in [s]$ is the largest integer such that $|A_{p}|\leq n$, 
   %\sout{ then $|V(\mathcal{H}) \cap A_i| = q$ for each $i\in [p+1,s]$, and}{(in this case $p=s$)} 
   then $|V(\mathcal{H}) \cap A_i| = q-q_i$ for $i\in [s]$ and some nonnegative integer $q_i$ such that $q_i\leq n-|A_i|$ and $\sum_{i\in [p]}q_i=q$.
    \end{enumerate}
%{Note that $|A_i|-(q-q_i)\le n-q$ with each $|A_i|<n$.}
Indeed, if $|A_s|\leq n$, then $p=s$ and  $\sum_{i\in [s]}(n-|A_i|)=|B|-(r-s)n\geq q$. Hence, for each $i\in [s]$, there exists $q_i\in \mathbb{N}$ such that $q_i\leq n-|A_i|$ and $\sum_{i\in [p]}q_i=q$. Thus, it suffices to choose for each $i\in [s]$,  $q_i$ copies of $K_r$ such that each of them contains no vertices from $A_i$. 



Let $\hat{\mathcal{P}}:=\{\hat{A_1},\ldots,\hat{A_s},\hat{B}\}$ be the resulting partition of $\hat{G}:=G-V(\mathcal{K}\cup \mathcal{H})$, each part being obtained by removing all vertices in $V(\mathcal{K}\cup \mathcal{H})$. Then, $(r-s)\mid |\hat{B}|$. {Moreover,
\begin{itemize}
    \item if $|A_s|>n$, then $|\hat{A_i}|-\frac{|\hat{G}|}{r}=|A_i|-n$ for each $i\in [s-1]$ and $|\hat{A_s}|-\frac{|\hat{G}|}{r}=|A_s|-n+q$;
    \item if $|A_s|\leq n$, then $|\hat{A_i}|-\frac{|\hat{G}|}{r}=|A_i|-n+q_i\leq 0$ for each $i\in [s]$. % and $|\hat{A_i}|-\frac{|\hat{G}|}{r}=|A_i|-n$ for each $i\in [p+1,s]$. 
\end{itemize} 
It follows that $|\hat{A_i}|>\frac{|\hat{G}|}{r}$ {if and only if} $|A_i|>n$ for each $i\in [s]$.} %Therefore, $|\hat{A_i}|\leq \frac{|\hat{G}|}{r}$ for all $i\in [p]$.}

{Let $a:=\frac{|\hat{G}|}{r}-\frac{|\hat{B}|}{r-s}$. Suppose that  $a>0$. Note that $a\leq 2r\alpha n$ and  
$$
\sum_{i\in [p+1,s]}\Big(|\hat{A_i}|-\frac{|\hat{G}|}{r}\Big)\geq\sum_{i\in [s]}\Big(|\hat{A_i}|-\frac{|\hat{G}|}{r}\Big)=|\hat{G}|-|\hat{B}|-s\frac{|\hat{G}|}{r}=(r-s)\Big(\frac{|\hat{G}|}{r}-\frac{|\hat{B}|}{r-s}\Big)=(r-s)a.
$$
%Thus, 
%$$
%\sum_{i\in [p+1,s]}\big(|A_i|-n+q\big)\geq \sum_{i\in [p+1,s]}\Big(|\hat{A_i}|-\frac{|\hat{G}|}{r}\Big)\geq (r-s)a.
%$$
Recall that for each $i\in [p+1,s]$,  $G[\hat{A_i}\cap U]$ contains a matching of size $|A_i|-n+r\geq |\hat{A_i}|-\frac{|\hat{G}|}{r}+r-q$. For  each such $i$, choose $p_i\in \mathbb{N}$ such that $p_i\leq |\hat{A_i}|-\frac{|\hat{G}|}{r}$ and  $\sum_{i\in [p+1,s]}p_i=(r-s)a$. Let $U'\subseteq U$ be a subset for which each induced subgraph $G[\hat{A_i}\cap U']$ contains $p_i$ edges. %Notice that every vertex in $A_i\cap U'$ has at least $|D|-2\beta' n$ neighbors in every part $D\in \mathcal{P}\setminus \{A_i\}$.
Hence, by the same argument as in Claim \ref{covering-v}, one may find a $K_r$-tiling $\mathcal{R}$ consisting of exactly $(r-s)a$ copies of $K_r$ each containing one edge in $G[U']$, $r-s-1$ vertices from $B$ and at least one vertex from each $A_i$ with $i\in[s]$, whilst avoiding $V(\mathcal{K}\cup \mathcal{H})\cup (U\setminus U')$.  If $a\leq 0$, then take $\mathcal{R}=\emptyset$ and $p_i=0$ for each $i\in [p+1,s]$. Thus, $|V(\mathcal{K}\cup \mathcal{H}\cup \mathcal{R})|\leq 3r^3\alpha n$.}

{Let $\mathcal{P}':=\{A_1',\ldots,A_s',B'\}$ be the resulting partition of $G':=G-V(\mathcal{K}\cup \mathcal{H}\cup \mathcal{R})$, each part being obtained by removing all vertices in $V(\mathcal{K}\cup \mathcal{H}\cup \mathcal{R})$. Then, $(r-s)\mid |B'|$. 
Furthermore, if $a>0$, we have 
\begin{align}\notag
&\frac{|B'|}{r-s}=\frac{|\hat{B}|-(r-s-1)(r-s)a}{r-s}=\frac{|\hat{G}|}{r}-(r-s)a=\frac{|G'|}{r},
\\\notag
&|A_i'|=|\hat{A_i}|-(r-s)a\leq \frac{|\hat{G}|}{r}-(r-s)a=\frac{|G'|}{r}\ \text{for}\ i\in [p],\\\notag
&|A_i'|=|\hat{A_i}|-p_i-(r-s)a\ge \frac{|\hat{G}|}{r}-(r-s)a=\frac{|G'|}{r}\ \text{for}\ i\in [p+1,s].
\end{align}
Thus, in a summary, we have
\begin{align}\label{G'}
\frac{|B'|}{r-s}\ge \frac{|G'|}{r},\ |A_i'|\leq \frac{|G'|}{r}\ \text{for}\ i\in [p],\ \text{and}\ 
\frac{|G'|}{r}\leq |A_i'|\leq |\hat{A_i}|-\frac{|V(\mathcal{R})|}{r}\ \text{for}\ i\in [p+1,s]. 
\end{align}
Moreover, for each $i\in [p+1,s]$, $G'[A_i'\cap U]$ contains a matching of size} \begin{align}\label{matching'}
    |A_i|-n+r-p_i\geq |\hat{A_i}|-\frac{|\hat{G}|}{r}-p_i+r-q\geq  |A_i'|-\frac{|G'|}{r}+1.
\end{align}
%= |A_i|-n+q-p_i+(r-q)
%\footnote{Actually plus $r-q$}
%\begin{itemize}
%    \item if $|A_s| > n$, then $V(\mathcal{H}) \cap A_s = \emptyset$;
%    \item if $|A_s| <n$, then $|V(\mathcal{H})\cap A_i|=q-1\ \text{if}\ i\in [t],\ \text{and}\ |V(\mathcal{H})\cap A_i|=q\ \text{if}\ i\in [t+1,s], $ where $q(s-1)\equiv t\pmod{s}$;
%\item if $|A_s|=n$, then $|V(\mathcal{H})\cap A_i|=q-(n-|A_i|)$ for each $i\in [s]$. % and $p\in [p]$ is the largest integer such that $|A_{p}|<n$, then $|V(\mathcal{H})\cap A_i|=q-(n-|A_i|)$ for each $i\in [p]$ and $|V(\mathcal{H})\cap A_i|=q$ for each $i\in [p+1,s]$. 
%$0 \le |V(\mathcal{H}) \cap A_j| - |V(\mathcal{H}) \cap A_i| \le 1$ for all $1 \le i \le j \le s$.
%\end{itemize}
%Notice that when $|A_s|\leq n$, we have 
%$$
%|V(\mathcal{H})\cap A_i|=q-1\ \text{if}\ i\in [1,t],\ \text{and}\ |V(\mathcal{H})\cap A_i|=q\ \text{if}\ i\in [t+1,s],
%$$
   %\medskip %After the covering procedures, we obtain a $K_r$-tiling of size at most $10r\alpha n$ in $G$ covering all kinds of bad vertices. We remove such $K_r$-tiling and get a new unbalanced partition $\mathcal{P}$.  To obtain a balanced one as in Lemma~\ref{lem:balance}, we aim to find a $K_{r-s}$-factor in $B$ and identify each copy of $K_{r-s}$ with a single vertex. Let $A_{s+1}$ be a resulting set of size exactly $\frac{|B|}{r-s}$.
%{%If $B'$ is a large part in $\mathcal{P}'$, that is, $|B'|>(r-s)\frac{|G'|}{r}$, then we shall take extra copies of $K_{r-s+1}$ from $H$ which would be contracted into edges so as to achieve \ref{A2}.}
% Consider the case that $|B'|>(r-s)\frac{|G'|}{r}$.
%Before constructing a $K_{r-s}$-factor in $G'[B]$, we shall notice the possibility that $|B'|>(r-s)\frac{|G'|}{r}$. 
% To balance the size of each part, we find some copies of $K_{r-s+1}$ in $G'[B']$; these play a role analogous to that of the matching edges in {$G'[A_i']$}.

Let $b:=\frac{|B'|}{r-s}- \frac{|G'|}{r}$. Clearly, $b\geq 0$. If $b>0$, then we shall take extra copies of $K_{r-s+1}$ from $\mathcal{H}_0$ which would be contracted into edges so as to achieve~\ref{A2}. It follows from~\ref{B1} that $b\leq r\alpha n$.  
%If $r-s=1$, then there are at least $\eps n\ge (r-s)b$ matching edges in $G'[B']$ as {$e(G[B])\ge \gamma_2n^2$ and $\eps \ll \gamma_2$}.
%If $r-s\ge 2$, t
Then as $\alpha \ll \eps$, we can greedily take $(r-s)b$ vertex-disjoint copies of $K_{r-s+1}$ from $\mathcal{H}_0-V(\mathcal{H}\cup \mathcal{R})$, denoted as $\mathcal{F}$. If $b= 0$, then take $\mathcal{F}=\emptyset$. Let $B'':=B'\setminus V(\mathcal{F})$. Clearly, $(r-s)\mid |B''|$. 



%{Tell what you are doing next}
\medskip
\noindent {\bf Step 2.2. Construct a $K_{r-s}$-factor in $G'[B'']$.} 

The case $r-s=1$ is trivial. For $r-s\ge 2$, it follows from~\ref{B07} that $G'[B'']$ admits no ${\gamma}^2$-independent set of size $\frac{|B''|}{r-s}$. It is routine to check that 
\begin{align}\label{mini-degree}  \de(G'[B''])\ge \Big(1-\frac{1}{r-s}-\beta\Big)\cdot|B''|.
\end{align}
If $r-s\ge 3$, then Theorem~\ref{thm:Non-extremal case r>3} implies that there is a $K_{r-s}$-factor in $G'[B'']$ {since $\beta\ll\gamma$}, as desired.

Now, we consider \( r - s = 2 \). By Lemma~\ref{thm:r-s=2}, we conclude that either \(G'[B'']\) contains a perfect matching (in which case we are done), or it is disconnected with two odd components. %In fact, %we may assume the former occurs; otherwise 
%if the latter case holds, then a perfect matching can be obtained via a minor modification of the subset $B''$; as explained below.
Suppose we are in the latter case that $G'[B'']=G_1\cup G_2$, where both $|G_1|$ and $|G_2|$ are odd. It follows from~\eqref{mini-degree} that $|G_1|,|G_2|\geq \delta(G'[B''])\geq n-3\beta n$. Then each $G_i$ contains a triangle, say $T_i$ for $i\in [2]$.
Moreover,  by the argument in Claim~\ref{claim:good Kr-s in B}, each $T_i$ can be extended to a copy of \(K_r\), denoted by {\(K^{i,j}\)}, by using exactly one $(\alpha^{1/5},A_\ell)$-good vertex from each \({A_{\ell}'}\) with {\(\ell \in [s]\setminus\{j\}\)}, whilst avoiding vertices in ${U}$. Next we proceed according to the following two possible cases. 

\medskip
{\bf Case 1. $|A_s|\geq n$.} % or $|A_s|=n$ and $A_s\cap U$ contains no vertices of \ref{v1}.}
\medskip

Recall that $G'[A_s'\cap U]$ contains a matching of size {$|A_s'|-\frac{|G'|}{r}+1>0$} when $|A_s|>n$. Together with \ref{B2}, there is an edge $v_sv_s'$ in $G'[A_s'\cap U]$. {Notice that $A_s'\cap U$ may contain a vertex that is not 
$(\alpha^{1/5},A_s)$-good  if $|A_s|=n$.  Without loss of generality, if such a vertex exists, let it be $v_s$. %Notice that  $\Delta(G[A_s])\leq \beta'n$ if $|A_s|>n$.
By~\ref{B2} and~\ref{B04}, it holds that   $d(v_s,A_i)\geq \beta  n$ and $d(v_s',A_i)\geq |A_i|-2\beta' n$ for all $i\in [s-1]$.}  Following the argument in Claim~\ref{covering-v}, by avoiding vertices in $V(\{K^{1,s},K^{2,s}\})\cup U$, we can find an \((\alpha^{1/5},A_j)\)-good vertex   
\(v_j\in A_j'\) for each \(j\in[s-1]\) and a vertex \(w\in B''\) such that  
\(G'[\{v_1,\dots,v_{s-1},v_s,v_s',w\}]\) forms a clique \(K^1\).  Without loss of generality, assume that $w\in V(G_1)$. We then add $K^1\cup K^{2,s}$ back into $\mathcal{K}$. %This preserves the `balance' condition $2|V(K^1\cup K^{2,s})\cap A_i| = |V(K^1\cup K^{2,s})\cap B|$ for every $i\in [s]$.
By Lemma~\ref{thm:r-s=2}, the updated subgraph $G'[B''] = (G_1 - w) \cup (G_2 - V(K^{2,s}))$  contains a perfect matching, as desired.

%\medskip
%{\bf Case 2. $|A_s|=n$ and $A_s\cap U$ contains a vertex of \ref{v1}.}
%\medskip


%Suppose that $v_sv_s'$ is an edge in $G[A_s\cap U]$, where $v_s$ is the vertex  of~\ref{v1} and $v_s$ is $(\alpha^{1/5},A_s)$-good.  
%{Recall that $\mathcal{Q}$ is a $K_r$-tiling covering all vertices of~\ref{v1}.}
%Then $\mathcal{Q}$ contains a copy of $K_r$, say $K^2$ containing $v_s$. Note that $d(v_s,A_s)\geq \alpha^{1/5}n$. 
%{Together with~\ref{B03} and the fact $|V(\mathcal{K\cup H}_1)\cup U|\le 4r^2\alpha n$}, there is an  $(\alpha^{1/5},A_s)$-good vertex $v_s'\in (A_i\cap N(v_i))\setminus V(\mathcal{K}\cup \mathcal{H})$. 
%Again as in Step~1, by avoiding vertices in \(V(\mathcal{K} \cup \mathcal{H}\cup\{K^{1,s},K^{2,s}\})\), we can find an $(\alpha^{1/5},A_j)$-good vertex $v_j\in A_j$ for each $j\in [s-1]$ and a vertex $w\in B''$ such that $G[\{v_1,\ldots,v_s,v_s',w\}]$ induces a clique $K_r$, denoted as $K^2$. Without loss of generality, assume that $w\in V(G_1)$. We then add $K^2\cup K^{2,s}$ to $\mathcal{Q}$. This replacement also preserves the balance condition $2|V(K^2\cup K^{2,s})\cap A_i| = |V(K^2\cup K^{2,s})\cap B|$ for every $i\in [s]$.
%Now, the updated subgraph $G'[B''] = (G_1 - w) \cup (G_2 - V(K^{2,s})) \cup G[B\cap V(K^2)]$  contains a perfect matching by Lemma~\ref{thm:r-s=2}, as desired. % Consequently, by Lemma~\ref{thm:r-s=2}, the graph $G'[B'']$ contains a perfect matching, as desired. % Then the above $K_r$ together with $K^2$ are % By a similar discussion as above, we obtain that there are two vertex-disjoint copies of $K_r$ avoiding vertices in $B\cap V(K^4)$ and two cliques with index vectors $\mathbf{1}_{s+1}+\mathbf{u}_i$ and $\mathbf{1}_{s+1}-\mathbf{u}_i+\mathbf{u}_{s+1}$ respectively. By Lemma~\ref{thm:r-s=2}, we obtain that the remaining graph of $G_1\cup G_2$ together with the edge in $B\cap V(K^4)$ contains a perfect matching, as desired. 

 % by a similar discussion as Step 1, there exists an $(\alpha^{1/5},A_i)$-good vertex $v_i\in A_i$ for each $i\in [s-1]$ such that $G[\{v_1,\ldots,v_{s-1},v_s,v_s'\}]$ forms a clique. Furthermore, there exists a  vertex $w\in V(G_j)$ for some $j\in [2]$ such that $w\in N(\{v_1,\ldots,v_{s-1},w_1,w_2\})$, thus forms a clique $K_r:=K^3$. This together with $K^{2-j}$ forms two vertex disjoint copies of $K_r$ with index vectors $\mathbf{1}_{s+1}+\mathbf{u}_s$ and $\mathbf{1}_{s+1}-\mathbf{u}_s+\mathbf{u}_{s+1}$ respectively.  %Moreover, there exists a clique $K_3$ in $G_{2-j}$, say $w_1w_2w_3$. Similarly, there exists an $(\alpha^{1/5},A_i)$-good vertex $v_i\in A_i$ for each $i\in [s]$ such that $G[\{v_1,\ldots,v_{s},w_1,w_2,w_3\}]$ forms a clique $K_r:=K^2$. 
%It is routine to check that $2|V(K^3\cup K^{2-j})\cap A_i|=|V(K^3\cup K^{2-j})\cap B|$, and $|G_1-V(K^1\cup K^{2-j})|$ and $|G_2-V(K^1\cup K^{2-j})|$ are even. Together with Lemma~\ref{thm:r-s=2}, we obtain that $G_1\cup G_2-V(K^1\cup K^{2-j})$ contains a perfect matching, as desired. 


\medskip
{\bf Case 2. $|A_s|<n$.}
\medskip

{Recall that $|A_1|\le \cdots\le |A_s|$. In this case,  $|B|>(r-s)n$. Observe that the $K_r$-tiling  $\mathcal{K}$ has index vector $|\mathcal{K}|(1,\ldots,1,r-s)$. Consequently, in Step $2.1$, either $q>0$  or $|B'|>(r-s)\frac{|G'|}{r}$ holds, which means \(\mathcal{H} \cup \mathcal{F} \ne \emptyset\).}

{Recall that $\mathcal{H}$ and $\mathcal{F}$ are families of vertex-disjoint copies of $K_r$ and $K_3$ obtained in Step~$2.1$, respectively.}
Choose a triangle $K\in \mathcal{F}$ (if $\mathcal{F}\neq \emptyset$, otherwise, choose from a {copy of} $K_r$ in $\mathcal{H}$) with vertex set $\{w_1,w_2,w_3\}\<B$. By~\eqref{mini-degree-K}, we know that $w_1$ is adjacent to some vertex in $V(G_1\cup G_2)$. Without loss of generality, assume that $N(w_1)\cap V(G_1)\neq \emptyset$.  If $K\in\mathcal{F}$, then update $\mathcal{F}$ by replacing the triangle $K$ with $T_2$ in $G_{2}$. By Lemma~\ref{thm:r-s=2}, the updated subgraph induced by  {$ (V(G_1\cup G_2)\setminus   V(T_2))\cup \{w_1,w_2,w_3\}$} has a perfect matching, as desired. If $K$ is chosen from some $K_r$ in $\mathcal{H}$, denoted as $K^2$, then update $\mathcal{H}$ by replacing $K^2$ with $K^{2,j}$ for some $j$ such that they have the same index vector {with respect to $\mathcal{P}'$}. By Lemma~\ref{thm:r-s=2}, we obtain that the current subgraph induced by  $(V(G_1\cup G_2)\setminus V(K^{2,j}))\cup \{w_1,w_2,w_3\}$ contains a perfect matching, as desired. \medskip

It is easy to see that the updated tiling  $\mathcal{K}$ has index vector $|\mathcal{K}|(1,\ldots,1,r-s)$, while the index vectors of updated $\mathcal{H}$ and $\mathcal{F}$ stay unchanged. Moreover,  $\mathcal{K}\cup \mathcal{H}\cup \mathcal{R}$ uses at most one edge in $G'[U\setminus U']$. In all possible cases, \(G'[B'']\) contains a $K_{r-s}$-factor, which is denoted as $\{K^1,\ldots,K^{t}\}$ for $t:=\frac{|B''|}{r-s}$. Let  $\mathcal{F}=\{K^{t+1},\ldots,K^{t
+(r-s)b}\}$ (could be empty if $b=0$).

\medskip
\noindent {\bf Step 2.3. Contracting and then verifying~\ref{A1}-\ref{A2}.} 


 % there exist $t:=\frac{|B'|}{r-s}$ vertex-disjoint copies of $K_{r-s}$, say $K^1,\ldots,K^{t}$, in $G'[B']$, and $q$ vertex-disjoint copies of $K_{r-s+1}$, say $K^{t+1},\ldots,K^{t+q}$, in $H$. 
Contracting each $K^i$ ($i\in [t]$) into a vertex $w_i$ and each $K^j$ ($j\in [t+1,t+(r-s)b]$) into an edge $w_jw_j'$, yields a new vertex set $B^*$ of size $t+2(r-s)b$. 
We then construct an auxiliary graph $G^*$ with $V(G^*)=\big (\bigcup_{i\in[s]}A_i'\big)\cup B^*$, and $xy\in E(G^*)$ if and only if one of the following holds:
\begin{itemize}
    \item $x\in A_i'$, $y\in A_j'$, and $xy\in E(G)$ for distinct $i,j\in[s]$; 
    \item $x=w_i\in B^*$, $y\in A_j'$, and $y\in N_G(V(K^i))$ for $j\in [s]$ and $i\in[t]$;
    \item $x\in \{w_j,w_j'\}\subseteq B^*$, $y\in A_i'$, and $y\in N_G(V(K^j))$ for $i\in [s]$ and $j\in[{t+1,t+(r-s)b}]$;
    \item $\{x,y\}=\{w_j,w_j'\}$ for some $j\in[t+1,t+(r-s)b].$
\end{itemize}
Note that $\mathcal{P}^*=\{A_1',\ldots,A_s',B^*\}$ is a partition of $G^*$.  We first estimate the sizes of $|G^*|$ and $|B^*|$ to verify~\ref{A1}. 
\medskip

% It is routine to check that 
\textbf{(i)~Verify $|B^*|$.}\medskip

%If $b<0$, then $\mathcal{F}=\emptyset$ and \begin{align}\label{G^*}
%|B^*|=\frac{|B'|}{r-s}<\frac{|G'|}{r},~~
%|G^*|=\sum_{i\in [s]}|A_i'|+|B^*|=|G'|-|B'|+|B^*|=|G'|-\frac{r-s-1}{r-s}|B'|.
%\end{align}
%It follows from~\eqref{G^*} that
%$$
%\frac{|G^*|}{s+1}-\frac{|G'|}{r}=\Big(\frac{1}{s+1}-\frac{1}{r}\Big){|G'|}-\frac{r-s-1}{(r-s)(s+1)}|B'|=\frac{r-s-1}{s+1}\Big(\frac{|G'|}{r}-\frac{|B'|}{r-s}\Big)=-\frac{(r-s-1)b}{s+1}\geq 0,
%$$
%which implies that $\frac{|G^*|}{s+1}\geq \frac{|G'|}{r}$. Thus, $|B^*|\leq \frac{|G^*|}{s+1}$. \medskip

%For the case $b\ge 0$, we have
Notice that 
\begin{align*}%\label{B^*}
|B^*|=\frac{|B''|}{r-s}+2(r-s)b
&=\frac{|B'|-(r-s+1)(r-s)b}{r-s}+2(r-s)b \nonumber\\ 
&=\frac{|B'|}{r-s}+(r-s-1)b\nonumber \\ 
&=\frac{|G'|}{r}+(r-s)b.
\end{align*}
%Note that $|B^*|=\frac{|B''|}{r-s}+2(r-s)b$ and $|B''|=|B'|-(r-s+1)(r-s)b$. Together with $b=\frac{|B'|}{r-s}- \frac{|G'|}{r}$, we have
%\begin{align*}
%|B^*|=\frac{|B'|-(r-s)b\cdot(r-s+1)}{r-s}+2(r-s)b=\frac{|B'|}{r-s}+(r-s-1)b=\frac{|G'|}{r}+(r-s)b.
%\end{align*}
It follows that 
\begin{align}\label{b>0}
|G^*|&=\sum_{i=1}^s|A_i'|+|B^*|=\sum_{i=1}^s|A_i'|+\frac{|G'|}{r}+(r-s)\cdot\Big(\frac{|B'|}{r-s}- \frac{|G'|}{r}\Big) \nonumber\\
&=\sum_{i=1}^s|A_i'|+|B'|+\frac{1-r+s}{r}|G'|=|G'|+\frac{1-r+s}{r}|G'|=(s+1)\frac{|G'|}{r}.
\end{align}
 %If $b<0$, then there is no $F$. Together with $b=\frac{|B'|}{r-s}- \frac{|G'|}{r}<0$, we have $|B^*|=\frac{|B'|}{r-s}<\frac{|G'|}{r}$.
%Note that $|G^*|=\sum_{i=1}^s|A_i'|+\frac{|B'|}{r-s}$ and $|G'|=\sum_{i=1}^s|A_i'|+|B'|$.
%This implies $|G^*|=|G'|-\frac{r-s-1}{r-s}|B'|$.
%, where $n^*=|G^*|<\frac{2(s+1)n}{r}.$    Therefore, $|B^*|\leq \frac{|G^*|}{s+1}$. Furthermore,
%$$
%|B^*|-\frac{|G^*|}{s+1}=\frac{|B'|}{r-s}-\frac{|G^*|}{s+1}=\frac{|G'|-|G^*|}{r-s-1}-\frac{|G^*|}{s+1}=\frac{1}{r(r-s-1)}\Big(\frac{|G'|}{r}-\frac{|G^*|}{s+1}\Big)\geq-\frac{2r\alpha n'}{r-s-1}.
%$$
As $b\leq r\alpha n$ and $\frac{|G|}{r}\leq 2\frac{|G^*|}{s+1}$, \eqref{b>0} implies that 
\begin{align}\label{equal}
\frac{|G^*|}{s+1}= \frac{|G'|}{r}\ \text{and}\  |B^*|-\frac{|G^*|}{s+1}=(r-s)b\leq r^2\alpha n\leq 2r^2\alpha \frac{|G^*|}{s+1}.
\end{align}
Hence, $B^*$ satisfies~\ref{A1} {with parameter $2r^2\alpha$ in place of $\alpha$}.\medskip
%-2r^3\alpha n'\leq |B^*|-\frac{|G^*|}{s+1}=(r-s)b\leq 2r^3\alpha n'$. 

%By the construction of $\mathcal{K}\cup \mathcal{H}$, one has $\frac{|G'|}{r}=\frac{|G|-|V(\mathcal{K}\cup \mathcal{H})|}{r}=n-q-\frac{|V(\mathcal{K})|}{r}$. Moreover, the following hold:
\textbf{(ii)~Verify $|A_i'|$.}\medskip

%By the construction of $\mathcal{K}\cup \mathcal{H}$, one has $$\frac{|G^*|}{s+1}\geq \frac{|G'|}{r}=\frac{|G|-|V(\mathcal{K}\cup \mathcal{H})|}{r}=n-q-\frac{|V(\mathcal{K})|}{r}.$$ Therefore, 
By \eqref{G'} and \eqref{equal}, for each $i\in [p]$ we have 
\begin{align*}%\label{eq:A1}
|A_i'|-\frac{|G^*|}{s+1}\leq \frac{|G'|}{r}- \frac{|G^*|}{s+1}=0,%|A_i|-\frac{|V(\mathcal{K}\cup \mathcal{H})|}{r}-\frac{|G'|}{r}=,
\end{align*}
for each $i\in [p+1,s]$  we have 
\begin{align}\notag
|A_i'|-\frac{|G^*|}{s+1}&=|A_i'|-\frac{|G'|}{r}\leq |A_i|-\frac{|V(\mathcal{K})|}{r}-\frac{|V(\mathcal{R})|}{r}-\frac{|G|-|V(\mathcal{K}\cup \mathcal{H}\cup \mathcal{R})|}{r}\\\notag
&=|A_i|-n+q\leq 2\alpha n\leq 3\alpha  \frac{|G^*|}{s+1}. 
\end{align}

In a summary, $G^*$ satisfies $|A_1'|,\ldots, |A_{s}'|,|B^*|\leq \frac{|G^*|}{s+1}+2r^2\alpha \frac{|G^*|}{s+1}$ as desired in~\ref{A1} under the partition $\mathcal{P}^*$.  %{Tell what you are doing next}

%Observe that $|V(\mathcal{K}\cup\mathcal{H})\cap B|\geq (r-s)|V(\mathcal{K}\cup\mathcal{H})\cap A_i|$ for each $i\in [s]$, and $|V(\mathcal{K}\cup\mathcal{H})\cap A_i|\geq |V(\mathcal{K}\cup\mathcal{H})\cap A_j|$ for distinct $i,j\in [s]$ with $|A_i|\leq |A_j|$. 

\medskip


%\item if $|A_s|=n$, then for each $i\in [s]$ we have 
%$$
%|A_i'|-\frac{|G^*|}{s+1}=|A_i|-\frac{|V(\mathcal{K})|}{r}-(q-(n-|A_i|))-\frac{|G|-|V(\mathcal{K}\cup \mathcal{H})|}{r}=0.
%$$
%\end{itemize}

\medskip
\textbf{(iii)~Verify~\ref{A3} and~\ref{A2}.}\medskip

Notice that after covering all bad or exceptional vertices in Step 1, each $D'\in \mathcal{P}'$ satisfies $\De(G'[D'])\le \max\{\be' n, \alpha^{1/5}n\}$. 
% For each $D'\in \mathcal{P}'$ and each $v\notin D'$, as $\alpha\ll\beta'$, we have 
% \begin{align*}%\label{DofQ'}
% d_{G'}(v,D')\ge |D'|-\beta' n -3r^3\alpha n\ge |D'|-2{\beta'} n.
% \end{align*} 
{As $\alpha\ll\beta'$}, for each vertex $v\in A_i'$ with $i\in [s]$ and each $D^*\in \mathcal{P}^*\setminus \{A_i'\}$ we have 
$$
d_{G^*}(v,D^*)\geq |D^*|-\beta' n-|V(\mathcal{K}\cup \mathcal{H}\cup \mathcal{R})|\geq \big(1-2\beta'\big)|D^*|; 
$$
for each vertex $v\in B^*$ and each $i\in [s]$ we have 
$$
d_{G^*}(v,A_i')\geq |A_i'|-(r-s)\beta' n-|V(\mathcal{K}\cup \mathcal{H}\cup \mathcal{R})|\geq \big(1-2(r-s)\beta'\big)|A_i'|. 
$$
Therefore,~\ref{A3} holds with parameter $2(r-s)\beta'$ in place of $\beta$. 

%for each $D^*\in \mathcal{P}^*$ and each $v\notin D^*$, 
% we have
%\begin{align}\label{D^*-degree}
%d_{G^*}(v,D^*)\geq |D^*|-{2(r-s)\beta' n~(more~details)}\geq \big(1-3(r-s)\beta'\big)|D^*|. 
%\end{align}


It follows from the discussion in \textbf{(ii)} that $|A_i'|>\frac{|G^*|}{s+1}$ only if $|A_i|>n$. 
%If  $|A_s|\leq n$, then $|A_i'|>\frac{|G^*|}{s+1}$ only if $|A_i|>n$ or $|A_i|=n$ and $i\in [t]$. Let $p\in [p]$ be the largest integer such that $|A_i|<n$. Then $q=\sum_{i\in [p]}(n-|A_i|)$. 
%Hence if $A_i$ is a small part in $\mathcal{P}$, then we have $|A_i'|\leq \frac{|G^*|}{s+1}$, that is, $A_i'$ is a small part in $\mathcal{P}'$. 
%Moreover, since $|B|\leq (r-s)n$ if $r-s\geq 2$, one has $|B^*|\leq \frac{|G^*|}{s+1}$, that is, $B^*$ is a small part of $\mathcal{P^*}$ if $r-s\geq 2$.
%{Clearly, in constructing \(\mathcal{K} \cup \mathcal{K}'\), condition \ref{B2} allows us to avoid using vertices from at least \(|A_i| - n\) matching edges in each \(G[A_i]\) with \(|A_i| > n\), \(i \in [s]\).}
Therefore,  
\begin{itemize}    
\item\label{C2} for each $A_i'$ with $|A_i'|>\frac{|G^*|}{s+1}$, $G[A_i']$ contains a matching of size at least $|A_i'|-\frac{|G'|}{r}= |A_i'|-\frac{|G^*|}{s+1}$, {by combining \eqref{matching'}, \eqref{equal} and the fact that $\mathcal{K}\cup \mathcal{H}$ uses at most one edge in $E(G'[U\setminus U'])$};
    \item $|B^*|= \frac{|G^*|}{s+1}+(r-s)b$ and $G^*[B^*]$ contains a matching of size $(r-s)b$. % if $r-s=1$, then there are at least $r\alpha n$ matching edges inside $G'[B']$ as $|B'|\geq n-100r^2\alpha n$ and $G'[B']$ contains no $\gamma'$-independent $n$-set. 
\end{itemize}
Thus \ref{A2} holds. Notice that $|B^*|\geq  \frac{|G^*|}{s+1}$. By Lemma~\ref{lem:balance}, $G^*$ has a  $K_{s+1}$-factor in which each copy of $K_{s+1}$ satisfies either $V(K_{s+1})\cap B^*=\{w_i\}$ for some $i\in [t]$,  or $V(K_{s+1})\cap B^*=\{w_j,w_j'\}$ for some $j\in [t+1,t+(r-s)b]$. Thus, such a $K_{s+1}$-factor in $G^*$ corresponds to a $K_r$-factor in $G'$. 
%{which yields a $K_r$-factor in $G'$.(??seems we should require that all edges in $B^*$ are covered by disjoint copies of $k_{s+1}$)} 
Together with $\mathcal{K}\cup \mathcal{H}\cup \mathcal{R}$, we obtain a $K_r$-factor in $G$, as desired.   % and we complete the proof.
\end{proof}
%\medskip
%{\bf Remark.} As the proof demonstrates of Lemma~\ref{good-partition}, condition~\ref{B2} is overly strong. In fact, Lemma~\ref{lem:balance} implies that the matching of size $|A_i|-n$ in each large part $A_i$ is sufficient to find a $K_r$-factor. Furthermore, condition~\ref{B6} can be relaxed to only require the presence of at least one edge between any two sets of size $n$ in $B$, which is enough to guide us in finding a perfect matching in $B$ when $r-s=2$.

\subsection{Proof of Lemma~\ref{vertex-cover}}
In this subsection, we prove Lemma~\ref{vertex-cover}, which establishes the existence of a subgraph with a large minimum vertex cover.
\begin{proof}[\bf Proof of Lemma~\ref{vertex-cover}]
Let $G$ be an $((r-1)n+1)$-regular graph on $rn$ vertices, and let $\{A_1,A_2,\ldots,A_r\}$ be a balanced partition of $V(G)$. Suppose that $A_i^1$ is a minimum vertex cover of $G[A_i]$ for each $i\in [r]$. Without loss of generality, assume that $|A_1^1|\geq |A_2^1|\geq \cdots \geq |A_r^1|$. Clearly, there exists a vertex cover $C_i$ of $G[A_i]$ with size $|A_2^1|$ for each $i\in [2,r]$. For convenience, let $C_1:=A_1^1$. Denote $x_i:=|C_i|$ and $B_i:=A_i\setminus C_i$ for each $i\in [r]$. Hence, $G[B_i]=\emptyset$ for each $i\in [r]$. We proceed by considering the following two possible cases. 
% ~{(I suggest to update $C_i$ instead of $A_i^1$, and use $B_i$ for $A_i\setminus C_i$)}

\medskip
{\bf Case 1. $e\Big(B_1,\bigcup_{j\neq 2}C_j\Big)\leq e\Big(C_2,\bigcup_{j\neq 1}B_j\Big)$.}
\medskip

Since $G$ is  $((r-1)n+1)$-regular, one has 
$$
e(C_2,B_1)\leq |C_2|((r-1)n+1)-e\Big(C_2,\bigcup_{j\neq 1}B_j\Big).
$$
Note that as $x_2=x_3=\cdots=x_r$,
$$
e\Big(B_1,\bigcup_{j\neq 1}B_j\Big)\le\sum_{j\neq1}(n-x_1)(n-x_j)=(n-x_1)(r-1)(n-x_2).
$$
Therefore, we have
\begin{align*}
e\Big(B_1,\bigcup_{j\neq 2}C_j\Big)&=|B_1|((r-1)n+1)-e(B_1,C_2)-e\Big(B_1,\bigcup_{j\neq 1}B_j\Big)\\
&\geq (n-x_1)((r-1)n+1)-x_2((r-1)n+1)+e\Big(C_2,\bigcup_{j\neq 1}B_j\Big)-e\Big(B_1,\bigcup_{j\neq 1}B_j\Big)\\
&\geq (n-x_1-x_2)((r-1)n+1)+e\Big(C_2,\bigcup_{j\neq 1}B_j\Big)-(n-x_1)(r-1)(n-x_2).
\end{align*}
%where the last equality holds as $x_j=x_2$ for each $j\in [2,r]$.
By our assumption in this case, one has
\begin{align*}
(n-x_1-x_2)((r-1)n+1)\leq (n-x_1)(r-1)(n-x_2).
\end{align*}
It follows that $n\leq (r-1)x_1x_2+(x_1+x_2)\leq (r-1)x_1^2+2x_1.$ Thus, $x_1\geq \frac{1}{r}\sqrt{n}$, as desired.

\medskip
{\bf Case 2. $e\Big(B_1,\bigcup_{j\neq 2}C_j\Big)> e\Big(C_2,\bigcup_{j\neq 1}B_j\Big)$.}
\medskip

Notice that
$$
e\Big(\bigcup_{j\neq 2}C_j,\bigcup_{j\neq 1}B_j\Big)\leq \Big|\bigcup_{j\neq 2}C_j\Big|((r-1)n+1)-e\Big(\bigcup_{j\neq 2}C_j,B_1\Big).
$$
Therefore,\allowdisplaybreaks
\begin{align*}
e\Big(\bigcup_{j\neq 1}B_j,C_2\Big)=&\Big|\bigcup_{j\neq 1}B_j\Big|((r-1)n+1)-e\Big(\bigcup_{j\neq 1}B_j,\bigcup_{j\neq 2}C_j\Big)-e\Big(\bigcup_{j\neq 1}B_j,\bigcup_{t\in [r]}B_t\Big)\\
\geq &\Big|\bigcup_{j\neq 1}B_j\Big|((r-1)n+1)-\Big|\bigcup_{j\neq 2}C_j\Big|((r-1)n+1)+e\Big(\bigcup_{j\neq 2}C_j,B_1\Big)-\sum_{j\neq 1}\sum_{t\neq j}e(B_j,B_t)\\
\geq &((r-1)n+1)\Big((r-1)n-\sum_{j\neq 1}x_j-\sum_{j\neq 2}x_j\Big)+e\Big(\bigcup_{j\neq 2}C_j,B_1\Big)-\sum_{j\neq 1}\sum_{t\neq j}(n-x_j)(n-x_t)\\
=&((r-1)n+1)\Big((r-1)n-(r-1)x_2-x_1-(r-2)x_2\Big)+e\Big(\bigcup_{j\neq 2}C_j,B_1\Big)\\
&-(r-1)(n-x_2)\Big((r-1)n-x_1-(r-2)x_2\Big)\\
=&(r-1)n-x_1-(2r-3)x_2-(r-1)x_1x_2-(r-1)(r-2)x_2^2+e\Big(\bigcup_{j\neq 2}C_j,B_1\Big).
\end{align*}
By the assumption in this case, one has
$$
(r-1)n\leq x_1+(2r-3)x_2+(r-1)x_1x_2+(r-1)(r-2)x_2^2\leq (r-1)\big((r-1)x_1^2+2x_1\big).
$$
Thus, $x_1\geq \frac{1}{r}\sqrt{n}$, as desired. 
\end{proof}





%%%%%%%%%%%%%%%%%%%%%%%%%%%%%%%%%%%%%%%


\section{Concluding remarks}
%{Slightly discuss $G[p]$ for general $0<p<1$?}
%{A natural extension involves considering more general parameters for our problem. More precisely, let $G[p]$ be a random induced subgraph including each vertex independently with probability $p$ where $p\in (0,1)$ and we may consider a $d$-regular graph on $rn$ vertices and ask whether $G[p]$ admits a $K_r$-factor with a constant probability. On the one hand, we also consider an $((r-1)n+1)$-regular graph on $rn$ vertices and decrease $p$. Note that the probability that $G[p]$ admits a $K_r$-factor may be likely to asymptotically approach $0$ as $p\to 0$. Therefore, it is nature to ask for the point of $p$ where the behavior changes. On the other hand, we stick to $p=\tfrac{1}{2}$ and decrease $d$. In this case, we need to impose additional conditions to ensure that the random induced subgraph remains dense.}

{Dragani\'{c}, Keevash and M\"{u}yesser \cite{Keevash2025} proposed a plausible class of extremal construction for Conjecture~\ref{conj:K_r factor}: slightly unbalanced complete $r$-partite graphs with a suitable  factor in the largest part $A$.  They conjectured that the optimal constant $c$ should be  $\frac{1}{r^2}$, where in the random induced subgraph $G[S]$ we have one
probability factor of $\frac{1}{r}$ for $r$ dividing $|S|$ and another for $A\cap S$ being the largest part. In what follows, we present a construction to demonstrate that these two events alone are insufficient to guarantee the existence of a  $K_r$-factor in $G[S]$.}

{Assume that $r\geq 3$. We first consider an $rn$-vertex complete $r$-partite graph $G_0$ with vertex partition $A_1\cup \ldots \cup A_r$,  where $|A_1|=n+r-1$ and $|A_2|=\cdots=|A_r|=n-1$. Let $G$ be a graph obtained from $G_0$ by embedding an $r$-regular graph into  $A_1$ {such that $G[A_1]$ is triangle-free}. Clearly, $G$ is  $((r-1)n+1)$-regular and $G[A_i]$ is empty for every $i\in [2,r]$. Consider an arbitrary vertex subset $S\subseteq V(G)$. A necessary condition for $G[S]$ to contain a $K_r$-factor is that 
\begin{itemize}
    \item $r\mid |S|$;
    \item $\frac{|S|}{r}\le |A_1 \cap S|\le \frac{2|S|}{r}$;
   \item $|A_i \cap S|\le \frac{|S|}{r}$ for all $i \in [2,r]$. 
\end{itemize}
This implies that the optimal constant 
$c$ in Conjecture~\ref{conj:K_r factor} must be strictly less than $\frac{1}{r^2}$ {for every $r\ge 3$}. }

{In Theorem~\ref{thm:main thm} we assume that $G$ is $((r-1)n+1)$-regular, and we have remarked that the regularity requirement cannot be relaxed to a mere minimum degree condition. The need for regularity, however, arises only in Lemma~\ref{vertex-cover}. In fact, by carefully examining the proof of Lemma~\ref{vertex-cover}, we can replace the regularity assumption by a weaker degree condition: 
\begin{align}\label{mini-degree-G}
(r-1)n+1\leq \de(G)\leq \De(G)\leq (r-1)n+n^{0.5}.
\end{align}
Therefore, we can slightly extend Theorem~\ref{thm:main thm} under~\eqref{mini-degree-G}. Furthermore, one may also consider the more general setting of a $d$-regular graph where $d\ge(r-1)n+1$, and ask for determining the probability that its random induced subgraph $G[p]$ contains a $K_r$-factor, where $p\in(0,1)$.}

\vspace{0.5cm}
\textbf{Acknowledgement.}
The third author would like to thank Professor Jie Han for his valuable discussions and suggestions.
%{From here I wonder if we can prove Thm4.1 for almost regular graphs so as to extend our main result (Yang: after check the proof, here we can replace $n^{0.6}$ with $\alpha n$)} 

\bibliographystyle{abbrv}
\bibliography{ref.bib}
\end{document}


\section{Extremal case}
In this section, we study the extremal case. 
Let $G$ be an $((r-1)n+1)$-regular graph on $rn$ vertices and assume that $G$ has a $\gamma$-independent set of size $n$.
In order to obtain a clear picture of $G$, our first step is to find as many vertex-disjoint sparse $n$-sets as possible.
Let $s$ be the maximum number of sparse sets with $1\le s\le r$ and $A_1^0,A_2^0\cdots,A_s^0$ are pairwise disjoint sparse sets such that $B^0:=G\setminus \bigcup_{j=1}^{s}A_j^0$ has no sparse $n$-set.
Let $\mathcal{P}_0=\{A_1^0,\cdots,A_s^0,B^0\}$ be an obtained partition of $V(G)$ which can only capture the rough structure of $G$.
Motivated by the Hajnal-Szemer\'edi theorem in multipartite graph (see Lemma~\ref{lem:HS thm in multipartite graph}), we adjust this partition slightly by removing a small number of vertices.
This yields a new partition $\mathcal{P}=\{A_1,\cdots,A_s,B\}$ such that each part $A_i$ has bounded maximum degree, while every vertex in $A_i$ has large degree to the outside of $A_i$.
We randomly sample vertices from each part and obtain a partition $\mathcal{P}'=\{A_1',\cdots,A_s',B'\}$, one that with high probability inherits the degree distributions of $\mathcal{P}_1$.
However, $\{A_1',\cdots,A_s',B'\}$ may be unbalanced.
Lemma~\ref{lem:balance} implies that using the matching in each part one can remove a $K_r$-tiling of small size to make the partition balance again.

We begin this section by introducing a powerful tool for handling the extremal case.%极值情况的turan定理
\begin{theorem}

\end{theorem}


The following multi-partite version of Hajnal-Szemer\'edi theorem provides a degree condition that guarantees the existence of $K_r$-factors in a balanced $r$-partite graph, which follows from \cite[Lemma 6.3]{gan2024} by using Edmonds algorithm.

\begin{lemma}[\cite{gan2024}]\label{lem:HS thm in multipartite graph} 
For any $r,n\in \mathbb{N}$, let $G$ be an $r$-partite graph on vertex classes $V_1,\ldots,V_r$ where $|V_i|=n$ for each $1\le i\le r$. If  
$$
\min_{i\in [r]}\min_{v\in V_i}\min_{j\in [r]\setminus \{i\}} |N(v)\cap V_j|\geq \Big(1-\frac{1}{2r}\Big)n,
$$
%$$d(x,A_j) \ge (1-\frac{1}{r})n+1$$ where $x\in A_i$ and $i\ne j$, 
then $G$ has a  $K_r$-factor.
\end{lemma}

Note that $|A_1^0|=\cdots=|A_s^0|=n$ and $|B_0|=(r-s)n$.
In order to obtain a balance $(s+1)$-partite graph (may be minor adjustments based on the current partition $\mathcal{P}_0$), a natural idea is to contract each copy of $K_{r-s}$ in $B_0$ into a single vertex.
In other words, we need to find a $K_{r-s}$-factor in $B_0$ provided the given conditions of degree and independent set.
When $r-s\ge 3$, it follows from Theorem~\ref{thm:Non-extremal case r>3}.
When $r-s=2$, we use the following result of Treglown~\cite{} to obtain a perfect matching in $B_0$.

\begin{theorem}[\cite{}]\label{thm:r-s=2}
Let $\beta>0$ and $n\in \mathbb{N}$ be even such that $1/n\ll\beta$. Suppose that $G$ is a graph on $2n$ vertices with $$\de(G)\ge (1-2\beta)n.$$ Then at least one of the following conditions holds:
\begin{itemize}
    \item $G$ contains a $3\beta$-independent set of size at least $n$;
    \item $G$ is $3\beta$-closed to $2K_n$, {\color{red}where $n$ is odd};%Tre文章里面没有n是奇数的条件，参考我们均匀染色那篇
    \item $G$ contains a perfect matching.
\end{itemize}
\end{theorem}

\begin{lemma}\label{SuperT}
    Let $0< {1}/{n} \ll \eps \ll \alpha \ll {1}/{r}$ where $r\in \mathbb{N}$ and $r\ge 2$. Suppose that $G$ is an $n$-vertex graph that contains at most $\eps n^r$ copies of $K_r$ and 
$$\delta(G)\ge \Big(1-\frac{1}{r-1}-\alpha\Big)n.$$
%and so that $G$ contains at most $\eps n^r$ copies of $K_r$. 
Then $G$ admits a $\sqrt{\alpha}$-independent set of size at least $\frac{n}{r-1}$.
\end{lemma}

In the following sections, we are ready to construct a partition of independent sets firstly (in section 4.1) and then transform it into a balanced one through a series of appropriate operations (in section 4.2).

\subsection{Constructing sparse sets and adjustment}
\begin{definition}\label{def:bad vertex for Ai}
Let $G$ be a graph and $\mathcal{P}=\{A_1,\ldots,A_k,B\}$ be a partition of $V(G)$. For any $i\in [k]$, we say that 
\begin{itemize}
    \item a vertex $v\in V(G)\setminus A_i$ is $(\beta, A_i)$-bad (or simply $(\beta, A_i)$-bad) if $d(v,A_i)\le \beta n$;
    \item a vertex $v\in A_i$ is $(\beta, A_i)$-good (or simply $(\beta, A_i)$-good) if $d(v,A_i)\le \beta n$;
    \item a vertex $v\in B$ is $(\beta,B)$-exceptional if $d(v,A_i)\le |A_i|-\beta n$ for some $i\in [k]$.
\end{itemize} 
\end{definition}

\begin{definition}[Good partition]
    Let $r,s,n\in \mathbb{N}$, and  $0<\frac{1}{n}\ll \alpha\ll \beta\ll \gamma'\ll \frac{1}{r}$. Suppose that $G$ is a graph on $rn$ vertices, and   $\mathcal{P}=\{A_1,\ldots,A_s,B\}$ is a partition of $V(G)$. We say that $\mathcal{P}$ is a \text{$(\alpha,\beta,\gamma')$-good partition} of $G$ if all of the following hold.
    \begin{enumerate}
[label =\rm  (B\arabic{enumi})]
    \item\label{B1} $n-r\alpha n\leq |A_i|\leq n+r\alpha n$ for each $i\in [s]$, $(r-s)n-r\alpha n\leq |B|\leq (r-s)n$ if $s<r-1$ and $n-r\alpha n\leq |B|\leq n+r\alpha n$ if $r-1\le s\le r$;
    \item\label{B2} $\Delta(G[A_i])\leq \beta n$,  for each $i\in[s]$ and the matching number of $G[A_i]$ is at least $\frac{|A_i|-n}{4\beta}$ if $|A_i|>n$, and the number of $(\alpha^{1/5},A_i)$-good vertices is at least $|A_i|-2\alpha n$ if $|A_i|\leq n$; 
    \item\label{B04} $d(v,A_i)\geq \beta n$ for each vertex $v\in V(G)\setminus A_i$;
    \item\label{B05} $d(v,B)\geq \beta n$ for each vertex $v\in V(G)\setminus B$ and $\de(G[B])\ge (r-1)n+1-sn-r\alpha n$;
    \item\label{B07}  the number of $(\beta,B)$-exceptional vertices is at most $2r\alpha n$;
\item\label{B07} for any two sets of size $n$ in $B$, there are at least $\beta n$ edges between them; 
    \item\label{B6} $d(v,D)\geq |D|-\alpha^{1/5} n$ for every $(\alpha^{1/5},A_i)$-good vertex $v\in A_i$ with $A_i\neq 
    D$;
    \item\label{B07} $A_i$ is a $\gamma'$-independent set for each $i\in [s]$, and $B$ has no $\gamma'$-independent set of size $n$.
\end{enumerate}
\end{definition}
\begin{lemma}\label{good-partition}
    Let $r,n$ be positive integers with $r\ge 3,\ 1\le s\le r$. Assume that $0<\frac{1}{n}\ll \gamma\ll\alpha\ll\beta\ll\gamma'\ll \frac{1}{r}$. Let $G$ be an $((r-1)n+1)$-regular graph on $rn$ vertices. Suppose that $G$ has a $\gamma$-independent set of size $n$. Then there exists an $(\alpha,\beta,\gamma')$-good partition of $G$. %Then there exists a constant $\gamma'\ll \frac{1}{r}$ and a  partition $\mathcal{P}=\{A_1,\ldots,A_s,B\}$ for some $s\in [r]$ such that
    %\begin{enumerate}
%[label =\rm  (B\arabic{enumi})]
   % \item\label{B1} $n-r\alpha n\leq |A_i|\leq n+r\alpha n$ for each $i\in [s]$, $(r-s)n-r\alpha n\leq |B|\leq (r-s)n$ for $s<r-1$ and $n-r\alpha n\leq |B|\leq n+r\alpha n$ for $r-1\le s\le r$;
   % \item\label{B2} $\Delta(G[A_i])\leq \beta n$ for each $i\in[s]$ and the matching number of $G[A_i]$ is at least $\frac{|A_i|-n}{4\beta}$ if $|A_i|>n$, and the number of $(\alpha^{1/5},A_i)$-good vertices is at least $|A_i|-2\alpha n$ if $|A_i|\leq n$; 
    %\item\label{B04} $d(v,A_i)\geq \beta n$ for each vertex $v\in V(G)\setminus A_i$;
   % \item\label{B05} $d(v,B)\geq \beta n$ for each vertex $v\in V(G)\setminus B$ and $\de(G[B])\ge (r-1)n+1-sn-r\alpha n$;
   % \item\label{B07}  the number of $(\beta,B)$-exceptional vertices is at most $2r\alpha n$;
%\item\label{B07} for any two sets of size $n$ in $B$, there are at least $n$ edges between them; 
   % \item\label{B6} $d(v,D)\geq |D|-\alpha^{1/5} n$ for every $(\alpha^{1/5},A_i)$-good vertex $v\notin D$;
  %  \item\label{B07} $A_i$ is a $\gamma'$-independent set for each $i\in [s]$, and $B$ has no $\gamma'$-independent set of size $n$.
%\end{enumerate}
\end{lemma}
\begin{proof}
    Choose constants $\gamma, \gamma_1,\gamma_2,\dots,\gamma_r, \gamma_{r+1}$ satisfying 
\begin{align}\label{eq:p2}
    0<\frac{1}{n}\ll \gamma \ll \gamma_1\ll \dots \ll \gamma_r\ll\frac{1}{r}.
\end{align}
Since $G$ contains a $\gamma$-independent $n$-set, we may greedily choose vertex-disjoint $\gamma_i$-independent $n$-sets $A_1^0, \dotsc, A_s^0$, until no $\gamma_{s+1}$-independent $n$-set remains. 
Let $B^0 := V(G) \setminus \bigcup_{i \in [s]} A_i^0$ and we obtain a partition $\mathcal{P}_0:=\{A_1^0,\ldots,A_s^0,B^0\}$ of $V(G)$.
Note that $B^0$ has no $\gamma_{s+1}$-independent $n$-set and $|B^0|=(r-s)n$.

We further define 
\begin{align}\label{eq:p3}
    \gamma_s\ll \alpha \ll\beta\ll \gamma_{s+1}.
\end{align}
We next estimate the number of bad vertices under the partition $\mathcal{P}_0$. 
\begin{claim}\label{fact:bad vertices for Ai}
For each $i\in [s]$, we have the following properties:
\begin{enumerate}[label=(\arabic*)]
    \item\label{large-degree vtx} the number of $(\alpha^{1/5},A_i)$-good vertices in $A_i^0$ is at least $|A_i^0|-\alpha n$;
    \item\label{bad vtx2} the number of $(\beta,A_i)$-bad vertices in $V(G)\setminus A_i^0$ is at most $\alpha n$, the number of $(\beta,B)$-bad vertices in $V(G)\setminus B^0$ is at most $r\alpha n$;
    \item\label{bad vtx3} any vertex $v\in V(G)$ cannot be both $(\beta,A_i)$-bad and $(\beta,A_j)$-bad for distinct $i,j\in [s]$;
    \item %\begin{claim}\label{fact:bad vertices in B}
the number of $(\beta,B^0)$-exceptional  vertices in $B^0$ is at most $\alpha^{1/4} n$.
\end{enumerate}
\end{claim}

\begin{proof}[Proof of Claim \ref{fact:bad vertices for Ai}]
(1) Suppose there are at least $\alpha n$ vertices in $A_i^0$ such that $d(v,A_i^0)\ge \beta n$.
As $G$ is $((r-1)n+1)$-regular, for such vertex $v$ in $A_i^0$ we have
\[
d(v,A_i^0)\ge (r-1)n+1-\alpha^{1/5} n-(r-2)n\ge (1-\alpha^{1/5})n.
\]
As $\gamma_i\ll\gamma_s\ll\alpha\ll\beta$ for each $i\in [s]$, we have
\[
e(G[A_i^0])\ge \frac{1}{2}\cdot(1-\alpha^{1/5} n)\cdot \alpha n\ge \gamma_i n^2.
\]
This is a contradiction with the fact that $A_i^0$ is a $\gamma_i$-independent set.

(2) Since $G$ is $((r-1)n+1)$-regular, the number of edges in $G[A_i^0,V(G)\setminus A_i^0]$ is at least
$$\underset{v\in A_i^0}{\sum}d(v)-2e(G[A_i^0])\ge 
|A_i^0| \cdot ((r-1)n+1)-2 e\left(G[A_i^0]\right) \ge |A_i^0||V(G)\setminus A_i^0|-2\gamma_i n^2.
$$
This implies that there are at most $2\gamma_i n^2$ non-edges between $A_i^0$ and $V(G)\setminus A_i$.
Note that a $(\beta,A_i)$-bad vertex in $V(G)\setminus A_i^0$ contributes at least $(1-\beta) n$ non-edges between $A_i^0$ and $V(G)\setminus A_i^0$.
As $\gamma_i \ll \alpha \ll \beta$, for each $i\in [s]$, we have that the number of $(\beta,A_i)$-bad vertices in $V(G)\setminus A_i^0$ is at most
\[
\frac{2\gamma_i n^2}{(1-\beta) n}\le\alpha n.
\]

If $s< r-1$, then $d(v,B^0)\ge |B^0|-n+1\ge n+1$ for every $v\in V(G)$. Therefore, $G$ contains no $(\beta,B^0)$-bad vertex. If $s=r-1$ and $v\in A_i^0$ is a $(\beta,B^0)$-bad vertex, then $v$ cannot be a $(\beta,A_i)$-good vertex. Together with (1), we conclude that the number of $(\beta,B^0)$-bad vertices is at most $r\alpha n.$
% $$|V_{ex}(\mathcal{P}_0,\beta,i)|\leq |V_{ne}(\mathcal{P}_0,\alpha^{1/5},i)|\le \frac{2\gamma_i n^2}{\alpha^{1/5}n}\leq  \alpha n.$$

(3) Let $u\in A_i^0$ be a $(\beta,A_j)$-bad vertex for distinct $i,j\in [s]$. Then $d(u,A_{\ell}^0)\geq n-\beta n$ for all $\ell\in [s]\setminus \{j\}$ due to the regularity of $G$. 

(4) For each $i\in [s]$ with $s\le r-1$, let $B_i$ be the set of bad vertices in $B^0$ satisfying $d(v,A_i^0)\le |A_i^0|-\beta n$ and $B_b=\bigcup_{i\in [s]}B_i$. Now we consider  the size of $B_i$.
Note that $|B^0|=(r-s)n$.
Thus we have that
\[
e(G[A_i^0,B^0])\le |B_i|\cdot (|A_i^0|-\beta n)+((r-s)n-|B_i|)\cdot|A_i^0|=|A_i^0|(r-s)n-|B_i|\beta n.
\]
Recall that $G$ is $(r-1)n+1$-regular and $A_i^0$ is a $\gamma_i$-independent set of $G$. Then % Combing with $d(v)=(r-1)n+1$ for each $v\in A_i$, we obtain that
\[
\gamma_in^2\geq e(G[A_i^0])\ge|A_i^0|\cdot\big((r-1)n+1-(s-1)n\big)-e(G[A_i^0,B^0])\ge |A_i^0|+|B_i|\cdot\beta n=n+|B_i|\cdot\beta n,
\]
which implies $|B_i|\leq \alpha n$. Thus 
%Recall that $A_i^0$ is a $\gamma_i$-independent set of $G$ and at most $r\alpha n$ vertices are moved from other part to $A_i^0$.
%As $\gamma_i\ll\alpha\ll\frac{1}{r}$, we have
%$$
%e(G[A_i])\leq (\gamma_i+2r\alpha) n^2<\sqrt{\alpha}n^2.
%$$
%Thus $|A_i|+|B_i|\cdot\beta n\le e(G[A_i])\le \sqrt{\alpha}n^2$, which implies that
$
|B_b|=\sum_{i\in [s]}|B_i|\le r\alpha n.$ 
%as $\alpha\ll\beta$.
\end{proof}

Now, we move all $(\beta, D)$-bad vertices from $V(G)\setminus D$ to $D$ for each $D\in \mathcal{P}_0$, and denote the resulting partition by $\mathcal{P}_1:=\{A_1^1,\ldots,A_s^1,B^1\}$. Based on Claim \ref{fact:bad vertices for Ai}, we obtain that \ref{B1}, \ref{B04} and \ref{B07} hold. 
%Next, we will proceed with the following two-step operations of moving at most $\alpha n$ vertices. The first step aims to obtain a new partition satisfying the degree condition in Lemma~\ref{lem:balance}. It is obvious that the new partition may be unbalanced. So the purpose of Step 2 is to identify a collection of matchings in the \textit{large} parts, whereby we adjust the partition to be balanced.

%\medskip
%{\bf Step 1. Move all $(\beta, D)$-bad vertices from $V(G)\setminus D$ to $D$ for each $D\in \mathcal{P}_0$.}
%\medskip

%For each $D\in \mathcal{P}$, if $v\in V(G)\setminus D$ is $(\beta,D)$-bad, then we move $v$ to $D$. 
%In the process of moving vertices, there is a tiny difference between the cases $s< r-1$ and the others in whether to move vertices to $B^0$. To be precise, i
%Notice that if $s< r-1$, then $d(v,B^0)\ge |B^0|-n+1\ge n+1$ for every $v\in V(G)$. In this case, we do not need to move any vertices to $B^0$. If $s\ge r-1$, then there may exists $v\in A_i$ such that $d(v,B^0)\le \beta n$ for some $i\in [s]$ (such a vertex cannot be $(\beta,A_i)$-good). In this case, we move such vertex $v$ into $B^0$. Note that the resulting partition after the above moving procedure may not be balanced and we denote it as $\mathcal{P}_1:=\{A_1^1,\ldots,A_s^1,B^1\}$, in which every vertex in $A_i^1$ has large degree to the outside of $A_i^1$. 
For each $i \in [s]$, a part $A_i^1$ is called \textit{large} if $|A_i^1| > n$, and \textit{small} otherwise. Next, we will move vertices in large parts to satisfy \ref{B2}.  
%\medskip
%{\bf Step 2. Move not $(\beta, A_i)$-bad vertices from large part $A_i^1$ to small one $A_j^1$.}
%\medskip
Let $A_i^1$ be a set with size greater than $n$. If $\Delta(G[A_i^1])\leq \beta n$, then we are done. Otherwise, move a vertex in \(A_i^1\) which has at least \( \beta n \) neighbors in \(A_i^1\) to the other side \(A_j^1\) if  $|A_j^1|<n$, or $B^1$ if $|B^1|<(r-s)n$, and update \(\mathcal{P}^1\) accordingly. %Assume that $|A_i^1|>n$ and $\Delta(G[A_i^1])\leq \beta n$ for some $i\in [s]$.  arbitrarily move each vertex in $A_i^1$ which has at least $\beta n$ neighbors in $A_i^1$ to other side $A_j^1$, where $j\in [s]$ with $|A_j^1|<n$, and update $\mathcal{P}^1$ accordingly. 
Repeat this until we reach the point where $|A_i^1|=n$ for each $i\in [s]$ and $|B^1|=(r-s)n$, or $G[A_i^1]$ has maximum degree at most $\beta n$ whenever $|A_i^1|>n$, denote the resulting partition by $\mathcal{P}:=\{A_1,\ldots,A_s,B\}$. 

Recall that $G$ is $(r-1)n+1$ regular. Let $|A_i|=n+k_i$ for each $i\in [s]$. If $k_i>0$, then together with $\Delta(G[A_i])\leq \beta n$, we obtain that the number of matching edges in $G[A_i]$ is at least
$$
\frac{\big((r-1)n+1-(rn-|A_i|)\big)|A_i|}{4\Delta(G[A_i])}\geq \frac{(k_i+1)(n+k_i)}{4\beta n}\geq \frac{k_i}{4\beta}.
$$
It is straightforward to check that \ref{B1}-\ref{B6} holds  and \ref{B07} holds with $\gamma'=\frac{\gamma_{s+1}}{2}$. 
%\medskip
%After the two-step moving operation, we state the following fact that a large matching can be found in a large part.
\end{proof}
%Let $r,n,s$ be positive integers with $r\ge 3,\ 1\le s\le r$ and let $\gamma$ be a positive constant. Define constants $\gamma, \gamma_1,\gamma_2,\dots,\gamma_r, \gamma_{r+1}$ satisfying 
%\begin{align}\label{eq:p2}
%    0<\frac{1}{n}\ll \gamma \ll \gamma_1\ll \dots \ll \gamma_r\ll\gamma_{r+1}=\frac{1}{r}.
%\end{align}

%Since $G$ contains a $\gamma$-independent $n$-set, we may greedily choose vertex-disjoint $\gamma_i$-independent $n$-sets $A_1^0, \dotsc, A_s^0$, until no $\gamma_{s+1}$-independent $n$-set remains. 
%Let $B^0 := V(G) \setminus \bigcup_{i \in [s]} A_i^0$ and we obtain a partition $\mathcal{P}_0:=\{A_1^0,\ldots,A_s^0,B^0\}$ of $V(G)$.
%Note that $B^0$ has no $\gamma_{s+1}$-independent $n$-set and $|B^0|=(r-s)n$.

%When the constant $s$ is fixed, we define 
%\begin{align}\label{eq:p3}
%    \gamma_s\ll \alpha \ll\beta\ll \gamma_{s+1}.
%\end{align}




In the following, we give a sufficient condition that guarantees the existence of $K_r$-factors in $G$.
%Applying with Lemma~\ref{lem:HS thm in multipartite graph}, we have the following result to guarantee the existence of $K_r$-factors in $G$. 


%Without loss of generality, assume that $|A_i|=n+k_i$ for each $i\in [s]$ and $k_1\leq \cdots\leq k_s$. Note that $\sum_{i\in [s]}k_i=0$. We have the following result for unbalanced partition.
\begin{lemma}\label{lem:balance}
Let $0 < \alpha \ll \beta \ll 1$. Suppose $G$ is a graph on $rn$ vertices with a partition $V(G) =\{A_1, \ldots, A_r\}$ satisfying 
\begin{enumerate}
[label =\rm  (A\arabic{enumi})]
    \item\label{A1} $n-\alpha n\leq |A_1|\leq \cdots\leq |A_r|\leq n+\alpha n$;
    \item\label{A2} for each $i \in [r]$, $G[A_i]$ contains a matching of size $k_i:=|A_i|-n$;
    \item\label{A3} for any distinct $i,j \in [r]$ and any $v \in A_i$, we have $d(v, A_j) \ge n - \beta n$.
\end{enumerate}
Then $G$ contains a $K_r$-factor.
\end{lemma}


%We only consider the case that the following operations yield an unbalanced partition, otherwise we can greedily find a $K_r$-factor. The following section details how to adjust the partition satisfying Lemma~\ref{lem:balance} and how to find a $K_r$-factor in an unbalanced partition. For the current $(r,s)$-partition $\mathcal{P}_0$, we define the \textit{bad} vertices motivated by Lemma~\ref{lem:balance} and then move them to appropriate parts.


% we define the following types of vertices. 
% \begin{itemize}
%     \item For a vertex $x\in A_i$ with $i\in [s]$, we say $x$ is $(\beta, A_i)$-bad (or simply {$(\beta, A_i)$-bad}) if $d(x,A_i)\ge \beta n$. %Otherwise, we say $x$ is \textit{$(\beta,A_i)$-good}. 
%     Let $V_b(\mathcal{P},\beta,i)$ be the set of $(\beta,A_i)$-bad vertices.
%     \item For a vertex $x\in V(G)\setminus A_i$ with $i\in [s]$, we say that $x$ is $(\beta, A_i)$-exceptional (or simply {$(\beta, i)$-exceptional}) if $d(x,A_i)\le \beta n$. %Otherwise we say $x$ is \textit{$(\beta,i)$-acceptable}. 
%     Let $V_{ex}(\mathcal{P},\beta,i)$ be the set of $(\beta,i)$-exceptional vertices. 
%     \item For a vertex $x\in V(G)\setminus A_i$ with $i\in [s]$, we say that $x$ is $(\beta, A_i)$-excellent (or simply {$(\beta, i)$-excellent}) if $d(x,A_i)\ge |A_i|-\beta n$. Let $V_e(\mathcal{P}, \beta,i)$ be the set of $(\beta ,i)$-excellent vertices, and let $V_{ne}(\mathcal{P}, \beta,i):=V(G)\setminus (A_i\cup V_e(\mathcal{P}, \beta,i))$. % be the set of vertices which are not $(\beta, i)$-excellent vertices and not belonged to $A_i$. 
    
%     \item A vertex $x\in V(G)\setminus B$ is called $(\beta,B)$-excellent if $d(x,B)\ge |B|-\beta n$.  Denote by $V_e(\mathcal{P},\beta, B)$ the set of all such  $(\beta,B)$-excellent vertices, and let $V_{ne}(\mathcal{P},\beta, B):=V(G)\setminus (B\cup V_e(\mathcal{P}, \beta,B))$.
% \end{itemize}

%\end{definition}

%Since $G$ is $((r-1)n+1)$-regular, we estimate the number of bad vertices under the partition $\mathcal{P}^0$.




% For each $i\in [s]$, since $A_i^0$ is a $\gamma_i$-independent set in $G$ and $\gamma_i\ll \alpha \ll \beta$, we have 
% $$|V_b(\mathcal{P}_0,\beta,i)|\le\frac{2e(G[A_i^0])}{\beta n} \le  \frac{2\gamma_i n^2}{\beta n}\le \alpha n.$$ %Observe that at least $\binom{|A_i^0|}{2}-\gamma_i n^2$ edges are missing in  $G[A_i^0]$. 


%\begin{fact}
%For each $i\in [s]$, if $|A_i|>n$ then there exists a matching of size at least $\frac{k_i}{2\beta}$ where $k_i:=|A_i|-n$.
%\end{fact}

%\begin{proof}
%Note that $\Delta(G[A_i])\leq \beta n$ and $\de(G[A_i])\ge k_i+1$ as $G$ is $((r-1)n+1)$-regular.
%Let $M_i$ be a maximum matching of $A_i$. By Vizing's theorem, we have
%\begin{align*}
%|M_i|\ge \frac{e(G[A_i])}{\Delta(G[A_i])+1}\ge \frac{|A_i|\cdot (k_i+1)}{2(\beta n+1)}\ge \frac{k_i}{2\beta}.
%\end{align*}
%\end{proof}




%Now we are ready to prove Lemma~\ref{lem:balance}, which establishes the existence of $K_r$-factors.

\begin{proof}%[Proof of Lemma~\ref{lem:balance}]
Suppose that $|A_j| = |A_1| + a_j$ for each $j \in [2, r]$, and let $p \in [s]$ be the largest integer such that $|A_p| \le n$. We proceed in two steps:
\begin{itemize}
    \item For each $i \in [2, p]$, use $n - |A_i|$ copies of $K_r$ with index vector $\mathbf{1}_r - \mathbf{u}_i + \mathbf{u}_\ell$ for some $\ell \in [p+1, r]$. Let $k_i'$ ($\le k_i$) be the number of matching edges used in this step.

    \item  For each $i \in [p+1, r]$, use $k_i - k_i'$ copies of $K_r$ with index vector $\mathbf{1}_r - \mathbf{u}_1 + \mathbf{u}_i$.
\end{itemize}

The degree condition allows a greedy vertex-selection procedure for both steps, which yields a $K_r$-tiling denoted by $\mathcal{K}$. We now show that the resulting partition is balanced.
\begin{claim}\label{result-balance}
    For each $\ell\in [r]$, we have $|A_\ell\setminus V(\mathcal{K})|=a_1 + \sum_{i=2}^p k_i$.
\end{claim}



\begin{proof}[Proof of Claim \ref{result-balance}]
Notice that $\sum_{i\in [r]}k_i=0$ and $\sum_{i=p+1}^{r}k_i'=\sum_{i=2}^p(n-|A_i|)$. Hence $-\sum_{i\in [p]}k_i=\sum_{i\in [p+1,r]}k_i$. We now verify the number of remaining vertices in each part by direct calculation: 
\begin{itemize}
    \item $|A_1\setminus V(\mathcal{K})|= a_1-\sum_{i=2}^p(n-|A_i|)=a_1+\sum_{i=2}^pk_i$;
    \item for each $\ell\in [2,p]$, 
        \begin{align*}
        |A_\ell\setminus V(\mathcal{K})|=&a_1+a_\ell-\sum_{i=2}^p(n-|A_i|)+(n-|A_\ell|)-\sum_{i=p+1}^r(k_{i}-k_i')\\
        =&n-\sum_{i=p+1}^rk_{i}=a_1-k_1+\sum_{i\in [p]}k_{i}=a_1+\sum_{i=2}^pk_i;
        \end{align*}
        
    \item for each $\ell\in [p+1,r]$,                  \begin{align*}
        |A_\ell\setminus V(\mathcal{K})|=&a_1+a_\ell-\sum_{i=2}^p(n-|A_i|)-k_{\ell}'-\sum_{i=p+1}^r(k_{i}-k_i')-(k_{\ell}-k_{\ell}')\\
        =&|A_{\ell}|-\sum_{i=p+1}^rk_{i}-k_{\ell}=n-\sum_{i=p+1}^rk_{i}=a_1+\sum_{i=2}^pk_i,
        \end{align*}
\end{itemize}
as desired. 
\end{proof}

Based on Claim \ref{result-balance}, we know $|A_{\ell}\setminus V(\mathcal{K})|=a_1+\sum_{i=2}^pk_i$ for all $\ell\in [r]$. Together with Theorem \ref{lem:HS thm in multipartite graph}, we know that $G-V(\mathcal{K})$ has a $K_r$-factor. Thus, $G$ admits a $K_r$-factor. 
\end{proof}
%In the rest of the proof, our main challenge is to adjust the partition $\mathcal{P}$ slightly to satisfy the following degree condition.



\subsection{Covering bad vertices for balance}
In this section, in order to derive the model in Lemma~\ref{lem:balance}, we need to cover bad vertices that violate the degree condition~\ref{A3} by $K_r$-tiling. Since $G$ is $((r-1)n+1)$-regular and $\Delta(G[A_i])\leq \beta n$ for each large part $A_i$, every vertex belonging to $A_i$ already meets condition~\ref{A3}.
Thus, we only need to address two types of bad vertices:
\begin{itemize}
    \item vertices of large degree within small parts. Claim~\ref{fact:bad vertices for Ai}~\ref{large-degree vtx} implies that there are few bad vertices of the first type.
    \item vertices in $B$ with many non-neighbors in some $A_i$, i.e., $(\beta,B)$-bad vertices. %By Lemma~\ref{good-partition}~\ref{B04}, every vertex $v\in B$ has at least $\beta n$ neighbors in $A_i$ for each $i\in[s]$. However, there may exist some vertices in $B$ which have many non-neighbors in some $A_i$, i.e., .
\end{itemize}
For such bad vertices, we attempt to cover them via a $K_r$-tiling so that the remaining partition satisfies a degree condition in Lemma~\ref{lem:balance}. A key requirement of this covering process is to ensure that the number of vertices that used in each part is balanced. 
We first prove that the number of $(\beta,B)$-bad vertices is small. 


%Note that each vertex in large part satisfies condition~\ref{A3} as $G$ is $((r-1)n+1)$-regular and $\Delta(G[A_i])\leq \beta n$. Thus, we are left to handle only the two types of bad vertices: one is large-degree vertices in small parts and the other is vertices in $B$ that do not have a sufficiently high degree. Crucially, this covering process should guarantee that the resulting partition is balanced again. We first show that the number of those bad vertices is also small.

% 

%For such bad vertices, we attempt to cover them via a $K_r$-tiling so that the remaining partition satisfies a degree condition analogous to that in Lemma~\ref{lem:balance}. To achieve this goal, we define bad vertices in $B$ and show that their number is small.

%\begin{definition}\label{def:bad vertex for B}

%We say a vertex $v\in B$ is bad if $d_{A_i}(v)\le |A_i|-\beta n$ for some $i\in [s]$.
%\end{definition}
%counting bad vertices


% Suppose that the number of bad vertices in $B$ is at least $\alpha^{1/4} n$.
% By the Pigeonhole Principle, there exists some $i\in [s]$ such that there are at least $\frac{|B_b|}{s}$ vertices $v\in B$ satisfying $d_{A_i}(v)\le |A_i|-\beta n$.
% Since $d(v)=(r-1)n+1$ for each $v\in A_i$, we obtain that 
% $$|E(G[A_i])|\geq \frac{1}{2}\cdot\frac{\alpha^{1/4} n}{s} \cdot 2\alpha^{1/5}n\ge \sqrt{\alpha}n^2.$$ 
% Recall that $A_i^0$ is a $\gamma_i$-independent set of $G$ and at most $r\alpha n$ vertices are moved from other part to $A_i^0$. Hence 
% $$
% |E(G[A_i])|\leq (\gamma_i+2r\alpha) n^2<\sqrt{\alpha}n,
% $$
% a contradiction. 
% Therefore, 
% \begin{align}\label{size-of-ne}
%    |V_{ne}(3\alpha)^{1/5}^{1/5},B)|\le \alpha^{1/4} n\ \text{and}\ |V_{ne}(3\alpha)^{1/5}^{1/5})|\leq 2\alpha^{1/4}n.
% \end{align}






The following lemma provides a method for covering the two types of bad vertices with a $K_r$-tiling and adjusting the current partition into a balanced one.
\begin{lemma}\label{lem:covering bad vertices for balance}
Let $r,n$ be positive integers with $r\ge 3,\ 1\le s\le r$. Assume that $0<\frac{1}{n}\ll \gamma\ll\alpha\ll\beta\ll1$. Let $G$ be a graph on $rn$ vertices with $(r-1)n+1-n^{0.1}\leq \delta(G)\leq \Delta(G)\leq (r-1)n+1+n^{0.1}$, and $\mathcal{P}=\{A_1,\cdots,A_s,B\}$ be a partition of $V(G)$ satisfying~\ref{B1}-\ref{B07} in Lemma~\ref{good-partition}. Then $G$ has a $K_r$-factor. % partition $\mathcal{P}$ can be transformed into a balanced one by removing a $K_r$-tiling of seize at most ?? vertices.
\end{lemma}

\begin{proof}
Choose additional constants $\alpha^{1/5}$ and $\eps$ such that
\[
\alpha\ll\beta\ll\eps\ll\frac{1}{r}.
\]
Our goal is to use Lemma \ref{lem:balance} to construct a $K_r$-factor in $G$. For this, we first  consider the following two types of vertices:
\begin{enumerate}[label =\rm  {\bf Type \arabic{enumi}}, leftmargin=6em]
    \item\label{v1} $v\in A_i$ and $d(v,A_i)\geq \alpha^{1/5} n$ for some $i\in [s]$ with $|A_i|\leq n$;
    \item\label{v2} $v\in B$ and $d(v,A_i)\leq |A_i|-\alpha^{1/5} n$ for some $i\in [s]$, i.e., $v$ is $(\alpha^{1/5},B)$-exceptional. 
\end{enumerate}

Let $v\in V(G)$ be a vertex of \ref{v1}. It follows from \ref{B04}-\ref{B05} that $d(v,D)\geq \beta n$ for each $D\in \mathcal{P}\setminus \{A_i\}$. Together with \ref{B6} and $\alpha\ll \beta$, one may choose a copy of  $K_s$ that covers $v$ and $V(K_s)\cap A_j$ consists of a unique $(\alpha^{1/5},j)$-good vertex for each $j\in [s]\setminus \{i\}$. Moreover, each of such $(\alpha^{1/5},j)$-good vertex has at least $(\beta-(r-1)\alpha^{1/5})n$ choices. Next, we are to find a copy of $K_{r-s}$ in $N_B(\{v_1,\ldots,v_s\})$. Let $G'=N_B(\{v_1,v_2,\cdots,v_s\})$. If $s=\{r,r-1\}$, then we are done. In what follows, we only consider $s\leq r-2$. 
Recall that \ref{B05} implies that
\[
\de(G[B])\ge (r-1)n+1-sn-r\alpha n\ge \left(1-\frac{1}{r-s}-\alpha^{1/5}\right)\cdot|B|.
\]


%\medskip
%{\bf Covering bad vertices in $A_i$ with $|A_i|\le n$.}
%\medskip

%We cover each worse vertex $v_i$ in $A_i$ with a copy of vertex-disjoint $K_r$ if $d_{A_i}(v_i)\ge\alpha^{1/5} n$.
%For each bad vertex $v_i\in A_i$, the condition~\ref{B04} and~\ref{B05} in Lemma~\ref{lem:partition after adjust vertices} imply that $v_i$ has at least $\alpha^{1/5} n$ neighbors in each of the other parts of the partition. By Fact~\ref{fact:bad vertices for Ai}~\ref{large-degree vtx}-\ref{bad vtx2}, one can choose an $(\alpha^{1/5}, j)$-bad vertex $v_j$ in $A_j$ for each $j\neq i$ such that $v_1,v_2,\cdots,v_s$ induces a $K_s$ as $\alpha\ll\alpha^{1/5}$. Note that such vertex $v_j$ is adjacent to all but at most $\alpha^{1/5} n$ vertices in $V(G)\setminus A_j$. We now try to find a copy of $K_{r-s}$ in the common neighborhood of $v_1,v_2,\cdots,v_s$ in $B$. If $s=r$, then we are done since we have identified such $r$ vertices. Let $G'=N_B(\{v_1,v_2,\cdots,v_s\})$. Since $G$ is $((r-1)n+1)$-regular and $\alpha\ll\alpha^{1/5}$, we have
%\[
%\de(G[B])\ge (r-1)n+1-sn\ge \left(1-\frac{1}{r-s}-\alpha^{1/5}\right)\cdot|B|.
%\]
Since $s\leq r-2$, one has $(r-s)n\ge |B|\ge (r-s-r\alpha)n>n$. Let $x:=|B|-|V(G')|$. Clearly, $x\leq n-2+(r-1)\alpha^{1/5} n$. % , then $d(v_i,B)\ge |B|-n+1$ for each $i\in [s]$. Let $x$ be the number of non-neighbors of $v_1,v_2,\cdots,v_s$. As $v_j$ is a $(\beta,A_j)$-bad vertex in $A_j$ for each $i\neq j$, we have
%\[
%|B|-x=|G'|\ge |B|-n+1-(s-1)\alpha^{1/5} n,
%\]
%which means that $x\le n+(s-1)\beta n$.
Thus, %as $\alpha\ll\alpha^{1/5}\ll\eps\ll 1$ we have
\begin{align*}
\de(G')&\ge \de(G[B])-x\ge \left(1-\frac{1}{r-s}-\alpha^{1/5}\right)\cdot|B|-x=|B|-x-\frac{|B|}{r-s}-\alpha^{1/5}|B|\\
&\ge|B|-x-\frac{|B|-x}{r-s-1}-(|B|-x)\eps 
=\left(1-\frac{1}{r-s-1}-\eps\right)\cdot(|B|-x)\\
&=\left(1-\frac{1}{r-s-1}-\eps\right)\cdot|V(G')|.
\end{align*}
where the last inequality follows from $\alpha^{1/5}\ll\eps$ and the fact $\frac{|B|}{r-s}\le\frac{|B|-x}{r-s-1}$ as $x\le n+(s-1)\beta n$.
By Theorem~\ref{SuperT}, we can find a $K_{r-s}$ in $G'$, as desired.%独立集版本的Turan thm

%If $s=r-1$, that is, we already have the vertices $v_1,v_2,\cdots,v_{r-1}$, then we only need to choose a vertex in the common neighborhood of $v_1,v_2,\cdots,v_{r-1}$ in $B$. Note that $B$ has no $\gamma_{r}$-independent set and $d(v,B)\geq \beta n$ for each vertex $v\in V(G)$. As $v_j\in A_j$ is a $(\beta,A_j)$-bad vertex for each $j\neq i$ and $\beta$, we have \[
%|G'|=d_B(v_1,v_2,\cdots,v_{r-1})\ge\beta n-(r-1) \alpha^{1/5}n>\frac{\beta n}{2}>0,
%\]
%which implies that we can find $v_r\in B$ such that $\{v_1,v_2,\cdots,v_r\}$ induced a $K_r$ in $G$.

Since there are at most $r\alpha n$ vertices of \ref{v1} and $\alpha\ll\beta\ll\eps$, the above operations yield a $K_r$-tiling that covers all  vertices of \ref{v1} .

%\medskip
%{\bf Covering bad vertices in $B$.}
%\medskip
Now, we consider $(\alpha^{1/5},B)$-exceptional vertices. The following claim guarantees the existence of a $K_{r-s}$-tiling that covers all these vertices. Moreover, each clique in this tiling contains exactly one $(\alpha^{1/5},B)$-exceptional vertex.%containing $v$ in $B$ such that the remaining $r-s-1$ vertices are not bad in $B$. % let $v\in V(G)$ be a vertex of \ref{v2}. 

%We cover each worse vertex $v$ in $B$ with a copy of vertex-disjoint $K_r$ if $d_{A_i}(v)\le |A_i|-\alpha^{1/5} n$ for some $i\in [s]$. 
\begin{claim}\label{claim:good Kr-s in B}
For any $W\subseteq B$ with $|W|\le\eps n$, every vertex $w\in B\setminus W$ is covered by a $K_{r-s}$.
\end{claim}
\begin{proof}[Proof of Claim \ref{claim:good Kr-s in B}]
The case for $r-s=1$ is trivial, so we only consider $r-s\geq 2$. Let $G_1:=G[B\setminus W]$.
It suffices to show that for each every  vertex $w\in B$, there is a copy of $K_{r-s-1}$ in $G_1[N_{G_1}(w)]$.
Note that $$|N_{G_1}(w)|\ge\de(G_1)\ge \de(G[B])-\eps n\ge \left(1-\frac{1}{r-s}-2\eps\right)\cdot|B|.$$
For any vertex $u\in N_{G_1}(w)$, we have 
\begin{align*}
&|N_{G_1}(w)\cap N_{G_1}(u)|
\ge |N_{G_1}(w)|+|N_{G_1}(u)|-|G_1|\\
\ge& |N_{G_1}(w)|+\left(1-\frac{1}{r-s}-2\eps\right)\cdot|B|-(|B|-|W|)\\
\ge& |N_{G_1}(w)|-\left(\frac{1}{r-s}+2\eps\right)\cdot|B|+|W|\\
\geq &\left(1-\frac{1}{r-s-2}-\eps\right)\cdot|N_{G_1}(v)|+|W|,
\end{align*}
which implies $\de(G[N_{G_1}(v)])\ge \left(1-\frac{1}{r-s-2}-\eps\right)\cdot|N_{G_1}(v)|$.
Therefore, by Theorem~\ref{SuperT}, we can find a $K_{r-s}$ in $G_1$, as desired.
\end{proof}

Claim~\ref{claim:good Kr-s in B} implies that there exists a $K_{r-s}$-tiling that covers $(\alpha^{1/5},B)$-exceptional  vertices. Moreover, each clique in this tiling contains exactly one $(\alpha^{1/5},B)$-exceptional vertex. For the $(\alpha^{1/5},B)$-exceptional vertex $v\in B$, let $K^v$ be the $K_{r-s}$ copy that covers $v$, with $V(K^v)=\{v,u_1,\ldots,u_{r-s-1}\}$. Observe that $d(u_i,A_j)\geq |A_j|-\alpha^{1/5}n$ for each $i\in [r-s-1]$ and each $j\in [s]$. It follows from Lemma \ref{good-partition} \ref{B04} that $d(v,A_j)\geq \beta n$ for each $j\in [s]$. Together with \ref{B2}, \ref{B07} and $\alpha\ll \beta$, there exists a copy of \( K_r \) covering \( K^v \) such that all vertices in \( V(K_r) \setminus V(K^v) \) are \( (\alpha^{1/5}, i) \)-good. Moreover, the \( K_r \)-copies corresponding to distinct \( (\alpha^{1/5}, B) \)-exceptional vertices are pairwise disjoint.
%For each $(\alpha^{1/5},B)$-exceptional vertex $v\in B$, Claim~\ref{claim:good Kr-s in B} implies that there exists a copy of $K_{r-s}$ covering $v$ in $B$ external to the vertices already used in the covering. Moreover, the other $r-s-1$ vertices in this $K_{r-s}$ are all good and we denote them by $u_1,\cdots,u_{r-s-1}$. Note that each $u_i$ has at least $|A_i|-\alpha^{1/5}n$ vertices in $A_i$. Combing Lemma~\ref{lem:partition after adjust vertices}~\ref{B04} and Fact~\ref{fact:bad vertices for Ai}~\ref{large-degree vtx}\ref{bad vtx3}, for each $j\in[s]$ we can choose a $(\beta,A_j)$-bad vertex $v_j\in A_j$ in the common neighbor of $v,u_1,\cdots,u_{r-s-1}$ such that $\{v_1,v_2,\cdots,v_s\}$ induces a $K_s$ as $\alpha\ll\beta$. Therefore, $\{v_1,\cdots,v_s,v,u_1,\cdots,u_{r-s-1}\}$ is a copy of $K_r$ covering $v$.

Let $\mathcal{K}$ be the $K_r$-tiling that used to covering vertices of \ref{v1} and \ref{v2}. By \ref{B2}, we have $|V(\mathcal{K})|\leq r^2\alpha n$. %Let $\mathcal{P}':=\{A_1',\ldots,A_s',B'\}$ be a partition of $G':=G-V(\mathcal{K})$, which is obtained from $\mathcal{P}$ by deleting all vertices in $V(\mathcal{K})$.
Assume that $|B\setminus V(\mathcal{K})|\equiv q \pmod{r-s}$. %The final step is to deal with the case that the number of vertices in $B'$ is not divisible by $r-s$.  %\medskip
%{\bf Covering non-divisible vertices in $B$.}
%\medskip
%Fact~\ref{fact:bad vertices for Ai} and Fact~\ref{fact:bad vertices in B} imply that the covering process in the first two steps requires no more than $10r\alpha n$ vertices.
As $\alpha\ll\beta$, we have
\[
\de(G[B\setminus V(\mathcal{K})])\ge (r-1)n+1-sn-2r^2\alpha n\ge \left(1-\frac{1}{r-s}-\beta\right)\cdot|B\setminus V(\mathcal{K})|.
\]
By Theorem~\ref{SuperT}, there are $q$ vertex-disjoint copies of $K_{r-s+1}$, say $H$, in $G[B\setminus V(\mathcal{K})]$. %Denote $B''=B'\setminus V(H)$.  
Note that each vertex $v$ in $B\setminus V(\mathcal{K})$ satisfies $d(v,A_i)\ge |A_i|-\alpha^{1/5} n$ for all $i\in [s]$. Thus, for each $i \in [s-1]$, we can greedily choose a vertex $v_i \in A_i$ that lies in the common neighborhood of the vertices from the $K_{r-s+1}$ in $H$. It is obvious that $K_{r-s+1}\cup\{v_1,\cdots,v_{s-1}\}$ forms a copy of $K_r$ in $G$. Let $\mathcal{K}'$ be the $K_r$-tiling that used to covering cliques in $H$. Let $\mathcal{P}':=\{A_1',\ldots,A_s',B'\}$ be a partition of $G':=G-V(\mathcal{K}\cup H)$, which is obtained from $\mathcal{P}$ by deleting all vertices in $V(\mathcal{K}\cup H)$. %\medskip %After the covering procedures, we obtain a $K_r$-tiling of size at most $10r\alpha n$ in $G$ covering all kinds of bad vertices. We remove such $K_r$-tiling and get a new unbalanced partition $\mathcal{P}$.  To obtain a balanced one as in Lemma~\ref{lem:balance}, we aim to find a $K_{r-s}$-factor in $B$ and identify each copy of $K_{r-s}$ with a single vertex. Let $A_{s+1}$ be a resulting set of size exactly $\frac{|B|}{r-s}$.

Next, we are to find a $K_{r-s}$-factor in $G'[B']$. It is trivial for $r-s=1$. Note that $\de(G'[B'])\ge \left(1-\frac{1}{r-s}-\beta\right)|B'|$ and $B'$ has no $\gamma'$- independent set of size $n$. 
If $r-s\ge 3$, Theorem~\ref{thm:Non-extremal case r>3} implies that there is a $K_{r-s}$-factor in $G'[B']$, as desired.
Suppose \( r - s = 2 \). Then by Lemma~\ref{thm:r-s=2}, either \( B' \) contains a perfect matching (in which case we are done), or it is \( 3\beta \)-closed to \( 2K_{|B'|/2} \) with odd $|B'|/2$. In the latter case, \( B' \) can be partitioned into two disjoint sets \( B_1 \) and \( B_2 \), each being \( 3\beta \)-closed to \( K_{|B'|/2} \). Observe that \ref{B07} implies that there is an edge, say $ab$, between $B_1$ and $B_2$. Applying Lemma~\ref{thm:r-s=2} again yields that $G[B'\setminus \{a,b\}]$ contains a perfect matching, as desired. 

%{\color{red}If $|B'|/2$ is even, then it is easy to find a perfect matching in $B'$. If $|B'|/2$ is odd, we claim that there is a crossing edge between $B_1$ and $B_2$. Indeed, since $G$ is $((r-1)n+1)$-regular, for any balanced partition $\{B_1^0,B_2^0\}$ of $B^0$, there are at least $n$ edges between $B_1^0$ and $B_2^0$. The above procedure preserves a linear number of edges between $B_1$ and $B_2$. Then we choose a crossing edge $e=\{ab\}$ such that both $|B_1\setminus\{a\}|$ and $|B_2\setminus\{b\}|$ are even and we can find a $K_r$-factor in $B$.}

Notice that for each $D'\in \mathcal{P}'$ and each $v\notin D'$ we have 
\begin{align}\label{DofQ'}
d(v,D')\ge |D'|-\beta n -r^3\alpha n\ge |D'|-2\beta n.
\end{align}
Notice that there exist $t:=\frac{|B''|}{r-s}$ vertex-disjoint copies of $K_{r-s}$, say $K^1,\ldots,K^{t}$, in $G'[B']$, and $q$ vertex-disjoint copies of $K_{r-s+1}$, say $K^{t+1},\ldots,K^{t+q}$, in $H$. Contracting each $K^i$ (where $i\in [t+q]$)  into a vertex $w_i$ yields a new vertex set $B^*$ of size $t+q$. 
We are to construct an auxiliary graph $G^*$ with $V(G^*)=\big (\bigcup_{i\in[s]}A_i'\big)\cup B^*$, and $xy\in E(G^*)$ if and only if 
\begin{itemize}
    \item $x\in A_i'$, $y\in A_j'$, and $xy\in E(G)$ for $i,j\in[s]$ with $i\ne j$; or 
    \item $x=w_i\in B^*$, $y\in A_j'$, and $y\in N(V(K^i))$ for $j\in [s]$ and $i\in[t]$.
\end{itemize}
Thus, $\mathcal{Q}^*=\{A_1',\ldots,A_s',B^*\}$ is a partition of $G^*$. 
Then  \eqref{DofQ'} implies that  
$$
d_{G^*}(v,D^*)\geq |D^*|-(r-s)\beta n\geq \Big(1-\frac{1}{2(s+1)}\Big)|D^*| \ \text{for each}\ D^*\in \mathcal{Q}^* \ \text{and each}\ v\notin D^*.
$$
Moreover, 
\begin{itemize}
    \item\label{C1} $n-100r^2\alpha n\leq |A_1'|,\ldots, |A_{s}'|,|B^*|\leq n+100r^2\alpha n$;
    \item\label{C2} for each $A_i'$ with $|A_i'|>n$, $G[A_i]$ contains a matching of size at least $|A_i'|-\frac{|G^*|}{s+1}$;
    \item if $r-s=1$, then there are at least $r\alpha n$ matching edges inside $G'[B']$ as $B'$ is not a $\gamma'$-independent set. 
\end{itemize}
By Lemma~\ref{lem:balance}, there is a $K_r$-factor in $G$ and we complete the proof.
\end{proof}
\medskip
{\bf Remark.}
As the proof demonstrates of Lemma~\ref{good-partition}, condition~\ref{B2} is overly strong.
In fact, Lemma~\ref{lem:balance} implies that the matching of size $|A_i|-n$ in each large part $A_i$ is sufficient to find a $K_r$-factor.
Furthermore, condition~\ref{B6} can be relaxed to only require the presence of at least one edge between any two sets of size $n$ in $B$, which is enough to guide us in finding a perfect matching in $B$ when $r-s=2$.




% \begin{definition}
% Let \( r \geq s+1 \geq 2 \) and \(\varepsilon < 1/r\). Let \( t_1, t_2, \ldots, t_{s+1} \) be \( s+1 \) integers such that \( |t_i| \leq \varepsilon n \) for each \( i \in [s+1] \) and \(\sum_{i=1}^{s+1} t_i = 0\). Let \( G \) be an \( rn \)-vertex graph and \((V_1, V_2, \ldots, V_{s+1})\) be a partition of \( V(G) \) with \( |V_i| = n + t_i \) for each \( i \in [s] \) and \( |V_{r'}| = (r - s + 1)n+ t_{s+1} \). Let \( \mathbf{b} = (1, \ldots, 1, r - s) \) be an \( (s+1) \)-dimensional vector. For any \( i \in [s+1] \), let \( \mathbf{u}_i \) be the \( i \)-th unit vector, i.e.\ \( \mathbf{u}_i \) has 1 on the \( i \)-th coordinate and 0 on the other coordinates. Let \( F \) be a \( K_r \)-tiling of size \( f \) in \( G \).

% \begin{itemize}
%     \item The \textup{index vector of \( F \)} is \( i(F) = (x_1, x_2, \ldots, x_{s+1}) \) by choosing \( x_i = |V(F) \cap V_i| \) for every \( i \in [s+1] \).

%     \item The \textup{slack} of \( F \) is \( {\ell}(F) = (\ell_1, \ell_2, \ldots, \ell_{s+1}) := i(F) - f\mathbf{b} \). Note that \(\sum_{i=1}^{s+1} \ell_i = 0\) and \(\sum_{\ell_i > 0} \ell_i \leq (r - 1)f\).

%     \item The  \textup{slack of the partition} \((V_1, V_2, \ldots, V_{r'})\) is \((t_1, t_2, \ldots, t_{r'}) = (|V_1|, |V_2|, \ldots, |V_{r'}|) - n\mathbf{b}\). The partition \((V_1, V_2, \ldots, V_{r'})\) is \textup{balanced} if \( |V_1| = \cdots = |V_{s}| = |V_{s+1}|/(r -s)\).

%     \item A copy of \( K_r \) in \( G \) is \textup{balanced} if its index vector is \( \mathbf{b} \). Compared with a balanced copy of \( K_r \), an \((i, j)\)-copy of \( K_r \) has index vector \( \mathbf{b} + \mathbf{u}_i - \mathbf{u}_j \), where \( i, j \in [r'] \) are distinct.

%     \item For a copy of \( K_r \), we refer to its index vector as its \textup{type}. %Note that there are \( y := \binom{r+s}{s} \) different types of \( K_r \).

%     \item Let \( I \) be a subset of \( [s+1] \) and \( \mathbf{v} \) be an \( (s+1) \)-dimensional vector. We use \( \mathbf{v}|_I \) to denote the 
% $|I|$-dimensional vector obtained by restricting ${\bf v}$ to the coordinates in $I$.
% \end{itemize}
% \end{definition}



\section{Random sample and the robustness}
%希望intro.介绍该问题与robustness
\subsection{Proof of Lemma~\ref{lem:extremal-case}}
In this section we prove our main result Lemma~\ref{lem:extremal-case} using the sufficient conditions for $K_r$-factors in Lemma~\ref{lem:covering bad vertices for balance}.
As the size of a random subset of a set $A$ has the binomial $B(|A|, 1/2)$ distribution, we need the following observation on the difference of independent binomials.

\begin{lemma}\label{lem:binomial distribution}
If random sets $X\sim B(n,1/2)$ and $Y\sim B(m,1/2)$ are independent, then $X+m-Y\sim B(n+m,1/2)$.
\end{lemma}

\begin{lemma}\label{lem-int}
    If $X\sim B(n,1/2)$, then $\mathbb{P}(a\sqrt{n}/2\leq X-n/2\leq b\sqrt{n}/2)=\int_a^b \frac{1}{\sqrt{2\pi}}e^{-t^2/2}\,dt+o(1)$.
\end{lemma}
The following obvious fact illustrates that minimum degrees can be preserved with high probability in the random vertex samples.
\begin{fact}
Given a bipartite graph $G=(A_1,A_2)$ on $2n$ vertices, if $\de(G)\ge \beta n$, then $\de(G[\frac{1}{2}])\ge \frac{\beta n}{4}$ with probability $1-o(1)$.
\end{fact}
\begin{proof}
Consider $G[\frac{1}{2}]=G[S]$ where $S$ is a random subset of $V(G)$.
For every vertex $v \in S$, since $d_G(v) \ge \beta n$, we have $\mathbb{E}[d_{G[S]}(v)]\ge \frac{1}{2}\cdot \beta n$.
By chernoff bound, w.h.p.~we have $d_{G[S]}(v)\ge \frac{\beta n}{4}$ for every $v \in S$.
\end{proof}

The required matchings in Lemma~\ref{good-partition} will be provided by the following two lemmas.
The first claim gives that the matching in each large part can be preserved with high probability after the random vertex sampling, which plays an essential role in this paper.

\begin{lemma}\label{lem:inherit matching}
For any part $A_i$ with $|A_i|=\frac{n}{r}+k_i$, if $k_i>0$ then the random set $A_i \cap S$ admits a matching of size at least $\frac{k_i}{32\beta}$ with probability $1-o(1)$.
\end{lemma}


The second lemma gives that for a balanced partition, there is a matching in some part, which uses regularity.
\begin{lemma}\label{vertex-cover}
Let $G$ be an $((r-1)n+1)$-regular graph on $rn$ vertices and $(A_1,A_2,\ldots,A_r)$ be a balanced partition. Suppose $A_i'$ is a minimum vertex cover of $G[A_i]$ for each $i\in [r]$. Then $\max\{|A_i'|:i\in [r]\}\geq \frac{1}{r}\sqrt{n}$.
\end{lemma}

\begin{lemma}\label{lem:extremal-case}
For any $r\ge 3$, let $\gamma>0$ and $G$ be an $((r-1)n+1)$-regular graph on $rn$ vertices. Suppose that $G$ contains a  $\gamma$-independent set of size $n$. Then  $G[\frac{1}{2}]$ admits a $K_r$-factor with probability at least $\frac{1}{200r^2}$.
\end{lemma}

\begin{proof}[Proof of Lemma~\ref{lem:extremal-case}]
Let $G$ be an $((r-1)n+1)$-regular graph on $rn$ vertices and $G$ contains a $\gamma$-independent set of size $n$.
Based on the above discussion, there exists a partition $\mathcal{P}:=\{A_1,\cdots,A_s,B\}$ of $V(G)$ satisfying~\ref{B1}-\ref{B07}.

Consider $G[\frac{1}{2}]=G[S]$ where $S$ is a random subset of $V(G)$ with $r\mid |S|$.
Note that $S\sim B(rn,1/2)$.
By Chernoff bounds, w.h.p.~$|S|=\frac{rn}{2}\pm n^{0.1}$, all vertex degrees are $(r-1)\frac{|S|}{r}\pm |S|^{0.1}$, $|S\cap A_i|=\frac{|A_i|}{2}\pm n^{0.1}$ for each $i\in [s]$, and $|S\cap B|=\frac{|B|}{2}\pm n^{0.1}$.
Moreover, by chernoff bound, we have the following result.
\begin{claim}\label{claim:B1-B06 in S}
Given constants $\frac{1}{n}\ll\gamma\ll\alpha\ll\beta\ll1$, there exists a constant $\gamma'\ll\frac{1}{r}$ such that the random partition $\mathcal{P}'=\{A_1\cap S,\cdots,A_s\cap S,B\cap S\}$ w.h.p. has the following properties:
\begin{enumerate}
[label =\rm  (D\arabic{enumi})]
    \item\label{D1} $|A_i\cap S|\le\frac{n}{2}+ n^{0.1}$ if $|A_i|\le n$ and $||A_i\cap S|-\frac{n+k_i}{2}|\le n^{0.1}$ if $|A_i|=n+k_i$ where $k_i>0$;
    \item\label{D2} $\Delta(G[A_i\cap S])\leq \beta n$ for each $i\in [s]$;
    \item\label{D3} $d(v,A_i\cap S)\geq \frac{\beta n}{4}$ for each vertex $v\in V(G)\setminus A_i$;
    \item\label{D4} $d(v,B\cap S)\geq \frac{\beta n}{4}$ for each vertex $v\in V(G)$ and $\de(G[B\cap S])\ge $;
    \item\label{D5} the number of $(\frac{\beta}{2},B\cap S)$-exceptional vertices is at most $\frac{r\alpha n}{2}$;
    \item\label{D6} for any two sets of size $\frac{|S|}{r}$ in $B\cap S$ there are at least $\frac{n}{4}$ edges between them;
    \item\label{D7} $d(v,D)\ge |D|-\alpha^{1/5}n$ for every $(\alpha^{1/5},A_i)$-good vertex $v\notin D$ where $D\in \mathcal{P}'$;
    \item\label{D8} $A_i\cap S$ is a $\gamma'$-independent set for each $i\in [s]$ and $B\cap S$ has no $\gamma'$-independent set of size $n/2$.
\end{enumerate}
\end{claim}


Recall that $A_i\cap S$ is called a large part if $|A_i\cap S|>\frac{|S|}{r}$, and a small part otherwise.
By Lemma~\ref{lem:covering bad vertices for balance} and Lemma~\ref{good-partition}, as all vertex degrees are $(r-1)\frac{|S|}{r}\pm |S|^{0.1}$, it suffices to find a matching of size at least $|A_i\cap S|-\frac{|S|}{r}$ in each large part $A_i\cap S$ to obtain a $K_r$-factor.


% \begin{proof}
% Given $A_i$ with $|A_i|=\frac{n}{r}+k_i$, we know that $\de(G[A_i])\ge k_i+1$ and $e(G[A_i])\ge \frac{1}{2}|A_i|(k_i+1)$.
% Let $X=\sum_{e\in E(G[A_i])}X_e$ where $X_e=1$ if $e\in E(G[A_i])$, otherwise $X_e=0$.
% Note that $X\sim B(e(G[A_i]),1/2)$ and $E[X]=\frac{1}{4}e(G[A_i])$.
% Let $S_i$ be an indicator variable which denotes whether $v_i$ is selected, that is, 
% \[
% S_i = 
% \begin{cases} 
% 1, & v_i \in S; \\
% 0, & v_i \notin S .
% \end{cases}
% \]
% Note that $\{S_i\}_{i=1}^{|A_i|}$ are independent and identically distributed random variables defined on a complete probability space.
% Let $\mathcal{F}_k=\sigma(S_1,S_2,\cdots,S_k)$ and $Z_k=E[X\mid S_1,\cdots,S_k]$.
% Since
% \[
% E[Z_k\mid \mathcal{F}_{k-1}]=E[E[X\mid \mathcal{F}_k]\mid \mathcal{F}_{k-1}]=E[X\mid \mathcal{F}_{k-1}]=Z_{k-1},
% \]
% we obtain that $\{Z_k\}_{k\ge 0}$ is a martingale.
% Note that $Z_0=E[X]$, $Z_{|A_i|}=X$ the martingale difference is $$|Z_k-Z_{k-1}|\le\Delta(G[A_i])\le \beta n.$$
% Applying with Azuma-Hoeffding inequality, we have
% \[
% P[|X-E[X]|\ge t]\le 2\exp\left(-\frac{2t^2}{|A_i|\cdot(\beta n)^2}\right)=o(1),
% \]
% which implies that there are at least $\frac{1}{4}e(G[A_i])-n$ edges in $A_i\cap S$ with probability $1-o(1)$.
% Therefore, we have
% \[
% \frac{1}{4}e(G[A_i])-\beta n^{2/3}-2ck_i\cdot\Delta(G[A_i])\ge k_i\cdot(\frac{1}{4}|A_i|-2c\beta n)+\frac{1}{4}|A_i|-\beta n^{2/3}
% \]


% \end{proof}


% Since the partition $\mathcal{P}'$ inherits the conditions on degree and independent sets with high probability, we can also find a $K_r$-tiling of size at most $10r\alpha n$ that cover all bad vertices, just as in subsection 4.2.
% Removing such $K_r$-tiling, we obtain a new partition $\mathcal{P}=\{A_1\cap S,\cdots,A_s\cap S,B\cap S\}$ (adhere to the previous notation for convenience) satisfying
% \begin{enumerate}
% [label =\rm  (E\arabic{enumi})]
%     %\item\label{E1} $||A_i\cap S|-\frac{|S|}{r}|\le |S|^{0.1}$ if $|A_i|\le n$ and $||A_i\cap S|-\frac{|S|}{r}|\le k_i+$ if $|A_i|>n$ where $k_i=|A_i|-\frac{|S|}{r}$;
%     \item\label{E1} $\Delta(G[A_i\cap S])\leq \beta |A_i\cap S|+n^{0.1}$ for each $i\in [s]$;
%     \item\label{E2} $d(v,A_i\cap S)\geq |A_i\cap S|-\beta n-n^{0.1}$ for each vertex $v\in S\setminus A_i$;
%     \item\label{E3} $d(v,B\cap S)\geq |B\cap S|-\beta n-n^{0.1}$ for each vertex $v\in S\setminus B$;
%     \item\label{E4} $(r-s)\mid |B|$ and $B\cap S$ has no $\frac{\gamma_{s+1}}{2}$-independent set of size $n/2$.
% \end{enumerate}

% %Without loss of generality, assume that $|A_i\cap S|=\frac{|S|}{r}+k_i$ for each $i\in [s]$ and $k_1\leq \cdots\leq k_s$ where $\sum_{i\in [s]}k_i=0$.
% To complete the proof via Lemma~\ref{lem:balance} and Claim~\ref{lem:inherit matching}, it remains to show that there exists a large matching in each large part.
Now we consider two cases according to the size of $k_i=|A_i|-n$.

\medskip
{\bf Case 1. $k_i>\de n^{0.2}$ for some $i\in[s]$.}
\medskip

Note that $k_1\le k_2\le\cdots \le k_s$ and $k_s\ge \de n^{0.2}$.
Together with $|S|=\frac{rn}{2}\pm n^{0.1}$, we know that $\frac{|S|}{r}\ge \frac{n}{2}+\de n^{0.1}$.
Thus, Claim~\ref{claim:B1-B06 in S}\ref{D1} gives that if $A_i$ is a small part or $k_i\le\de n^{0.2}$ with respect to $\mathcal{P}$, then w.h.p.~$A_i\cap S$ is also a small part with respect to $\mathcal{P}'$.
If $k_i>\de n^{0.2}$, Lemma~\ref{lem:inherit matching} implies that there is a matching of size at least $\frac{k_i}{32\beta}$ in $A_i\cap S$ with probability $1-o(1)$.
Therefore, we verify that condition~\ref{B2} in Lemma~\ref{good-partition} holds.
Thus, we have
\[
P[G[\frac{1}{2}]~\text{contains a $K_r$-factor}]\ge P[r\mid |S|]-o(1)\ge \frac{1}{r}-o(1).
\]

\medskip
{\bf Case 2.} 
$k_s\leq \de n^{0.2}$.
\medskip

{\bf Subcase 2.1.} $r-s=0$.
\medskip

We consider any balanced partition $\{A_1^*,A_2^*,\ldots,A_r^*\}$ obtained from $\{A_1,\ldots,A_r\}$ by moving at most $r n^{0.2}$ vertices from large parts to small parts. Let $A_i'$ be a minimum vertex cover of $G[A_i^*]$ for each $i\in [r]$. By Lemma \ref{vertex-cover}, we have $\max\{|A_1'|,\ldots,|A_r'|\}\geq \frac{1}{r}\sqrt{n}$. Without loss of generality, assume that $|A_1'|\geq \frac{1}{r}\sqrt{n}$. Then $G[A_1^*]$ has a matching $M$ of size at least $\frac{\sqrt{n}}{2r}$. By Chernoff we have $\mathbb{P}[|M[S]|>\frac{\sqrt{n}}{9r}]>0.99$. Observe that at most $n^{0.2}$ edges in $M[S]$ contained in a moved vertex, so this event gives a matching of size $\frac{\sqrt{n}}{10r}$ in $S\cap A_1$. By Lemma \ref{lem-int} and Lemma~\ref{lem:binomial distribution},
$$
\mathbb{P}[\frac{|A_1|}{2}+n^{0.4}\leq |A_1\cap S|\leq \frac{|A_1|}{2}+\frac{\sqrt{n}}{10r}]=\mathbb{P}[\frac{|A_1|}{2}+n^{0.4}\leq B(|A_1|,1/2)\leq \frac{|A_1|}{2}+\frac{\sqrt{n}}{10r}],
$$
which by Lemma \ref{lem-int} differs by at most $\delta$ from $$
\int:=\int_{0}^{\frac{1}{10r}}\frac{1}{\sqrt{2\pi}}e^{-t^2/2}\,dt=\frac{1}{\sqrt{2\pi}} \left[ \frac{1}{10r} - \frac{1}{6 \cdot (10r)^3} + O\left(\frac{1}{r^5}\right) \right]\geq \frac{1}{\sqrt{2\pi}}\frac{1}{20r}.
$$
Thus, $\mathbb{P}(E)\geq \frac{1}{150r}$. Moreover, $\mathbb{P}[r\mid |S|]=\frac{1}{r}$, and with high probability, for each $i\in [2,r]$, we have 
$
|A_i\cap S|=\frac{|A_i|}{2}\pm n^{0.1}.
$ Thus, with high probability, we have $|A_i\cap S|\leq \frac{|S|}{r}$ for each $i\in [2,r]$ and $G[A_1\cap S]$  has at least $|A_1\cap S|-\frac{|S|}{r}$ matching edges. Together with Lemma \ref{lem:balance}, we know that $G[S]$ has a $K_r$-factor with probability at least $\frac{1}{200r^2}$, as desired.

\medskip
{\bf Subcase 2.2} $r-s\geq 1$.
\medskip

In this subcase, $B$ is a set of size $(r-s)n$ containing no $\gamma'$-independent set of size $n$. Choose $\eps'\ll\eps\ll \gamma'$. Recall that 
\[
\de(G[B])\ge (r-1)n+1-sn-r\alpha n\ge \left(1-\frac{1}{r-s}-\alpha^{1/5}\right)\cdot|B|.
\]
By Theorem \ref{SuperT}, there are at least $\eps' n$ vertex-disjoint copies of $K_{r-s+1}$ in $G[B]$. By Chernoff, we know that with high probability, $G[B\cap S]$ has at least $\frac{\eps'}{3}n$ vertex-disjoint  copies of $K_{r-s+1}$. By a similar discussion as Subcase 2.1, we obtain that 
$$
\mathbb{P}[\frac{|B1|}{2}+n^{0.4}\leq |B\cap S|\leq \frac{|B|}{2}+\frac{\sqrt{n}}{10r}]\geq \frac{1}{150r},\ \mathbb{P}[r\mid |S|]=\frac{1}{r}, 
$$
and with high probability we have 
$
|A_i\cap S|=\frac{|A_i|}{2}\pm n^{0.1}
$ for each $i\in [s]$. Thus, with high probability, we have $|A_i\cap S|\leq \frac{|S|}{r}$ for each $i\in [s]$ and $G[B\cap S]$ has at least $\frac{\eps'}{3}n$ matching edges. Together with Lemma \ref{lem:balance}, we know that $G[S]$ has a $K_r$-factor with probability at least $\frac{1}{200r^2}$, as desired. 
\end{proof}


%In this section, we deal with the case that every part is small after the above removing procedures, namely, $|A_i|=n$ for each $i\in [s]$.
\subsection{Proof of two lemmas}
In this subsection, we give the proof of the two lemmas to find matchings.
\begin{proof}[Proof of Lemma~\ref{lem:inherit matching}]
Note that $G[A_i]$ contains a matching $M_i$ of size at least $\frac{k_i}{4\beta}$.
Let $X=\sum_{e=\{uv\}\in M_i}X_e$ where $$X_e= 
\begin{cases} 
1, & u\in A_i\cap S~\text{and}~v\in A_i\cap S \\
0, &~\text{otherwise}
\end{cases}$$ is an indicator variable.
Note that $X\sim B(|M_i|,1/2)$.
As $\mathbb{P}[X_e=1]=\mathbb{P}[u\in A_i\cap S]\cdot \mathbb{P}[v\in A_i\cap S]=\frac{1}{4}$, we have
\[
\mathbb{E}[X]=\sum_{e\in M_i}\mathbb{E}[X_e]=\frac{1}{4}|M_i|\ge \frac{k_i}{16\beta}.
\]
By chernoff bound, as $\beta\ll\frac{1}{r}$, we have
\[
\mathbb{P}[X\le \frac{1}{2}\mathbb{E}[X]]\le2\exp\left(-\frac{k_i}{160\beta}\right)=o(1),
\]
which implies that there exists a matching of size at least $\frac{k_i}{32\beta}$ in $A_i\cap S$ with probability $1-o(1)$.
\end{proof}



\begin{proof}[Proof of Lemma~\ref{vertex-cover}]
Assume that $|A_1'|\geq |A_2'|\geq \cdots \geq |A_r'|$. Clearly, there exists a vertex cover $A_i^1$ with size $|A_2'|$ for each $i\in [3,r]$. For convenience, let  $A_1^1:=A_1'$ and $A_2^1:=A_2'$. Denote $x_i:=|A_i^1|$ and $A_i^2:=A_i\setminus A_i^1$ for each $i\in [r]$. Hence, $G[A_i^2]=\emptyset$ for each $i\in [r]$. We proceed by considering the following two cases. 

\medskip
{\bf Case 1. $e(A_1^2,\bigcup_{j\neq 2}A_j^1)\leq e(A_2^1,\bigcup_{j\neq 1}A_j^2)$.}
\medskip

Recall that $G$ is an $((r-1)n+1)$-regular graph. Hence 
$$
e(A_2^1,A_1^2)\leq |A_2^1|((r-1)n+1)-e(A_2^1,\bigcup_{j\neq 1}A_j^2).
$$
Note that
$$
e(A_1^2,\bigcup_{j\neq 1}A_j^2)=\sum_{j\neq1} e(A_1^2,A_j^2)\le\sum_{j\neq1}(n-x_1)(n-x_j)=(n-x_1)\left((r-1)n-\sum_{j\neq1}x_j\right).
$$
Therefore, we have
\begin{align*}
e(A_1^2,\bigcup_{j\neq 2}A_j^1)&=|A_1^2|((r-1)n+1)-e(A_1^2,A_2^1)-e(A_1^2,\bigcup_{j\neq 1}A_j^2)\\
&\geq (n-x_1)((r-1)n+1)-x_2((r-1)n+1)+e(A_2^1,\bigcup_{j\neq 1}A_j^2)-e(A_1^2,\bigcup_{j\neq 1}A_j^2)\\
&\geq (n-x_1-x_2)((r-1)n+1)+e(A_2^1,\bigcup_{j\neq 1}A_j^2)-(n-x_1)\Big((r-1)n-\sum_{j\neq 1}x_j\Big)\\
&=(n-x_1-x_2)((r-1)n+1)+e(A_2^1,\bigcup_{j\neq 1}A_j^2)-(n-x_1)(r-1)(n-x_2),
\end{align*}
where the last equality holds as $x_j=|A_2'|$ for each $j\in [2,r]$.
By the assumption, one has
\begin{align*}
(n-x_1-x_2)((r-1)n+1)\leq (n-x_1)(r-1)(n-x_2).
\end{align*}
It follows that $n\leq (r-1)x_1x_2+(x_1+x_2)\leq (r-1)x_1^2+2x_1.$ Thus, $x_1\geq \frac{1}{r}\sqrt{n}$, as desired.

\medskip
{\bf Case 2. $e(A_1^2,\bigcup_{j\neq 2}A_j^1)> e(A_2^1,\bigcup_{j\neq 1}A_j^2)$.}
\medskip

Notice that
$$
e(\bigcup_{j\neq 2}A_j^1,\bigcup_{j\neq 1}A_j^2)\leq |\bigcup_{j\neq 2}A_j^1|((r-1)n+1)-e(\bigcup_{j\neq 2}A_j^1,A_1^2).
$$
Therefore,
\begin{align*}
e(\bigcup_{j\neq 1}A_j^2,A_2^1)=&|\bigcup_{j\neq 1}A_j^2|((r-1)n+1)-e(\bigcup_{j\neq 1}A_j^2,\bigcup_{j\neq 2}A_j^1)-e(\bigcup_{j\neq 1}A_j^2,\cup_{t\in [r]}A_t^2)\\
\geq &|\bigcup_{j\neq 1}A_j^2|((r-1)n+1)-|\bigcup_{j\neq 2}A_j^1|((r-1)n+1)+e(\bigcup_{j\neq 2}A_j^1,A_1^2)-\sum_{j\neq 1}\sum_{t\neq j}e(A_j^2,A_t^2)\\
\geq &((r-1)n-\sum_{j\neq 1}x_j)((r-1)n+1)-\sum_{j\neq 2}x_j((r-1)n+1)+e(\bigcup_{j\neq 2}A_j^1,A_1^2)\\
&-\sum_{j\neq 1}\sum_{t\neq j}(n-x_j)(n-x_t)\\
=&((r-1)n+1)\Big((r-1)n-\sum_{j\neq 1}x_j-\sum_{j\neq 2}x_j\Big)+e(\bigcup_{j\neq 2}A_j^1,A_1^2)-\sum_{j\neq 1}\sum_{t\neq j}(n-x_j)(n-x_t)\\
=&((r-1)n+1)\Big((r-1)n-(r-1)x_2-x_1-(r-2)x_2\Big)+e(\bigcup_{j\neq 2}A_j^1,A_1^2)\\
&-(r-1)(n-x_2)\big((r-1)n-x_1-(r-2)x_2\big)\\
=&(r-1)n-x_1-(2r-3)x_2-(r-1)x_1x_2-(r-1)(r-2)x_2^2+e(\bigcup_{j\neq 2}A_j^1,A_1^2).
\end{align*}
By the assumption, one has
$$
(r-1)n\leq x_1+(2r-3)x_2+(r-1)x_1x_2+(r-1)(r-2)x_2^2\leq (r-1)\big((r-1)x_1^2+2x_1\big).
$$
Thus, $x_1\geq \frac{1}{r}\sqrt{n}$, as desired. 
\end{proof}



\bibliographystyle{abbrv}
\bibliography{ref.bib}
\end{document}

    Our first step is to balance the number of vertices in $A_2,\ldots,A_p$. 

    Suppose that by using $x_i$ copies of $K_r$ with index vector ${\bf 1_{r}}-{\bf u_i}+{\bf u_j}$ for some $j\in [p+1,r]$, the left number of vertices in $A_2,\ldots,A_p$ are equal. We will prove the existence of $x_i$ for each $i\in [p]$. It suffices to consider the solution of the following equations:
    \begin{align*}
        (x_1,\ldots,x_p)(J_p-I_p)=(b_1,b_2,\ldots,b_p),
    \end{align*}
    where $a_1+a_2-b_2=\cdots=a_1+a_p-b_p:=b$. By a simple computation, one has
    $$
    x_1=\frac{1}{p-1}(b_1+\sum_{j=2}^p(a_1+a_j-b))-b,\ \text{and}\ x_i=\frac{1}{p-1}(b_1+\sum_{j=2}^p(a_1+a_j-b))-(a_1+a_i-b)\ \text{for each}\ i\in [2,p].
    $$
    In this process, for each $i\in [p+1,r]$, assume that $k_i'(\leq k_i)$ matching edges in $G[A_i]$ are used. Then the index vector of the above $K_r$-tiling is
    $$
(b_1,a_1+a_2-b,\ldots,a_1+a_p-b,k_{p+1}'+b_1,\ldots,k_{r}'+b_1).
    $$

    Now, the slack of the resulting partition is 
\begin{align*}
&(a_1-b_1,b,\ldots,b,a_1+a_{p+1}-(k_{p+1}'+b_1),\ldots, a_1+a_{r}-(k_{r}'+b_1))-(a_1-b_1){\bf 1_r}\\
=&(0,b-(a_1-b_1),\ldots,b-(a_p-b_p),a_{p+1}-k_{p+1}',\ldots,a_{r}-k_{r}').
\end{align*}
Next, we use $x_i$ copies of $K_r$ with index vector ${\bf 1_r-u_1+u_{i}}$ for each $j\in [p+1,r]$. It suffices to consider the solution of 
\begin{equation*}
\left\{
\begin{array}{ll}
	&x_{p+1}+\cdots+x_r=b-(a_1-b), \\
     &2x_{p+1}+x_{p+2}\cdots+x_r=a_{p+1}-k_{p+1}',\\
     &\cdots\\
     &x_{p+1}+\cdots+x_{r-1}+2x_r=a_{r}-k_{r}'.
\end{array}\right.
\end{equation*}

    Next, use $a_1-b_1$ copies of $K_r$ with index vector ${\bf 1_r}$ to cover all left  vertices in $A_1$. 
    %之前验证B06的证明%{Next we shall verify $B06$?}
{Let $B_1'\cup B_2'$ be a partition of $B$ with $|B_1'|=n$. 
For a random vertex set $B_2\subseteq B_2'$, we have}\todo{where is the randomness in the following event?}
\begin{align*}
&\mathbb{P}\big[E(G[B_1,B_2])\neq \emptyset~\text{for every}~B_1\subseteq B_1'~\text{with}~|B_1|=\frac{n}{2}\pm n^{0.7}\big]\\
&=\sum_{B_1\subseteq B_1'}\mathbb{P}[E(G[B_1,B_2])\neq \emptyset\mid B_1\subseteq B_1'~\text{with}~|B_1|=\frac{n}{2}\pm n^{0.7}]\cdot \mathbb{P}[B_1\subseteq B_1'~\text{with}~|B_1|=\frac{n}{2}\pm n^{0.7}].
\end{align*}
% if $B_1\subseteq B_1'$ is chosen uniformly at random with $|B_1|=\frac{n}{2}\pm n^{0.7}$, then  % we have  % with size  For every two disjoint sets $B_1,B_2\subseteq B$ with $|B_1|=|B_2|\leq \frac{2}{3}n$, there exist disjoint sets $B_1',B_2'$ in $B$ such that  $B_1'\supseteq B_1$, $|B_1'|=n$ and $B_2'\supseteq B_2$. Therefore,
Since $G$ is $((r-1)n+1)$-regular and $|B_1'|=n$, for every vertex $v\in B_2'$ we obtain that $$|N_G(v)\cap B_1'|\ge (r-1)n+1-((r-1)n-1)\ge 2.$$ \todo{more details} 
Thus we have
\begin{align*}
&\mathbb{P}[E(G[B_1,B_2])\neq \emptyset\mid\text{{conditional events?}} B_1\subseteq B_1'~\text{with}~|B_1|=\frac{n}{2}\pm n^{0.7}]\\
=&1-\mathbb{P}[E(G[B_1,B_2])=\emptyset\mid B_1\subseteq B_1'~\text{with}~|B_1|=\frac{n}{2}\pm n^{0.7}]\\=& 1-\frac{1}{\binom{|B_1'|}{|B_1|}}\cdot\binom{|B_1'\setminus N_{B_1'}(B_2)|}{|B_1|}\\
\geq &1-\frac{1}{\binom{|B_1'|}{|B_1|}}\cdot\binom{|B_1'|-2}{|B_1|}\\
=&1-\frac{(|B_1'|-|B_1|)(|B_1'|-|B_1|-1)}{|B_1'|(|B_1'|-1)}\\
\geq& \frac{3}{4}-o(1),
\end{align*}
where $\binom{|B_1'\setminus N_{B_1'}(B_2)|}{|B_1|}$ is the number of choice of $B_1$ such that there is no edge between $B_1$ and $B_2$.
%\begin{align*}
%&\mathbb{P}[\text{There exists at least one edge between}~B_1~\text{and}~B_2]\\
%&=1-\mathbb{P}[\text{There is no edge between}~B_1~\text{and}~B_2]\\
%&= 1-\frac{1}{\binom{|B|}{|S|/r}}\cdot\sum_{B_1\subseteq B}\frac{2^{n-|S|/r-|N(B_1)|}}{2^{n-|S|/r}}
%\end{align*}
Therefore, {as $B_2$ is chosen uniformly and randomly?}, with probability at least $\frac{3}{4}-o(1)$ the following holds: {double check the proof of D7}
\begin{enumerate}
    [label =\rm  (D7)]\item\label{D6} for any two sets of size $\frac{|S|}{r}$ in $B\cap S$,  there is at least one  edge between them.
\end{enumerate}






%Recall that $A_i\cap S$ is called a large part if $|A_i\cap S|>\frac{|S|}{r}$, and a small part otherwise.
By Lemma~\ref{lem:covering bad vertices for balance}, in order to obtain a $K_r$-factor in $G[S]$, it suffices to construct a $(3\alpha,\alpha^{1/5},\frac{\beta}{4},\frac{\gamma'}{2})$-good partition of $G[S]$ under the condition that $r\mid |S|$. Conditional on \ref{D2}–\ref{D6}, {it suffices to verify $B2$?} i.e., to find a matching of size at least {$|A_i \cap S| - \frac{|S|}{r}$??} in each large part of $\mathcal{P}'$. 

FKG对B06：{(B06) Next we use FKG Inequality to prove that $\mathcal{P}'$ w.h.p.~satisfies~\ref{B07}.}
{todo: Check the proof! (Yang: If B itself is not $o(1)$-close to the two disjoint copies of $K_n$, then maybe we can apply Lemma~3.3 to verify that w.h.p, $B_8$ happens. Otherwise $B$ has a sparse cut $\{B_1,B_2\}$, and we still know from the min degree that $e(B_1,B_2)\ge n/2$, while each $B_i$ is nearly complete. In this case, $B06$ fails (i.e, there is no edge in $X,Y\< B\cap S$) if and only if $X\<B1, Y\< B2$???, and essentially we have $|B1\cap S|\ge |X|, |B2\cap S|\ge |Y|$, and )}

{If $s=r-2$, then~\ref{B1} implies that the size of $B$ is $2n\pm2\alpha n$.
Recall that $|B\cap S|=\frac{|B|}{2}\pm n^{0.6}=n\pm 2\alpha n$ and $S$ is obtained by independently selecting vertices from $V(G)$ with probability $1/2$, without replacement. 
We may assume that the size of $B\cap S$ is exactly $n$ and let $B\cap S=\{x_1,x_2,\cdots,x_n\}$ where these $n$ vertices are obtained through sequential selection in the given order. In fact, the error term associated with $|B\cap S|$ has no impact on the validity of the subsequent probability calculations.
Without loss of generality, we define $X$ as the set of the first $\frac{n}{2}-\beta n$ elements in $B\cap S$, and let $Y$ be the set containing all the others. 
For each edge $e\in E(G[B])$, define the indicator
\[
I_e=
\begin{cases}
1, & \text{if } e \in E(G'[X,Y]), \\
0, & \text{if } e \notin E(G'[X,Y]).
\end{cases}
\]
For any set $X\subseteq B\cap S$ with $|X|=\frac{n}{2}-\beta n$, we choose a set $Y$ with $X\cap Y=\emptyset$ and $X\cup Y=B\cap S$. Note that $S$ and $Y$ share the same randomness.
By the FKG inequality, we have
\begin{align*}
&\mathbb{P}\left[\text{There exists}~S~\text{such that}~e_S(X,Y)=0~\text{for any}~X\subseteq B\cap S~\text{with}~|X|=\frac{n}{2}-\beta n\right]\\
&\le\binom{n}{\frac{n}{2}-\beta n}\cdot\mathbb{P}\left[\sum_{e\in E(G[X,B\setminus X])}I_e=0\right]{\ge ?(\frac{1}{2})^{e(X, B\setminus X)}}\\
&\le 2^n\cdot e^{-\Theta(\beta^2 n)}\cdot \left(\frac{1}{2}\right)^{n^2/4}=o(1),
\end{align*}
where the last inequality holds since every edge is chosen into $S$ with probability $1/2$ and there are at least $\frac{n^2}{4}$ edges between $X$ and $B\setminus X$.
Indeed, as $G$ is $((r-1)n+1)$-regular, for every vertex $v\in X$, we have
\[
d_B(v,B\setminus X)\ge (r-1)n+1-(r-2)n-|X|\ge \frac{n}{2}.
\]
Thus,~\ref{B07} holds with probability $1-o(1)$.}
\end{proof}

% Note that $\mathbb{P}[\sum_{e\in E(G[X,B\setminus X])}I_e=0]=\mathbb{E}[\prod_{e\in E(G[X,B\setminus X])}(1-I_{e_j})]$. Let $\mathcal{F}_k=\sigma(I_{e_1},\cdots,I_{e_k})$.
% By Fubini Theorem, 
% \[
% \mathbb{E}[\prod_{j=1}^{m}(1-I_{e_j})]=\mathbb{E}[\mathbb{E}[\prod_{j=1}^{m}(1-I_{e_j})\mid \mathcal{F}_{m-1}]]=\mathbb{E}[\prod_{j=1}^{m}(1-I_{e_j})\cdot\mathbb{E}[1-I_{e_m}\mid \mathcal{F}_{m-1}]].
% \]




% {Under the conditioning that the above events occur, we verify sequentially the second part of \ref{B2}, \ref{B07}, and finally \ref{B1}.}

% {We first consider \ref{B07} when $s=r-2$.  Pick uniformly at random two subsets  $X, Y\subseteq B\cap S$.} %By the Chernoff bound, w.h.p. we have
% %\[
% %|X|=\tfrac{n}{2}\pm n^{0.7} \text{ and } |Y|=\tfrac{n}{2}\pm n^{0.7}.
% %\]
% Since $G$ is $((r-1)n+1)$-regular, each vertex has
% at least $n-o(n)$ neighbours in $B$.
% Hence
% \[
% e(G[B]) \ge \tfrac12\cdot |B|\cdot (n-o(n)) = n^2-o(n^2).
% \]

% For each edge $e=uv\in E(G[B])$, define the indicator
% \begin{equation*}
% I_e= \left\{
% \begin{array}{ll}
% 	1 & \{u\in X,\, v\in Y\}\cup\{u\in Y,\, v\in X\}, \\
%      0 & {\text{otherwise}},
% \end{array}\right. \  \qquad Z = \sum_{e\in E(G[B])} I_e = |E(G[X,Y])|.
% \end{equation*}
% For each fixed edge $e=uv$ we have
% \[
% \mathbb{P}(I_e=1) = \mathbb{P}(u\in X, v\in Y) + \mathbb{P}(u\in Y, v\in X) = \tfrac14 + \tfrac14 = \tfrac12.
% \]
% Hence
% \[
% \mathbb E[Z] = \sum_{e\in E(G[B])} \mathbb E[I_e]
% = \tfrac12 e(G[B])
% \ge \tfrac12 (n^2-o(n^2)).
% \]

% Although the indicators $\{I_e\}$ are not mutually independent, two such variables
% are independent whenever the corresponding edges are vertex-disjoint.
% Hence only pairs of edges sharing a vertex contribute to the covariance terms. 
% There are at most $O(n^3)$ such pairs, while $e(G[B])=\Theta(n^2)$.
% Therefore
% \[
% \operatorname{Var}(Z)
% = \sum_e \operatorname{Var}(I_e)
% + \sum_{e\neq f}\operatorname{Cov}(I_e,I_f)
% = O(n^3),
% \]
% whereas $\mathbb E[Z]^2=\Theta(n^4)$.
% By Chebyshev's inequality,
% \[
% \Pr(Z=0) \le \frac{\operatorname{Var}(Z)}{\mathbb E[Z]^2} = O(\frac{1}{n})=o(1).
% \] 
% Therefore, $\mathbb{P}(E(G[X,Y])\neq\emptyset)=1-o(1)$, that is, w.h.p. \ref{B07} holds.

% {Next we shall verify $B06$?}
% {Let $B_1'\cup B_2'$ be a partition of $B$ with $|B_1'|=n$. 
% For a random vertex set $B_2\subseteq B_2'$, we have}\todo{where is the randomness in the following event?}
% \begin{align*}
% &\mathbb{P}\big[E(G[B_1,B_2])\neq \emptyset~\text{for every}~B_1\subseteq B_1'~\text{with}~|B_1|=\frac{n}{2}\pm n^{0.7}\big]\\
% &=\sum_{B_1\subseteq B_1'}\mathbb{P}[E(G[B_1,B_2])\neq \emptyset\mid B_1\subseteq B_1'~\text{with}~|B_1|=\frac{n}{2}\pm n^{0.7}]\cdot \mathbb{P}[B_1\subseteq B_1'~\text{with}~|B_1|=\frac{n}{2}\pm n^{0.7}].
% \end{align*}
% % if $B_1\subseteq B_1'$ is chosen uniformly at random with $|B_1|=\frac{n}{2}\pm n^{0.7}$, then  % we have  % with size  For every two disjoint sets $B_1,B_2\subseteq B$ with $|B_1|=|B_2|\leq \frac{2}{3}n$, there exist disjoint sets $B_1',B_2'$ in $B$ such that  $B_1'\supseteq B_1$, $|B_1'|=n$ and $B_2'\supseteq B_2$. Therefore,
% Since $G$ is $((r-1)n+1)$-regular and $|B_1'|=n$, for every vertex $v\in B_2'$ we obtain that $$|N_G(v)\cap B_1'|\ge (r-1)n+1-((r-1)n-1)\ge 2.$$ \todo{more details} 
% Thus we have
% \begin{align*}
% &\mathbb{P}[E(G[B_1,B_2])\neq \emptyset\mid\text{{conditional events?}} B_1\subseteq B_1'~\text{with}~|B_1|=\frac{n}{2}\pm n^{0.7}]\\
% =&1-\mathbb{P}[E(G[B_1,B_2])=\emptyset\mid B_1\subseteq B_1'~\text{with}~|B_1|=\frac{n}{2}\pm n^{0.7}]\\=& 1-\frac{1}{\binom{|B_1'|}{|B_1|}}\cdot\binom{|B_1'\setminus N_{B_1'}(B_2)|}{|B_1|}\\
% \geq &1-\frac{1}{\binom{|B_1'|}{|B_1|}}\cdot\binom{|B_1'|-2}{|B_1|}\\
% =&1-\frac{(|B_1'|-|B_1|)(|B_1'|-|B_1|-1)}{|B_1'|(|B_1'|-1)}\\
% \geq& \frac{3}{4}-o(1),
% \end{align*}
% where $\binom{|B_1'\setminus N_{B_1'}(B_2)|}{|B_1|}$ is the number of choice of $B_1$ such that there is no edge between $B_1$ and $B_2$.
% %\begin{align*}
% %&\mathbb{P}[\text{There exists at least one edge between}~B_1~\text{and}~B_2]\\
% %&=1-\mathbb{P}[\text{There is no edge between}~B_1~\text{and}~B_2]\\
% %&= 1-\frac{1}{\binom{|B|}{|S|/r}}\cdot\sum_{B_1\subseteq B}\frac{2^{n-|S|/r-|N(B_1)|}}{2^{n-|S|/r}}
% %\end{align*}
% Therefore, {as $B_2$ is chosen uniformly and randomly?}, with probability at least $\frac{3}{4}-o(1)$ the following holds: {double check the proof of D7}
% \begin{enumerate}
%     [label =\rm  (D7)]\item\label{D6} for any two sets of size $\frac{|S|}{r}$ in $B\cap S$,  there is at least one  edge between them.
% \end{enumerate}

\section{Probabilistic characterization}
For Conjecture \ref{conj:K_r factor}, Dragani\'{c}, Keevash and M\"{u}yesser \cite{Keevash2025} identified a natural extremal construction: slightly unbalanced complete $r$-partite graphs with a prescribed factor structure imposed on the largest part $A$.  This suggests the optimal constant $c$ should be  $\frac{1}{r^2}$, where in the random induced subgraph $G[S]$ we have one
probability factor of $\frac{1}{r}$ for $r$ dividing $|S|$ and another for $A\cap S$ being the largest part.

\begin{lemma}
    There exists an $((r-1)n+1)$-regular  graph $G$ on $rn$ vertices  such that at most $\frac{2}{3r^2}2^{rn}$ subsets of $V(G)$ induce a $K_r$-factor.  
\end{lemma}
% as the probability in $G[S]$ contains one factor of $\frac{1}{r}$ for $r$ dividing $|S|$ and another for $A\cap S$ being the largest part. 

%In~\cite{Keevash2025}, they conjecture that the optimal $c$ in Conjecture~\ref{conj:K_r factor} should be $\frac{1}{r^2}$. The speculation arises from an extremal construction: slightly unbalanced complete $r$-partite graph $G$ with a suitable factor in the larger part $A_1$. In the random induced subgraph $G[S]$, they suggest one probability factor of $1/r$ for $r\mid |S|$ and another for $|A_1\cap S|$ being the largest part.
\begin{proof}
We first consider an $rn$-vertex complete $r$-partite graph $G_0$ with vertex partition $A_1\cup \ldots \cup A_r$ with $|A_1|=n+r-1$ and $|A_2|=\cdots=|A_r|=n-1$. Let $G$ be a graph obtained from $G_0$ by embedding an $r$-regular Hamiltonian graph\footnote{may assume that $rn$ is even.} into the part $A_1$ of $G_0$ {such that $G[A_1]$ has no $K_3$}. Clearly, $G$ is an $((r-1)n+1)$-regular graph and $G[A_i]$ is empty for each $i\in [2,r]$. Let  $S\subseteq V(G)$. %with $r\mid |S|$ formed by including each vertex of $V(G)$ independently with probability $\frac{1}{2}$. %, and let $\mathcal{P}=\{A_1\cap S,\ldots,A_r\cap S\}$. 


{We claim that a sufficient and necessary condition for $G[S]$ to contain a $K_r$-factor is that 
\begin{enumerate}
    \item $r\mid |S|$;
    \item $\frac{|S|}{r}\le |A_1 \cap S|\le \frac{2|S|}{r}$;
   \item $|A_i \cap S|\le \frac{|S|}{r}$ for all $i \in [2,r]$.
\end{enumerate}
Next we count the total number of such sets $S$. Suppose $|S|=rm$ and $|A_i \cap S|=:m_i,i\in[r]$. Then 
there are at most\[
\sum_{m_2+\cdots+m_r\ge (r-2)m}\prod_{i\in[r]}\binom{n}{m_i}
\le m^{r-1}\binom{n}{m}^r\] choices for such $S$, where we use the log-concavity of the binomial coefficient. Then the total number of choices for $S$ is at most $\sum_{m\in[n]}m^{r-1}\binom{n}{m}^r$ (this is bad)}
Indeed, if $G[S]$ admits a $K_r$-factor, then since $G[A_i]$ is empty for each $i \in [2,r]$, each copy of $K_r$ can contain at most one vertex from each $A_i$. It follows that $|A_i \cap S|$ is at most the number of $K_r$-copies, which is $\frac{|S|}{r}$. %Otherwise, we call it \textit{space barrier} that prevent the existence of $K_r$-factors in $G[S]$.
This shows that requiring $A_1 \cap S$ to be the largest part alone is not sufficient. 
%Indeed, the condition we need is that the average $\frac{\sum_{i=1}^r|A_i\cap S|}{r}$ is at least $|A_j\cap S|$ for all $j\in \{2,3,\cdots,r\}$.
Furthermore, 
\[
\mathbb{P}\left[\frac{\sum_{i=1}^r|A_i\cap S|}{r}\ge |A_j\cap S|~\text{for all}~j\in [2,r]\right]< \mathbb{P}\left[|A_1\cap S|\ge |A_j\cap S|~\text{for all}~j\in [2,r]\right].
\]

Define $\mathbf{E}$ as the event that $\frac{\sum_{i=1}^r |A_i \cap S|}{r} \geq \max_{j \in [2,r]} |A_j \cap S|$. We proceed to  estimate the probability of $\mathbf{E}$ using the normal distribution. % that the average $\frac{\sum_{i=1}^r|A_i\cap S|}{r}$ is at least $|A_j\cap S|$ for all $j\in \{2,3,\cdots,r\}$
%Let $G$ be an $((r-1)n+1)$-regular graph on $rn$ vertices obtained from the complete $r$-partite graph $K_{n+r-1,n-1,\cdots,n-1}$ by adding an $r$-regular graph within a part of size $n+r-1$. Let $|A_1|=n+r-1, |A_i|=n-1$ for each $i\in \{2,3,\cdots,r\}$ and $a_i:=|A_i\cap S|$. 
Let $a_j:=|A_j\cap S|$ and  $U_j:=a_1+\sum_{i\in[r]\setminus\{1,j\}}a_i-(r-1)a_j$ for each $j\in [2,r]$. 
Note that 
\[
\mathbb{P}\left[
\mathbf{E}\right]=\mathbb{P}\left[U_j\ge 0~\text{for all}~j\in [2,r]\right].
\]
Let $a_1=m_1+\sigma_1X$ and $a_j=m+\sigma Y_j$ for each $j\in [2,r]$,  where $m_1=\frac{n+r-1}{2}, \sigma_1=\sqrt{\frac{n+r-1}{4}}, m=\frac{n-1}{2}, \sigma=\sqrt{\frac{n-1}{4}}$.
It is obvious that $X, Y_j\sim N(0,1)$ for each $j\in [2,r]$.
Thus 
\begin{align*}
U_j&=a_1+\sum_{i\in[r]\setminus\{1,j\}}a_i-(r-1)a_j\\
&=m_1+\sigma_1X+(r-2)m+\sigma\sum_{i\in[r]\setminus\{1,j\}}Y_i-(r-1)(m+\sigma Y_j)\\
&=\frac{r}{2}+\sigma_1X+\sigma\sum_{i\in [r]\setminus\{1,j\}}Y_i-(r-1)\sigma Y_j.
\end{align*}
Define  $W_j=\sigma_1X+\sigma\sum_{i\in [r]\setminus\{1,j\}}Y_i-(r-1)\sigma Y_j$ for each $j\in [2,r]$.
Note that $U_j\ge 0$ implies that $W_j\ge -\frac{r}{2}$. 
%To simplify the computations, we derive the results exclusively for $r=3$.
Moreover, 
\[
Var(W_j)=\sigma_1^2+\sigma^2(r^2-r-1)\approx \frac{nr(r-1)}{4}~\text{and}~\sigma_{W_j}\approx\sqrt{\frac{nr(r-1)}{4}}.
\]
We now turn to calculating the covariance for distinct $i,j\in [2,r]$,
\begin{align*}
Cov(W_i,W_j)=\sigma_1^2-(r+1)\sigma^2\approx -\frac{rn}{4}.
\end{align*}
Hence, the correlation coefficient between $W_i$ and $W_j$ is $$\rho_{i,j}:=\frac{Cov(W_i,W_j)}{\sigma_{W_i}\sigma_{W_j}}\approx\frac{-\frac{rn}{4}}{\frac{nr(r-1)}{4}}=-\frac{1}{r-1}.$$
Let $\xi_j=\frac{W_j}{\sigma_{W_j}}$ for each $j\in [2,r]$.
Note that $\xi_2,\cdots,\xi_r\sim N(0,1)$, and for any distinct indices $i, j$, the correlation coefficient between $\xi_i$ and $\xi_j$ is $-\frac{1}{r-1}$.
{Let $\sum$ be the covariance matrix of $\xi_2,\cdots,\xi_r$ with dimension $r-1$, whose diagonal elements are all $1$ and off-diagonal elements are all $-\frac{1}{r-1}$.
Note that $|\sum|=\frac{1}{r-1}\cdot(\frac{r}{r-1})^{r-2}.$
Define $\mathbf{X}$ as a random vector of dimension $r-1$ where $\mathbf{X}\sim N(0,\sum)$.
Since $n$ is sufficiently large, we have
\begin{align*}
\mathbb{P}\left[\mathbf{E}\right]&=\mathbb{P}\left[U_j\ge 0~\text{for all}~j\in [2,r]\right]\\&\approx \mathbb{P}\left[W_j\ge -\frac{r}{2}~\text{for all}~j\in [2,r]\right]\\&=\mathbb{P}\left[\xi_2\ge 0,\cdots,\xi_r\ge 0\right]=\mathbb{P}[\mathbf{X}\ge 0]\\
&=\int_{-\infty}^{\infty}\cdots\int_{-\infty}^{\infty}\phi(\mathbf{X})\mathbb{I}(\mathbf{X}\ge 0)d\mathbf{X}\approx ,
% &{=?}\int_{-\infty}^{\infty}\phi(z)\left[\Phi\big(\frac{z}{\sqrt{r}}\big)\right]^{r-1}dz \\
% &{\le~\text{how?}~}\frac{1}{2r}+\frac{1}{4\sqrt{\pi}r^{3/2}}+O\left(\frac{1}{r^{5/2}}\right),
\end{align*}
where $\phi(\mathbf{X})=\frac{1}{\sqrt{(2\pi)^{r-1}|\sum|}}\exp(-\frac{1}{2}\mathbf{X}^T\sum^{-1}\mathbf{X})$ and $\mathbb{I}(\mathbf{X}\ge 0)$ denotes the indicator function which equals $1$ if all components of $\mathbf{X}$ are nonnegative, and $0$ otherwise.}
\end{proof}

计算A_i'
\begin{itemize}
    \item if $|A_s|>n$, then  $A_s\cap V(\mathcal{H})=\emptyset$ by~\ref{E1} and  for each $i\in [s-1]$ we have 
\begin{align*}%\label{eq:A1}
|A_i'|-\frac{|G^*|}{s+1}\leq %|A_i|-\frac{|V(\mathcal{K}\cup \mathcal{H})|}{r}-\frac{|G'|}{r}=
|A_i|-\frac{|V(\mathcal{K}\cup \mathcal{H})|}{r}-\frac{|G|-|V(\mathcal{K}\cup \mathcal{H})|}{r}=|A_i|-n\leq 2\alpha \cdot \frac{|G^*|}{s+1},
\end{align*}
and  
$$
|A_s'|-\frac{|G^*|}{s+1}\leq |A_s|-\frac{|V(\mathcal{K})|}{r}-\frac{|G^*|}{s+1}=|A_s|-n+q\leq |A_s|-n+r-1\leq 2\alpha \cdot \frac{|G^*|}{s+1}; 
$$
\item if $|A_s|\leq n$, then by~\ref{E2}, for each $i\in [p]$,  we have   $|V(\mathcal{H})\cap A_i|=q-q_i$ and   
\begin{align*}
|A_i'|-\frac{|G^*|}{s+1}&\leq |A_i|-\frac{|V(\mathcal{K})|}{r}-|V(\mathcal{H})\cap A_i|-\frac{|G^*|}{s+1}= |A_i|-\frac{|V(\mathcal{K})|}{r}-(q-q_i)-\frac{|G^*|}{s+1}\\
&=|A_i|-n+q_i\leq 0,
\end{align*}
for each $i\in [p+1,s]$, we have  $|V(\mathcal{H})\cap A_i|=q$ and 
$$
|A_i'|-\frac{|G^*|}{s+1}\leq |A_i|-\frac{|V(\mathcal{K}\cup \mathcal{H})|}{r}-\frac{|G^*|}{s+1}=|A_i|-\frac{|V(\mathcal{K})|}{r}-q-\frac{|G^*|}{s+1}={|A_i|-n}=0.
$$
\end{itemize}

